\documentclass[11pt,a4paper]{article}

\usepackage{iftex}
\ifPDFTeX
  \usepackage[utf8]{inputenc}
  \usepackage[T1]{fontenc}
  \usepackage{lmodern}
\else
  \usepackage{fontspec}
  \setmainfont{Latin Modern Roman}
  \setsansfont{Latin Modern Sans}
  \setmonofont{Latin Modern Mono}
\fi
\usepackage[spanish,es-noindentfirst]{babel}
\usepackage{amsmath,amssymb,amsthm}
\usepackage{stmaryrd}
\usepackage[margin=3cm]{geometry}
\usepackage{enumitem}
\usepackage{tikz-cd}
\usetikzlibrary{arrows.meta,calc}
\usepackage[colorlinks=true,linkcolor=black,citecolor=black,urlcolor=black]{hyperref}

% ---------- Entorno austero: serif, sin color ----------
\theoremstyle{plain}
\newtheorem{theorem}{Teorema}[section]
\newtheorem{proposition}[theorem]{Proposición}
\newtheorem{lemma}[theorem]{Lema}
\newtheorem{corollary}[theorem]{Corolario}
\theoremstyle{definition}
\newtheorem{definition}[theorem]{Definición}
\newtheorem{example}[theorem]{Ejemplo}
\theoremstyle{remark}
\newtheorem{remark}[theorem]{Observación}
\newtheorem{note}[theorem]{Nota}
\newtheorem{question}[theorem]{Pregunta}

\newcommand{\RR}{\mathbb R}
\newcommand{\NN}{\mathbb N}
\newcommand{\ZZ}{\mathbb Z}
\newcommand{\QQ}{\mathbb Q}
\newcommand{\CC}{\mathbb C}
\DeclareMathOperator{\ran}{ran}
\DeclareMathOperator{\Stab}{Stab}
\DeclareMathOperator{\Aut}{Aut}
\DeclareMathOperator{\Lev}{Lev}
\DeclareMathOperator{\Mor}{Mor}
\DeclareMathOperator{\Homeo}{Homeo}
\DeclareMathOperator{\im}{im}
\DeclareMathOperator{\Hom}{Hom}
\newcommand{\CWq}{\mathrm{CW}}

\title{\normalfont\bfseries Topología algebraica\\ \large Observaciones, notas y ejercicios}
\author{César Galindo López\\ \normalsize Universidad Panamericana}
\date{\normalsize }

\begin{document}
\shorthandoff{"<>}
\maketitle

\begin{abstract}
\noindent Estas notas cubren tres bloques estándar de un primer curso de
topología algebraica --- homotopía y grupo fundamental, cubrientes, y
homología singular y celular --- con las demostraciones y una colección de
ejercicios resueltos que fijan el uso de cada resultado. No se sigue un
único libro de texto; el orden y las convenciones (coeficientes en un
anillo $R$ arbitrario para homología, cubrientes definidos vía trivialización
local) son los de un curso de posgrado típico.
\end{abstract}

\tableofcontents

\section{Homotopía}

\begin{definition}
\label{def:caminos}
Sean $x_0$ y $x_1$ puntos de un espacio $X$. Denotaremos por
\[
\Omega(X,x_0,x_1) = \{\alpha : I \to X \mid \alpha \text{ es continua y } \alpha(0)=x_0 \text{ y } \alpha(1)=x_1\}
\]
al conjunto de todas las trayectorias (o caminos) de origen $x_0$ y extremo $x_1$.

Cuando $x_0=x_1$, decimos que el espacio $(X,x_0)$ está basado en $x_0$ y
$\Omega(X,x_0)=\Omega(X,x_0,x_0)$ denota al conjunto de lazos en $x_0$.
\end{definition}

\begin{definition}
\label{def:homotopia-lazos}
Sean $\alpha,\beta\in\Omega(X,x_0)$. Decimos que $\alpha$ y $\beta$ son
\emph{homotópicos}, denotado $\alpha\simeq\beta$, si existe una función
$H : I\times I \to X$ continua tal que
\[
H(s,0)=\alpha(s), \qquad H(s,1)=\beta(s), \qquad H(0,t)=H(1,t)=x_0 \ \forall t\in I.
\]
A tal función se le llama \emph{homotopía relativa a $\{0,1\}$}.

La relación de homotopía (relativa) en $\Omega(X,x_0)$ es de equivalencia.
Al conjunto de clases de homotopía de lazos lo denotaremos por
\[
\pi_1(X,x_0) = \Omega(X,x_0)/\!\sim.
\]
A la clase de homotopía de un lazo $\alpha$ lo denotaremos por $[\alpha]$.

Sea $f:(X,x_0)\to(Y,y_0)$ una función continua, es decir $f:X\to Y$ tal que
$f(x_0)=y_0$. Definimos
\[
f_* : \pi_1(X,x_0)\to\pi_1(Y,y_0), \qquad [\alpha]\mapsto[f\circ\alpha].
\]
Se verifica que:
\begin{enumerate}[label=\arabic*.]
\item $\mathrm{id}_*[\alpha]=[\alpha]$.
\item Si $g:(Y,y_0)\to(Z,z_0)$ es continua, entonces $(g\circ f)_*=g_*\circ f_*$.
\end{enumerate}
Por lo tanto, $\pi_1(X,x_0)$ es un invariante topológico (para espacios basados).
\end{definition}

\begin{definition}[Homotopía]
\label{def:homotopia-func}
Sean $f,g:X\to Y$ funciones continuas. Decimos que $f$ es \emph{homotópica}
a $g$, denotado $f\simeq g$, si existe una función continua $H:X\times I\to Y$
tal que $H(x,0)=f(x)$ y $H(x,1)=g(x)$. A tal función $H$ se le llama una
\emph{homotopía}.

Ya que la homotopía es una relación de equivalencia en $\mathrm{Map}(X,Y)$,
denotaremos por $[X,Y]$ al conjunto de clases de homotopía de funciones
continuas de $X$ a $Y$.
\end{definition}

\begin{theorem}[Grupo fundamental]
\label{thm:pi1-grupo}
El conjunto $\pi_1(X,x_0)$ es un grupo con la operación
\[
[\alpha]\cdot[\beta]=[\alpha*\beta], \qquad \text{en donde} \qquad
(\alpha*\beta)(s)=\begin{cases}\alpha(2s) & 0\le s\le 1/2\\ \beta(2s-1) & 1/2\le s\le 1.\end{cases}
\]
El elemento identidad es la aplicación constante $c_{x_0}(s)=x_0$ para toda
$s\in I$ y la inversa de un elemento $\alpha\in\Omega(X,x_0)$ es
$\overline\alpha(s)=\alpha(1-s)\in\Omega(X,x_0)$.
\end{theorem}

\begin{proof}
Primero, $*$ está bien definida en clases de homotopía: si
$H:\alpha\simeq\alpha'$ y $K:\beta\simeq\beta'$ (relativas a $\{0,1\}$),
entonces $(s,t)\mapsto H(2s,t)$ para $s\le1/2$ y $(s,t)\mapsto K(2s-1,t)$
para $s\ge1/2$ define una homotopía $\alpha*\beta\simeq\alpha'*\beta'$.

\emph{Asociatividad.} Los caminos $(\alpha*\beta)*\gamma$ y
$\alpha*(\beta*\gamma)$ recorren $\alpha,\beta,\gamma$ en el mismo orden,
solo que en intervalos de distinta longitud: el primero usa
$[0,\tfrac14],[\tfrac14,\tfrac12],[\tfrac12,1]$ y el segundo
$[0,\tfrac12],[\tfrac12,\tfrac34],[\tfrac34,1]$. La función continua y
creciente por tramos
\[
\varphi(s)=\begin{cases}
s/2 & 0\le s\le1/2\\
s-1/4 & 1/2\le s\le3/4\\
2s-1 & 3/4\le s\le1
\end{cases}
\]
satisface $((\alpha*\beta)*\gamma)\circ\varphi=\alpha*(\beta*\gamma)$, y
$H(s,t)=((\alpha*\beta)*\gamma)\big((1-t)s+t\varphi(s)\big)$ es una
homotopía relativa a $\{0,1\}$ entre ambos caminos.

\emph{Elemento neutro.} La homotopía
$H(s,t)=\alpha\!\left(\dfrac{2s}{2-t}\right)$ para $s\le(2-t)/2$ y
$H(s,t)=x_0$ para $s\ge(2-t)/2$ es continua, verifica $H(\cdot,0)=\alpha*c_{x_0}$
y $H(\cdot,1)=\alpha$, y es constante en $s\in\{0,1\}$; así
$\alpha*c_{x_0}\simeq\alpha$. Análogamente $c_{x_0}*\alpha\simeq\alpha$.

\emph{Inverso.} La homotopía
$H(s,t)=\begin{cases}\alpha(2s) & 2s\le1-t\\ \alpha(1-t) & 1-t\le2s\le1+t\\ \alpha(2-2s) & 2s\ge1+t\end{cases}$
satisface $H(\cdot,0)=\alpha*\overline\alpha$ y $H(\cdot,1)=c_{x_0}$, y es
constante en $s\in\{0,1\}$; así $\alpha*\overline\alpha\simeq c_{x_0}$, y
simétricamente $\overline\alpha*\alpha\simeq c_{x_0}$.
\end{proof}

Consideremos a la pareja $(X,A)$ en donde $A\subset X$ es subespacio y
funciones $f:(X,A)\to(Y,B)$ (es decir, $f:X\to Y$ tal que $f(A)\subset B$).

\begin{definition}[Homotopía de parejas]
\label{def:homotopia-parejas}
Decimos que $f,g:(X,A)\to(Y,B)$ son homotópicos si existe una aplicación
continua $H:(X\times I, A\times I)\to(Y,B)$ tal que $H(x,0)=f(x)$ y
$H(x,1)=g(x)$.

A las clases de homotopía de funciones de parejas las denotamos por
$[(X,A),(Y,B)]$. De esta manera,
\[
\pi_1(X,x_0)=[(I,\partial I),(X,x_0)].
\]
\end{definition}

\begin{definition}[Homotopía superior]
\label{def:homotopia-superior}
El \emph{grupo de homotopía} de dimensión $q$ para un espacio $X$ basado se
define como
\[
\pi_q(X,x_0)=[(I^q,\partial I^q),(X,x_0)].
\]

La estructura de grupo es la siguiente: sean $\alpha,\beta:I^q\to X$.
Definimos
\[
\alpha*\beta(s_1,s_2,\dots,s_q)=\begin{cases}
\alpha(2s_1,s_2,\dots,s_q) & 0\le s_1\le 1/2\\
\beta(2s_1-1,s_2,\dots,s_q) & 1/2\le s_1\le 1
\end{cases}
\]
Definimos $\overline\alpha(s_1,s_2,\dots,s_q)=\alpha(1-s_1,s_2,\dots,s_q)$ y
$[\alpha][\beta]=[\alpha*\beta]$.
\end{definition}

\begin{proposition}
\label{prop:fpushforward-hom}
Sea $f:(X,x_0)\to(Y,y_0)$ una función continua. Entonces, la aplicación
\[
f_*:\pi_q(X,x_0)\to\pi_q(Y,y_0), \qquad [\alpha]\mapsto[f\circ\alpha]
\]
es un homomorfismo de grupos.
\end{proposition}

\begin{proof}
La aplicación está bien definida: si $H:\alpha\simeq\alpha'$ (relativa a
$\partial I^q$), entonces $f\circ H:f\circ\alpha\simeq f\circ\alpha'$. Es
morfismo de grupos pues, directamente de la definición de $*$ en
coordenadas,
\[
f\circ(\alpha*\beta)=(f\circ\alpha)*(f\circ\beta),
\]
ya que ambos lados aplican $f$ tras evaluar en $\alpha$ o $\beta$ según la
misma partición de $I^q$ en las coordenadas $s_1\le1/2$ y $s_1\ge1/2$. Por
lo tanto $f_*([\alpha][\beta])=f_*([\alpha])f_*([\beta])$.
\end{proof}

\begin{theorem}
\label{thm:pi_q-sphere}
Para $q\ge1$, se cumple que
\[
\pi_q(X,x_0)=[(S^q,1),(X,x_0)].
\]
\end{theorem}

\begin{proof}
La aplicación cociente $q:I^q\to I^q/\partial I^q$ identifica todo
$\partial I^q$ a un punto, y $I^q/\partial I^q$ es homeomorfo a $S^q$
(enviando el punto colapsado al polo norte $1\in S^q$) mediante un
homeomorfismo $\phi$. Dado $\alpha:(I^q,\partial I^q)\to(X,x_0)$, ya que
$\alpha$ es constante en $\partial I^q$, factoriza (de manera única, y
continua por la propiedad universal del cociente) como
$\alpha=\overline\alpha\circ q$ para una única $\overline\alpha:I^q/\partial I^q\to X$
con $\overline\alpha(\ast)=x_0$, donde $\ast$ es el punto colapsado.
Componiendo con $\phi^{-1}$ se obtiene una biyección
\[
\{\alpha:(I^q,\partial I^q)\to(X,x_0)\}\ \longleftrightarrow\ \{\beta:(S^q,1)\to(X,x_0)\}, \qquad \alpha\mapsto\overline\alpha\circ\phi^{-1},
\]
que es compatible con las homotopías relativas al borde correspondientes
en cada lado (una homotopía $I^q\times I\to X$ constante en
$\partial I^q\times I$ desciende, por la misma propiedad universal, a una
homotopía $S^q\times I\to X$ constante en $\{1\}\times I$, y
recíprocamente). Por lo tanto induce una biyección
$\pi_q(X,x_0)\cong[(S^q,1),(X,x_0)]$, la cual es de hecho isomorfismo de
grupos si se transporta la operación de concatenación de
$\pi_q(X,x_0)$ a través de $\phi$.
\end{proof}

\begin{definition}[Tipo de homotopía]
\label{def:tipo-homotopia}
Dos espacios $X$ y $Y$ tienen el \emph{mismo tipo de homotopía} si existen
aplicaciones continuas $f:X\to Y$ y $g:Y\to X$ tales que
\[
f\circ g\simeq\mathrm{id}_Y \qquad\text{y}\qquad g\circ f\simeq\mathrm{id}_X.
\]
Esta relación se denota $X\simeq Y$ y las aplicaciones $f$ y $g$ se llaman
\emph{equivalencias de homotopía}.
\end{definition}

\begin{example}
\label{ex:tipo-homotopia}
\
\begin{enumerate}[label=\arabic*.]
\item Los espacios $X=S^1$ y $Y=S^1\cup([1,2]\times\{0\})$ tienen el mismo
tipo de homotopía y no son homeomorfos.
\item Los espacios $\RR^2-\{0\}$ y $S^1$ tienen el mismo tipo de homotopía.
En efecto, consideremos la inclusión canónica $\iota:S^1\to\RR^2-\{0\}$ y la
aplicación $g:\RR^2-\{0\}\to S^1$ definida por $g(x)=x/|x|$. Se tiene que
$g\circ\iota=\mathrm{id}_{S^1}$ y $H:\iota\circ g\simeq\mathrm{id}_{\RR^2-\{0\}}$
por la homotopía
\[
H:(\RR^2-\{0\})\times[0,1]\to\RR^2-\{0\}, \qquad (x,t)\mapsto tx+(1-t)g(x).
\]
\end{enumerate}
\end{example}

\begin{definition}[Espacio contraíble]
\label{def:contraible}
Un espacio $X$ es \emph{contraíble} si la aplicación identidad sobre $X$ es
homótopa a una aplicación constante.
\end{definition}

\begin{example}
\label{ex:contraibles}
\begin{itemize}
\item Un subconjunto convexo $C$ de $\RR^n$ es contraíble. La homotopía
entre la identidad y la aplicación constante sobre el punto $x_0\in C$ es la
aplicación continua
\[
H:C\times[0,1]\to C, \qquad (x,t)\mapsto(1-t)x_0+tx.
\]
\item La bola $D^n$ y el espacio $\RR^n$ son contraíbles para $n\ge1$.
\end{itemize}
\end{example}

\begin{definition}[Retractos]
\label{def:retractos}
\begin{enumerate}[label=\arabic*.]
\item Un subconjunto $A$ de un espacio $X$ es un \emph{retracto} de $X$ si
existe una aplicación continua $r:X\to A$ tal que $r\circ\iota=\mathrm{id}_A$,
en donde $\iota:A\to X$ es la inyección canónica. La aplicación $r$ se llama
\emph{retracción}.
\item Un subconjunto $A$ de $X$ es un \emph{retracto por deformación} si
existe una retracción $r$ tal que $\iota\circ r\simeq\mathrm{id}_X$. En este
caso, las aplicaciones $\iota$ y $r$ son equivalencias de homotopía.
\end{enumerate}

En particular, un espacio $X$ es contraíble si, y solamente si, existe un
punto $x_0\in X$ tal que $\{x_0\}$ es un retracto por deformación de $X$. Se
tiene que
\[
\{0\}\simeq D^n\simeq\RR^n
\]
y también
\[
\RR^{n+1}-\{0\}\simeq S^n\simeq S^n\times[0,1].
\]
\end{definition}

\begin{example}
\label{ex:S3-S1}
$S^3-S^1\simeq S^1$. Escribimos
\[
S^3=\{(z_1,z_2)\in\CC^2 \mid |z_1|^2+|z_2|^2=1\}, \qquad
S^1=\{(z,0)\mid z\in S^1\}.
\]
El círculo $\{(0,z)\mid z\in S^1\}$ es un retracto por deformación de
$S^3-S^1$ gracias a la homotopía
\[
H:(S^3-S^1)\times[0,1]\to S^3-S^1, \qquad
((z_1,z_2),t)\mapsto\left(tz_1,\ \sqrt{1-t^2|z_1|^2}\cdot\frac{z_2}{|z_2|}\right).
\]
\end{example}

\subsection*{Ejercicios}

\begin{enumerate}[label=\arabic*.,leftmargin=*]

\item (\textbf{Propiedades de $\pi_1$})
\begin{enumerate}[label=\alph*)]
\item Sea $\alpha\in\Omega(X,x_0,x_1)$. El morfismo
\[
\alpha_*:\pi_1(X,x_0)\to\pi_1(X,x_1), \qquad [\sigma]\mapsto[\overline\alpha\sigma\alpha],
\]
es un isomorfismo de grupos.
\item Si $\alpha,\beta\in\Omega(X,x_0,x_1)$ con $\alpha\ne\beta$, los
isomorfismos $\alpha_*$ y $\beta_*$ son conjugados.
\item Si $X$ es conexo por arcos, y $x_0,x_1\in X$, entonces $\pi_1(X,x_0)$
y $\pi_1(X,x_1)$ son isomorfos.
\item Sean $(X,x_0)$ y $(Y,y_0)$ dos espacios. Entonces
\[
\pi_1(X\times Y,(x_0,y_0))\cong\pi_1(X,x_0)\times\pi_1(Y,y_0).
\]
\item Una aplicación continua $f:X\to Y$ induce un morfismo de grupos
\[
f_*:\pi_1(X,x)\to\pi_1(Y,f(x)), \qquad [\sigma]\mapsto[f\circ\sigma].
\]
\item Sean $f,g:X\to Y$ dos aplicaciones continuas homótopas. Para todo
$x\in X$, existe un isomorfismo de grupos $u:\pi_1(Y,g(x))\to\pi_1(Y,f(x))$
tal que el siguiente diagrama conmuta:
\[
\begin{tikzcd}
& \pi_1(Y,g(x)) \arrow[dd,"u"] \\
\pi_1(X,x) \arrow[ur,"g_*"] \arrow[dr,swap,"f_*"] & \\
& \pi_1(Y,f(x))
\end{tikzcd}
\]
Además, si $f$ y $g$ son homótopos relativos a $\{x\}$, entonces $f_*=g_*$.
\item Si $f:X\to Y$ es una equivalencia de homotopía, entonces
$f_*:\pi_1(X,x)\to\pi_1(Y,f(x))$ es un isomorfismo de grupos.
\item Todo grupo fundamental de un espacio contraíble es trivial (todo
espacio contraíble es simplemente conexo).
\item La inclusión $i:S^n\hookrightarrow\RR^{n+1}-\{0\}$ induce un
isomorfismo $i_*$.
\item Sea $r:X\to A$ una retracción, y sea $x\in A$. Sea $i:A\hookrightarrow X$
la inclusión de $A$ en $X$. Entonces, el morfismo $r_*:\pi_1(X,x)\to\pi_1(A,x)$
es sobreyectivo y el morfismo $i_*:\pi_1(A,x)\to\pi_1(X,x)$ es inyectivo.
\item Sea $r:X\to A$ un retracto por deformación, y sea $x\in A$. Entonces,
el morfismo $i_*:\pi_1(A,x)\to\pi_1(X,x)$ inducido por la inclusión
$i:A\hookrightarrow X$ es un isomorfismo.
\end{enumerate}

\noindent\textbf{Solución.}
\begin{enumerate}[label=\alph*)]
\item Veamos que $\alpha_*$ es morfismo de grupos. Sean
$\sigma_1,\sigma_2\in\Omega(X,x_0)$, entonces
\[
\alpha_*([\sigma_1\sigma_2])=[\overline\alpha\sigma_1\sigma_2\alpha]
=[(\overline\alpha\sigma_1\alpha)(\overline\alpha\sigma_2\alpha)]
=[\overline\alpha\sigma_1\alpha][\overline\alpha\sigma_2\alpha]
=\alpha_*([\sigma_1])\alpha_*([\sigma_2]).
\]
Para mostrar que es isomorfismo, encontremos una inversa para $\alpha_*$.
Si ponemos $\beta=\overline\alpha$, definimos $\beta_*=\overline\alpha_*$. Sea
$[\sigma_{x_1}]\in\pi_1(X,x_1)$. Entonces,
\[
\alpha_*\circ\beta_*([\sigma_{x_1}])=\alpha_*([\overline\beta\sigma_{x_1}\beta])
=\alpha_*([\alpha\sigma_{x_1}\overline\alpha])
=[\overline\alpha(\alpha\sigma_{x_1}\overline\alpha)\alpha]
=[\sigma_{x_1}].
\]
Ahora, tomemos $[\sigma_{x_0}]\in\pi_1(X,x_0)$. Entonces,
\[
\beta_*\circ\alpha_*([\sigma_{x_0}])=\beta_*([\overline\alpha\sigma_{x_0}\alpha])
=\beta_*([\overline\alpha\sigma_{x_0}\alpha])
=[\overline\beta(\overline\alpha\sigma_{x_0}\alpha)\beta]
=[\sigma_{x_0}].
\]
Por lo tanto, $\alpha_*$ es isomorfismo de grupos.
\item Se tiene que, si $[\sigma]\in\pi_1(X,x_0)$,
\[
\beta_*([\sigma])=[\overline\beta\sigma\beta]
=[\overline\beta\alpha\overline\alpha\sigma\alpha\overline\alpha\beta]
=[\overline\beta\alpha][\overline\alpha\sigma\alpha][\overline\alpha\beta]
=[\overline\beta\alpha][\overline\alpha\sigma\alpha][\overline\beta\alpha]^{-1}
=[\overline\beta\alpha]\alpha_*([\sigma])[\overline\beta\alpha]^{-1}.
\]
\item Es una consecuencia directa del primer inciso.
\item Sean $p:X\times Y\to X$ y $q:X\times Y\to Y$ proyecciones. Entonces,
se induce una aplicación
\[
\Phi=(p_*,q_*):\pi_1(X\times Y,(x_0,y_0))\to\pi_1(X,x_0)\times\pi_1(Y,y_0),
\qquad [\alpha]\mapsto(p_*[\alpha],q_*[\alpha])=([p\circ\alpha],[q\circ\alpha]).
\]
Veamos que $\Phi$ está bien definida y es morfismo de grupos.

$\Phi$ es sobreyectiva. Sea $\beta\in\Omega(X,x_0)$ y $\gamma\in\Omega(Y,y_0)$.
Definimos
\[
\alpha:I\to X\times Y, \qquad t\mapsto(\beta(t),\gamma(t)).
\]
Entonces $\alpha\in\Omega(X\times Y,(x_0,y_0))$ y
$\Phi([\alpha])=([p\circ\alpha],[q\circ\alpha])=([\beta],[\gamma])$.

$\Phi$ es inyectiva. Supongamos que $\alpha\in\Omega(X\times Y,(x_0,y_0))$
tal que $\Phi([\alpha])=([c_{x_0}],[c_{y_0}])$. Esto significa que
$G:p\circ\alpha\simeq c_{x_0}$ y $H:q\circ\alpha\simeq c_{y_0}$. Entonces,
defínase
\[
F:I\times I\to X\times Y, \qquad (s,t)\mapsto(G(s,t),H(s,t)).
\]
Se cumple que $F(s,0)=(p\circ\alpha,q\circ\alpha)$, $F(s,1)=(c_{x_0},c_{y_0})=e_{(x_0,y_0)}$
y $F(0,t)=F(1,t)=(x_0,y_0)$. Por lo tanto, $F:\alpha\simeq c_{(x_0,y_0)}$.
\item Si $\alpha\in\Omega(X,x_0,x_1)$, entonces $f\circ\alpha\in\Omega(Y,f(x_0),f(x_1))$.
La composición $\alpha\mapsto f\circ\alpha$ es compatible con la homotopía
(relativa a los extremos) de caminos: $\alpha\simeq\beta\Rightarrow
f\circ\alpha\simeq f\circ\beta$, y con la composición:
$f\circ(\alpha\beta)=(f\circ\alpha)(f\circ\beta)$.
\item Sea $H:f\simeq g$. Sea $\gamma(s)=H_s(x)=H(x,s)\in\Omega(Y,f(x),g(x))$.
Por un inciso anterior, se tiene un isomorfismo
\[
u:\pi_1(Y,g(x))\to\pi_1(Y,f(x)), \qquad [\sigma]\mapsto[\overline\gamma\sigma\gamma].
\]
Por demostrar que $f_*=u\circ g_*$. Para esto, basta demostrar que, para
todo $\alpha\in\Omega(X,x)$, se tiene la homotopía
$(\overline\gamma(g\circ\alpha))\gamma\simeq f\circ\alpha$ (relativo a
extremidades). Para todo $s\in[0,1]$, consideremos la aplicación
$\gamma_s:t\mapsto H(x,st)$. La aplicación
$(s,t)\mapsto((\overline\gamma_s(H_s\circ\alpha))\gamma_s)(t)$ es continua.
En $s=0$, vale $(c_{f(x)}(f\circ\alpha))c_{f(x)}$, que es homótopo a
$f\circ\alpha$. En $s=1$, vale $(\overline\gamma(g\circ\alpha))\gamma$.
\item Sea $g:Y\to X$ una aplicación continua tal que $f\circ g$ y $g\circ f$
son homótopos a la identidad. Por el inciso anterior y por funtorialidad,
$f_*\circ g_*$ y $g_*\circ f_*$ son isomorfismos de grupos.
\item Se sigue del inciso anterior.
\item Se tiene debido a que $S^n\simeq\RR^{n+1}-\{0\}$.
\item La aplicación $r\circ i$ coincide con $\mathrm{id}:A\to A$. Por lo
tanto, $r_*\circ i_*=\mathrm{id}_{\pi_1(A,x)}$. Por consiguiente, $r_*$ es
sobre y $i_*$ es inyectivo.
\item Ya que $i\circ r\simeq\mathrm{id}_X$, se sigue que
$i_*\circ r_*=\mathrm{id}_{\pi_1(X,x)}$. Por lo tanto, $i$ es sobre. Por el
inciso anterior, $i_*$ es isomorfismo.
\end{enumerate}

\item (\textbf{Grupo fundamental del círculo}) Sea $\alpha:I\to S^1$ un lazo
de base $1$ en $S^1$, y sea $x_0\in\ZZ\subset\RR$ un punto. Existe un camino
$\tilde\alpha:I\to\RR$ tal que $\exp\circ\tilde\alpha=\alpha$ y
$\tilde\alpha(0)=x_0$, en donde $\exp:\RR\to S^1$ es la aplicación definida
por $t\mapsto\exp(2\pi it)$. La diferencia $d(\alpha):=\tilde\alpha(1)-\tilde\alpha(0)$
es un número entero y no depende de la elección del punto de $\tilde\alpha$.
Se tiene la aplicación
\[
d:\pi_1(S^1,1)\to\ZZ, \qquad [\alpha]\mapsto\tilde\alpha(1)-\tilde\alpha(0).
\]
Dos lazos homótopos de base $1$ en $S^1$ tienen el mismo grado.
\begin{enumerate}[label=\alph*)]
\item La aplicación $d:\pi_1(S^1,1)\to\ZZ$ es un isomorfismo de grupos.
\item El círculo $S^1\subset D^2$ no es un retracto.
\item (\textbf{Teorema de Brouwer}) Toda aplicación continua $D^2\to D^2$
tiene a lo más un punto fijo.
\end{enumerate}

\noindent\textbf{Solución.}
\begin{enumerate}[label=\alph*)]
\item Primeramente, mostremos que $d:\pi_1(S^1,1)\to\ZZ$ es una biyección.
Para todo número entero $n$, podemos elegir un camino $\delta:I\to\RR$ que
conecta a $0$ y $n$, y $\exp\circ\delta$ es un lazo de base $1$ y de grado
$n$ (podemos poner también $\alpha^{(n)}(t)=\exp(2\pi int)$ que es un lazo de
grado $n$). Por lo tanto, la aplicación $d$ es sobreyectiva.

Mostremos que $d:\pi_1(S^1,1)\to\ZZ$ es inyectiva. En efecto, sea $n$ un
número entero, y sean $[\alpha_1]$ y $[\alpha_2]$ elementos de grado $n$ de
$\pi_1(S^1,1)$. Sean $\tilde\alpha_1$ y $\tilde\alpha_2$ los levantamientos
de $\alpha_1$ y $\alpha_2$, respectivamente, que conectan a $0$ con $n$.
Entonces $\tilde\alpha_1\simeq\tilde\alpha_2$: una homotopía está dada por
$H(s,t)=(1-t)\tilde\alpha_1(s)+t\tilde\alpha_2(s)$. Por consiguiente, la
aplicación $\exp\circ H$ certifica que los lazos $\alpha_1$ y $\alpha_2$ son
homótopos.

Mostremos ahora que $d:\pi_1(S^1,1)\to\ZZ$ es un morfismo de grupos.
Consideremos los lazos $\alpha_1,\alpha_2\in\Omega(S^1,1)$. Sea
$\tilde\alpha_1:I\to\RR$ el levantamiento de $\alpha_1$ tal que
$\tilde\alpha_1(0)=0$, y sea $\tilde\alpha_2:I\to\RR$ el levantamiento de
$\alpha_2$ tal que $\tilde\alpha_2(0)=\tilde\alpha_1(1)$. El grado
$d(\alpha_1\alpha_2)$ es igual a $\tilde\alpha_2(1)$. Por otra parte,
$d(\alpha_1)=\tilde\alpha_1(1)$ y
$d(\alpha_2)=\tilde\alpha_2(1)-\tilde\alpha_2(0)=\tilde\alpha_2(1)-\tilde\alpha_1(1)$.
Por lo tanto, $d(\alpha_1\alpha_2)=d(\alpha_1)+d(\alpha_2)$.
\item La bola $D^2$ es simplemente conexo. Si existiera $r:S^1\to D^2$,
entonces $r_*:\ZZ\to\{1\}$ sobreyectiva, lo cual no puede suceder.
\item Supongamos que $f:D^2\to D^2$ es una aplicación continua que no
admite ningún punto fijo. Entonces, para todo punto $x\in D^2$, existe un
número real $\lambda=\lambda_x\ge0$ tal que $\|x+\lambda(x-f(x))\|=1$: el
vector $x+\lambda(x-f(x))$ corresponde al punto de intersección del círculo
$S^1$ y el semirrayo que tiene al punto $f(x)$ por extremidad y pasa por
$x$. La aplicación que a todo punto $x\in D^2$ asocia $x+\lambda(x-f(x))$ es
continua. Además, la restricción de esta aplicación sobre $S^1$ es la
identidad $\mathrm{id}_{S^1}:S^1\to S^1$. Por lo tanto, esta aplicación es
un retracto de $D^2$ sobre $S^1$. Contradicción.
\end{enumerate}

\item Sea $\alpha:S^1\to X$ una aplicación continua tal que $\alpha(1)=x$.
Entonces, $[\alpha]=1$ en $\pi_1(X,x)$ si y solo si se extiende continuamente
en una aplicación continua $D^2\to X$.

\noindent\textbf{Solución.} Si $\alpha$ se extiende continuamente en
$f:D^2\to X$, si $i:S^1\to D^2$ es la inclusión; entonces, por
funtorialidad, el siguiente diagrama es conmutativo:
\[
\begin{tikzcd}
\pi_1(S^1,1) \arrow[rr,"\alpha_*"] \arrow[dr,swap,"i_*"] & & \pi_1(X,x) \\
& \pi_1(D^2,1) \arrow[ur,swap,"f_*"] &
\end{tikzcd}
\]
Si $\beta:[0,1]\to S^1$ es el lazo inducido por pasar al cociente, entonces
$\alpha_*([\beta])=[\alpha]$. Como $\pi_1(D^2,1)=\{1\}$, esto muestra el
resultado.

Recíprocamente, si $[\alpha]=1$, sea $H:S^1\times[0,1]\to X$, $\alpha\simeq c$.
Entonces, $H$ factoriza en una aplicación continua del espacio topológico
$Y=(S^1\times[0,1])/S^1\times\{1\}$ a $X$. Sin embargo $Y$ es homeomorfo al
disco.

\item (\textbf{Homotopía superior})
\begin{enumerate}[label=\alph*)]
\item Sean $f,g:(X,x_0)\to(Y,y_0)$ funciones continuas tales que $f\simeq g$.
Entonces $f_*=g_*$.
\item Si $(X,x_0)\simeq(Y,y_0)$, entonces $\pi_q(X,x_0)\cong\pi_q(Y,y_0)$.
\item Si $(X,x_0)$ es del mismo tipo de homotopía de un punto, entonces
$\pi_q(X,x_0)=0$ para toda $q$. En particular, si $X$ es contraíble,
$\pi_q(X,x_0)=0$ para toda $q$.
\item Sea $\{X_\alpha\}$ una familia de espacios topológicos conexos por
arcos, y considérese el producto $\prod_\alpha X_\alpha$. Probar que las
proyecciones $p_\alpha:\prod_\alpha X_\alpha\to X_\alpha$ inducen el
isomorfismo
\[
\pi_q\left(\prod_\alpha X_\alpha\right)\overset{\cong}{\to}\prod_\alpha\pi_q(X_\alpha).
\]
\end{enumerate}

\noindent\textbf{Solución.}
\begin{enumerate}[label=\alph*)]
\item Sea $[\alpha]\in\pi_q(X,x_0)$. Ya que $f\simeq g$, entonces
$g\circ\alpha\simeq f\circ\alpha$. Por lo tanto, $f_*[\alpha]=g_*[\alpha]$.
\item Por hipótesis, existen $f:(X,x_0)\to(Y,y_0)$ y $g:(Y,y_0,X,x_0)$
tales que $f\circ g\simeq\mathrm{id}_Y$ y $g\circ f\simeq\mathrm{id}_X$.
Entonces, $g_*\circ f_*=(g\circ f)_*=\mathrm{id}_{\pi_q(X,x_0)}$ y
$f_*\circ g_*=(f\circ g)_*=\mathrm{id}_{\pi_q(Y,y_0)}$. Por lo tanto, $f_*$
es isomorfismo.
\item Se sigue de los puntos anteriores.
\item Recordemos que $\pi_q\left(\prod_\alpha X_\alpha\right)=\left[S^q,\prod_\alpha X_\alpha\right]$.
Entonces, dar un elemento $\beta:S^q\to\prod_\alpha X_\alpha$ en
$[S^q,\prod_\alpha X_\alpha]$ es equivalente a dar elementos
$p_\alpha\circ\beta=\beta_\alpha:S^q\to X_\alpha$ en $[S^q,X_\alpha]$, para
cada $\alpha$. Entonces, por la propiedad universal del producto de grupos,
se tiene un morfismo de grupos
\[
p_*:\pi_q\left(\prod_\alpha X_\alpha\right)\to\prod_\alpha\pi_q(X_\alpha), \qquad
[\beta]\mapsto\{[\beta_\alpha]\}_\alpha.
\]
Supongamos que $[\beta]\in\ker p_*$. Entonces, $p_*([\beta])=c_\alpha$.
Entonces, para cada $\alpha$, existe una homotopía
$H_\alpha(x,t):S^q\times I\to X_\alpha$ entre $\beta_\alpha$ y $c_\alpha$.
Por lo tanto, por propiedad universal del espacio producto, existe una
homotopía
\[
H:S^q\times I\to\prod_\alpha X_\alpha, \qquad (x,t)\mapsto\{H_\alpha(x,t)\}_\alpha,
\]
entre $\beta$ y $e$.
\end{enumerate}

\item (\textbf{Simplemente conexo: Seifert--Van Kampen débil}).
\begin{enumerate}[label=\alph*)]
\item Sean $U$ y $V$ dos abiertos conexos por arcos de un espacio $X$ tales
que $X=U\cup V$ y tales que $U\cap V$ es conexo por arcos. Para todo
$x\in U\cap V$, si $i:U\to X$ y $j:V\to X$ son inclusiones, entonces
$i_*\pi_1(U,x)\cup j_*\pi_1(V,x)$ generan a $\pi_1(X,x)$. En particular, si
$U$ y $V$ son simplemente conexos de intersección no vacía, entonces $X$ es
simplemente conexo.
\item Para $n\ge2$, la esfera $S^n$ es simplemente conexo.
\item Para $n\ge1$, el espacio proyectivo complejo $\CC P^n$ es simplemente
conexo.
\end{enumerate}

\noindent\textbf{Solución.}
\begin{enumerate}[label=\alph*)]
\item Sea $\alpha$ un lazo en $x$. Por compacidad de $[0,1]$ y continuidad
de $\alpha$, existe $n\in\NN$ tal que $\alpha\left(\left[\frac{i}{n},\frac{i+1}{n}\right]\right)$
está contenido en $U$ o en $V$ para todo $i=0,\dots,n-1$. Para todo
$i=0,\dots,n$, sea $c_i$ un camino entre $x$ y $\alpha(i/n)$, contenido en
$U$, $V$, $U\cap V$ si $\alpha(i/n)$ pertenece a $U$, $V$, $U\cap V$
respectivamente, con $c_0$ y $c_n$ constantes. Denotamos por $\alpha_i$ al
camino $t\mapsto\alpha\left(\frac{i+t}{n}\right)$. Así, el lazo $\alpha$ en
$x$ es homótopo a
\[
(\overline{c_0}\alpha_0c_1)(\overline{c_1}\alpha_1c_2)\cdots(\overline{c_{n-2}}\alpha_{n-2}c_{n-1})(\overline{c_{n-1}}\alpha_{n-1}c_n).
\]
Para $i=0,\dots,n-1$, el camino $\overline{c_i}\alpha_ic_{i+1}$ es un lazo
en $x$ contenido en $U$ o en $V$. Por lo tanto $i_*\pi_1(U,x)\cup j_*\pi_1(V,x)$
genera a $\pi_1(X,x)$. Para la segunda aserción, si $U\cap V$ es no vacío,
entonces $X$ está conexo por arcos. Si $U$ y $V$ son simplemente conexos,
entonces cada $\overline{c_i}\alpha_ic_{i+1}$ es homótopo al lazo constante
en $x$. Por consiguiente, $\alpha$ es homótopo al lazo constante en $x$. Por
lo tanto, $X$ es simplemente conexo.
\item La esfera $S^n$ se escribe como una unión del abierto $U$
complementario al polo norte, y del abierto $V$ complementario al polo sur.
Los abiertos $U,V$ son contraíbles, y por tanto simplemente conexos. La
intersección $U\cap V$ se retracta por deformación sobre el ecuador
$S^{n-1}$, que es conexo por arcos si $n\ge2$.
\item Recordemos que para todo $n>0$,
\[
\CC P^n=(\CC^{n+1}-\{0\})/\CC^*.
\]
Se tiene que $\pi_1(\CC P^1)\cong\pi_1(S^2)=0$. Supongamos que $\CC P^n$ es
simplemente conexo. Si $f:S^{2n+1}\to\CC P^n$ es la proyección canónica,
entonces
\[
\CC P^{n+1}\cong\CC P^n\cup_f D^{2n+2}.
\]
Sea $U=\mathring D^{2n+2}$. Entonces, $U$ es un abierto de $\CC P^{n+1}$. El
complemento del origen de $D^{2n+2}$ es un abierto $V$ de $\CC P^{n+1}$.
Como $U$ es contraíble, y $V$ se retracta por deformación fuerte sobre
$\CC P^n$, por tanto simplemente conexo, entonces $\CC P^{n+1}$ es
simplemente conexo.
\end{enumerate}

\item (\textbf{Invarianza de la dimensión})
\begin{enumerate}[label=\alph*)]
\item Sea $n\ge0$ diferente de $1$. Entonces, $\RR^n$ no es homeomorfo a
$\RR^1$.
\item Sea $m\ge0$ un entero diferente de $2$. Entonces, $\RR^m$ no es
homeomorfo a $\RR^2$.
\end{enumerate}

\noindent\textbf{Solución.}
\begin{enumerate}[label=\alph*)]
\item Si $n=0$, entonces $\RR^n$ no es homeomorfo a $\RR^1$. Si $n>1$ y
$f:\RR^n\to\RR^1$ es un homeomorfismo, obtenemos un homeomorfismo entre
$\RR^n-\{0\}$ y $\RR^1-\{f(0)\}$, lo cual es imposible porque
$\RR^1-\{f(0)\}$ no es conexo y $\RR^n-\{0\}$ es conexo.
\item Si $m=0$, $\RR^m$ no es homeomorfo a $\RR^2$. Si $m>2$ y
$f:\RR^m\to\RR^2$ es un homeomorfismo, obtenemos un homeomorfismo entre
$\RR^m-\{0\}$ y $\RR^2-\{f(0)\}$, que es imposible debido a que
$\RR^2-\{f(0)\}$ tiene el mismo tipo de homotopía de $S^1$ (y por tanto
$\pi_1(\RR^2-\{f(0)\},*)\cong\ZZ$) y $\RR^m-\{0\}$ tiene el mismo tipo de
homotopía de $S^{m-1}$ (por tanto, $\RR^m-\{0\}$ es simplemente conexo).
Contradicción.
\end{enumerate}

\item (\textbf{Cálculos con van Kampen}) Calcular el grupo fundamental de
$X$ cuando
\begin{enumerate}[label=\alph*)]
\item $X=S^1\vee S^1$.
\item $X=S^1\vee S^2$.
\item $X=\RR P^2\vee\RR P^2$.
\end{enumerate}

\noindent\textbf{Solución.}
\begin{enumerate}[label=\alph*)]
\item El espacio $S^1\vee S^1$ se puede descomponer como la unión
$(S^1\vee(S^1-p))$ y $((S^1-q)\vee S^1)$. La intersección es contraíble y
por lo tanto
\[
\pi_1(S^1\vee S^1)\cong\pi_1(S^1)*\pi_1(S^1)\cong\ZZ*\ZZ.
\]
\item Sea $U=S^1\cup U_{x_0}$ en donde $U_{x_0}$ es una vecindad contráctil
de $x_0$ en $S^2$ y sea $V=S^2\cup V_{x_0}$ en donde $V_{x_0}$ es una
vecindad contráctil de $x_0$ en $S^1$. Entonces, $X=U\cup V$. Así, $U$ es
homótopo a $S^1$, $V$ es homótopo a $S^2$ y $U\cap V$ es la unión en un
punto de dos conjuntos contraíbles, y por tanto es simplemente conexo. Por
lo tanto,
\[
\pi_1(S^1\vee S^2)\cong\pi_1(U,*)*\pi_1(V,*)\cong\pi_1(S^1,*)*\pi_1(S^2,*)\cong\ZZ.
\]
\item Podemos descomponer a $\RR P^2\vee\RR P^2$ de la forma de una vecindad
homeomorfa a $\RR P^2\vee U$ y de una vecindad $U\vee\RR P^2$, en donde $U$
es una vecindad contráctil de un punto base de $\RR P^2$. Entonces, la
vecindad $U\vee U$ es contráctil. Por lo tanto,
\[
\pi_1(\RR P^2\vee\RR P^2)\cong\ZZ/2\ZZ*\ZZ/2\ZZ.
\]
\end{enumerate}

\item Si los espacios $(X,x_0)$ y $(Y,y_0)$ son correctamente basados
\footnote{El espacio $(X,x_0)$ es correctamente basado si el punto $x_0$
admite una vecindad abierta $V$ tal que existe una homotopía
$H:V\times[0,1]\to V$ que verifica $H(x,0)=x$, $H(x,1)=x_0=H(x_0,t)$ para
todo $t\in[0,1]$ y todo $x\in V$.}, con $X$ y $Y$ conexos por arcos,
entonces
\[
\pi_1(X\vee Y)\cong\pi_1(X)*\pi_1(Y).
\]

\noindent\textbf{Solución.} Denotemos por $V_{x_0}$ y $V_{y_0}$ a las
vecindades abiertas de $x_0$ y $y_0$ y por $H_{x_0}$ y $H_{y_0}$ a las
homotopías correspondientes. El espacio $X\vee Y$ admite una cubierta
abierta formada por los abiertos $U=X\vee V_{y_0}$ y $V=V_{x_0}\vee Y$. La
aplicación $K:U\times[0,1]\to U$ definida por $K(x,t)=x$ para $x\in X$ y
$K(x,t)=H_{y_0}(y,t)$ para $y\in V_{y_0}$ muestra que $X$ es un retracto por
deformación fuerte de $U$. De igual forma, $Y$ es un retracto por
deformación de $V$, y el punto $\{[x_0]=[y_0]\}$ es un retracto por
deformación de $U\cap V$. Deducimos entonces del teorema de Seifert y Van
Kampen
\[
\pi_1(X\vee Y)=\pi_1(U)*\pi_1(V)=\pi_1(X)*\pi_1(Y).
\]

\end{enumerate}

\section{Cubrientes y grupo fundamental}

\subsection{Cubrientes}

\begin{definition}[Cubriente]
\label{def:cubriente}
Una aplicación continua $p:X\to B$ es un \emph{cubriente} si existe una
cubierta abierta $\{U_\alpha\}$ de $B$ tal que para cada $\alpha$,
$p^{-1}(U_\alpha)$ es una unión disjunta de conjuntos abiertos en $X$ que,
bajo $p$, se mapean de manera homeomorfa en $U_\alpha$.

Equivalentemente, un cubriente es una aplicación continua $p:X\to B$ tal
que, para todo $b\in B$, existe una vecindad abierta de $b$, $U\ni b$, un
espacio discreto no vacío $F$ y un homeomorfismo $\Phi:p^{-1}(U)\to U\times F$
tal que el siguiente diagrama conmuta:
\[
\begin{tikzcd}
p^{-1}(U) \arrow[rr,"\Phi",dashed] \arrow[dr,swap,"p"] & & U\times F \arrow[dl,"\mathrm{pr}"] \\
& U &
\end{tikzcd}
\]
\end{definition}

\begin{theorem}
\label{thm:accion-propia-cubriente}
Sea $G$ un grupo discreto actuando (continuamente) libre y propiamente
sobre un espacio topológico separado. Entonces
\begin{enumerate}[label=\arabic*.]
\item para todo $x\in X$, existe una vecindad $V$ de $x$ tal que
$gV\cap V=\varnothing$ para todo $g\in G-\{e\}$;
\item cada órbita de $G$ es discreta;
\item la proyección canónica $X\to X/G$ es un cubriente.
\end{enumerate}
\end{theorem}

\begin{proof}
Aquí \emph{propiamente} significa, como se precisa en el Ejercicio 3 de
esta sección: para cada $x\in X$ existe una vecindad $U=U_x$ tal que
$gU\cap g'U=\varnothing$ para $g\ne g'$ en $G$; en particular, con $g'=e$,
$gU\cap U=\varnothing$ para $g\ne e$, que es precisamente la afirmación 1.

\emph{2.} Sea $x\in X$ y $V$ una vecindad como en 1. Para $g\ne e$,
$gV\cap V=\varnothing$, así que $gx\notin V$ salvo si $g=e$. Por lo tanto
$V\cap Gx=\{x\}$: la órbita $Gx$ es discreta como subespacio de $X$
(cada uno de sus puntos $gx$ tiene la vecindad $gV$ que solo la interseca
en $gx$, trasladando el argumento por $g$).

\emph{3.} Sea $p:X\to X/G$ la proyección. Por 1, para $x\in X$ existe $V$
con $gV\cap V=\varnothing$ para $g\ne e$; como $gV\cap g'V=g(V\cap g^{-1}g'V)=\varnothing$
para $g\ne g'$ (aplicando la misma propiedad a $g^{-1}g'\ne e$), los
trasladados $\{gV\}_{g\in G}$ son disjuntos dos a dos. Sea $U=p(V)$; por el
Ejercicio 3a, $p$ es abierta, así $U$ es abierto. Además
$p^{-1}(U)=\bigsqcup_{g\in G}gV$ (Ejercicio 3b), y $p|_{gV}:gV\to U$ es
continua, biyectiva y abierta (Ejercicio 3b), luego homeomorfismo. Así,
$U$ es una vecindad evenly covered de $p(x)$, y $p$ es un cubriente.
\end{proof}

\begin{corollary}
\label{cor:accion-finita-cubriente}
Sea $G$ un grupo finito (discreto) actuando libremente (por homeomorfismos)
sobre un espacio separado $X$. Entonces $X/G$ es separado, y la proyección
canónica $X\to X/G$ es un cubriente con $\#G$ hojas.
\end{corollary}

\begin{proof}
Basta verificar que la acción es propiamente discontinua en el sentido del
Teorema~\ref{thm:accion-propia-cubriente}, pues entonces éste da el
cubriente y la separación de $X/G$ se sigue de que las órbitas (fibras de
$p$) son finitas y cerradas en un espacio Hausdorff. Sea $x\in X$ y
$G=\{e=g_0,g_1,\dots,g_n\}$. Como $X$ es Hausdorff y la acción es libre,
los puntos $g_0x,\dots,g_nx$ son distintos dos a dos, así que existen
vecindades disjuntas $U_0,\dots,U_n$ de cada uno. Sea
$V=\bigcap_{j=0}^ng_j^{-1}U_j$, vecindad abierta de $x$. Entonces
$g_jV\subset U_j$ para cada $j$ (pues $V\subset g_j^{-1}U_j$), y como
$U_0,\dots,U_n$ son disjuntos dos a dos, también lo son
$g_0V,\dots,g_nV$; en particular $g_jV\cap V=g_jV\cap g_0V=\varnothing$
para $j\ne0$, es decir, $g_jV\cap V=\varnothing$ para $g_j\ne e$: la
condición de discontinuidad propia se cumple. Por el
Teorema~\ref{thm:accion-propia-cubriente}, $p$ es cubriente; como
$\#G<\infty$, cada fibra $p^{-1}(p(x))=Gx$ tiene $\#G$ elementos, es
decir, $\#G$ hojas.
\end{proof}

\begin{corollary}
\label{cor:subgrupo-discreto-cubriente}
Sea $H$ un subgrupo discreto de un grupo topológico separado $G$. Entonces,
$H$ es cerrado, $G/H$ es separado y la proyección canónica $G\to G/H$ es un
cubriente.
\end{corollary}

\begin{proof}
$H$ actúa sobre $G$ libremente por traslación derecha, $h\cdot g:=gh^{-1}$
(libre porque $G$ es grupo: $gh^{-1}=g\Rightarrow h=e$). Como $H$ es
discreto, existe una vecindad $W$ de $e_G$ con $W\cap H=\{e\}$; por
continuidad de la multiplicación, existe una vecindad simétrica $V\ni e_G$
($V=V^{-1}$) con $V\cdot V\subset W$. Para $g\in G$, la vecindad $gV$
satisface: si $h\ne e$ está en $H$ y $(gV)\cap(gVh^{-1})\ne\varnothing$,
existen $v_1,v_2\in V$ con $gv_1=gv_2h^{-1}$, es decir $h=v_2^{-1}v_1\in V^{-1}V=VV\subset W$,
así $h\in W\cap H=\{e\}$, contradicción. Por lo tanto $gV\cap(gV)h^{-1}=\varnothing$
para $h\ne e$, que es la condición de discontinuidad propia (con la acción
escrita multiplicativamente). Por el Teorema~\ref{thm:accion-propia-cubriente},
$p:G\to G/H$ es cubriente y, en particular, abierta; como las fibras de
$p$ (las clases $gH$) son cerradas cuando $H$ es cerrado, basta ver que
$H$ es cerrado: si $g\notin H$, por Hausdorff separamos $g$ de $e$ con
vecindades disjuntas, y una vecindad evenly covered de $p(g)$ disjunta de
$p(e)$ (que existe pues $p$ es abierta y $G/H$ hereda de $G$ suficiente
separación vía las vecindades $gV$ de arriba) muestra que $G\setminus H$
es abierto. Finalmente, $G/H$ es separado porque las órbitas de una acción
propiamente discontinua sobre un espacio Hausdorff son cerradas y
disjuntas, y $p$ es abierta.
\end{proof}

\subsection{Teoremas de levantamientos}

\begin{definition}[Levantamiento]
\label{def:levantamiento}
Sea $p:X\to B$ un cubriente y $f:Y\to B$ una aplicación continua. Un
\emph{levantamiento} de $f$ es una aplicación continua $\tilde f:Y\to X$ tal
que $p\circ\tilde f=f$. Es decir, el siguiente diagrama conmuta:
\[
\begin{tikzcd}
& X \arrow[d,"p"] \\
Y \arrow[ur,dashed,"\tilde f"] \arrow[r,swap,"f"] & B
\end{tikzcd}
\]
\end{definition}

\begin{proposition}[Unicidad de levantamientos]
\label{prop:unicidad-levantamiento}
Si $Y$ es conexo, entonces dos levantamientos de $f$ que coinciden en un
punto son iguales.
\end{proposition}

\begin{proof}
Sean $\tilde f_1,\tilde f_2:Y\to X$ dos levantamientos de $f$ con
$\tilde f_1(y_0)=\tilde f_2(y_0)$ para algún $y_0\in Y$. Sea
$A=\{y\in Y\mid\tilde f_1(y)=\tilde f_2(y)\}$; $y_0\in A$, así que
$A\ne\varnothing$. Basta mostrar que $A$ es abierto y cerrado en $Y$
conexo, de donde $A=Y$.

Sea $y\in Y$ y $U\ni f(y)$ una vecindad evenly covered, con
$p^{-1}(U)=\bigsqcup_iV_i$ y $p|_{V_i}:V_i\xrightarrow{\cong}U$. Sean
$V_1\ni\tilde f_1(y)$ y $V_2\ni\tilde f_2(y)$ las hojas que contienen a
cada levantamiento; por continuidad, existe una vecindad $W$ de $y$ con
$\tilde f_1(W)\subset V_1$ y $\tilde f_2(W)\subset V_2$.

Si $y\in A$, entonces $\tilde f_1(y)=\tilde f_2(y)$, así que $V_1=V_2=:V$
(la misma hoja contiene a ambos). En $W$, $\tilde f_1=(p|_V)^{-1}\circ f=\tilde f_2$
(pues ambos son levantamientos de $f$ con valores en $V$, y $p|_V$ es
inyectiva), luego $W\subset A$: $A$ es abierto.

Si $y\notin A$, entonces $\tilde f_1(y)\ne\tilde f_2(y)$; si $V_1=V_2$,
como $p|_{V_1}$ es inyectiva y $p(\tilde f_1(y))=p(\tilde f_2(y))=f(y)$, se
tendría $\tilde f_1(y)=\tilde f_2(y)$, contradicción; así $V_1\ne V_2$, y
como las hojas son disjuntas, $\tilde f_1(w)\ne\tilde f_2(w)$ para todo
$w\in W$ (pues $\tilde f_1(w)\in V_1$, $\tilde f_2(w)\in V_2$ y
$V_1\cap V_2=\varnothing$): $W\subset Y\setminus A$, luego
$Y\setminus A$ es abierto y $A$ es cerrado.
\end{proof}

\begin{theorem}[Levantamientos de caminos y homotopías]
\label{thm:levantamiento-caminos}
Sea $p:X\to B$ un cubriente.
\begin{itemize}
\item Sea $f:Y\to B$ una aplicación continua que admite un levantamiento
$\tilde f:Y\to X\in\Lev(f)$. Para toda aplicación continua
$H:Y\times[0,1]\to B$ tal que $H(\bullet,0)=f(\bullet)$, existe un único
levantamiento $\tilde H\in\Lev(H)$ tal que $\tilde H(\bullet,0)=\tilde f(\bullet)$:
\[
\begin{tikzcd}
Y \arrow[r,"\tilde f"] \arrow[d,hook,swap,"{\bullet\mapsto(\bullet,0)}"] & X \arrow[d,"p"] \\
Y\times[0,1] \arrow[r,swap,"H"] \arrow[ur,dashed,"\exists!\,\tilde H"] & B
\end{tikzcd}
\]
\item Para todo camino $\alpha:I\to B$ con $\alpha(0)=b$, para todo
$x\in p^{-1}(b)$, existe un único levantamiento $\tilde\alpha\in\Lev(\alpha)$
tal que $\tilde\alpha(0)=x$:
\[
\begin{tikzcd}
& X \arrow[d,"p"] \\
I \arrow[ur,dashed,"{\exists!\,\tilde\alpha,\ \tilde\alpha(0)=x}"] \arrow[r,swap,"{\alpha,\ \alpha(0)=b}"] & B
\end{tikzcd}
\]
\item Si $\alpha$ y $\beta$ son dos caminos homótopos, y si
$\tilde\alpha\in\Lev(\alpha)$ y $\tilde\beta\in\Lev(\beta)$ con
$\tilde\alpha(0)=\tilde\beta(0)$, entonces $\tilde\alpha(1)=\tilde\beta(1)$
y $\tilde\alpha\simeq\tilde\beta$.
\item Sean $x\in X$ y $b=p(x)$. El morfismo $p_*:\pi_1(X,x)\to\pi_1(B,b)$
es inyectivo.
\end{itemize}
\end{theorem}

\begin{proof}
\emph{Levantamiento de caminos (segundo punto).} Sea $\alpha:I\to B$ con
$\alpha(0)=b$ y $x\in p^{-1}(b)$. Cubrimos $B$ por abiertos evenly
covered; por continuidad de $\alpha$ y compacidad de $I$, el número de
Lebesgue de la cubierta $\{\alpha^{-1}(U)\}$ da una partición
$0=t_0<t_1<\cdots<t_n=1$ tal que $\alpha([t_{i-1},t_i])$ está contenido en
un abierto evenly covered $U_i$, para cada $i$. Construimos $\tilde\alpha$
inductivamente: supongamos definida en $[0,t_{i-1}]$ con
$\tilde\alpha(0)=x$; sea $V_i$ la hoja de $p^{-1}(U_i)$ que contiene a
$\tilde\alpha(t_{i-1})$, y definamos
$\tilde\alpha|_{[t_{i-1},t_i]}:=(p|_{V_i})^{-1}\circ\alpha|_{[t_{i-1},t_i]}$.
Esto es continuo en cada pieza y coincide en los puntos de empalme, luego
$\tilde\alpha$ es continua en $I$ (lema del pegado) y $p\circ\tilde\alpha=\alpha$.
La unicidad es la Proposición~\ref{prop:unicidad-levantamiento} aplicada a
$Y=I$ (conexo).

\emph{Levantamiento de homotopías (primer punto).} Sea
$H:Y\times[0,1]\to B$ con $H(\bullet,0)=f(\bullet)=p\circ\tilde f(\bullet)$.
Fijemos $y\in Y$. Como antes, la compacidad de $\{y\}\times I$ y un
argumento de Lebesgue dan una vecindad $N$ de $y$ en $Y$ y una partición
$0=t_0<\cdots<t_n=1$ tales que $H(N\times[t_{i-1},t_i])$ está contenido en
un abierto evenly covered $U_i^y$, para cada $i$ (se cubre $\{y\}\times I$
por abiertos producto $N_j\times(s_j,s_j')$ con $H(N_j\times(s_j,s_j'))\subset$
evenly covered, se extrae una subcubierta finita de $\{y\}\times I$ por
compacidad, se toma $N=\bigcap N_j$ y se refina la partición para que cada
$[t_{i-1},t_i]$ quede en algún $(s_j,s_j')$). Construimos
$\tilde H|_{N\times[0,1]}$ inductivamente en $i$, exactamente como en el
levantamiento de caminos pero ahora usando, en cada paso, la hoja de
$p^{-1}(U_i^y)$ que contiene a $\tilde H(\bullet,t_{i-1})|_N$ (definida en
todo $N$ por inducción, con valor inicial $\tilde f|_N$ en $t_0=0$) y
definiendo $\tilde H|_{N\times[t_{i-1},t_i]}:=(p|_{V_i})^{-1}\circ H|_{N\times[t_{i-1},t_i]}$,
que es continua y coincide en $N\times\{t_{i-1}\}$ con lo ya construido.

Esto define $\tilde H$ en una vecindad $N\times I$ de $\{y\}\times I$ para
cada $y\in Y$; por la Proposición~\ref{prop:unicidad-levantamiento}
aplicada a $\{y'\}\times I$ (conexo) para $y'\in N\cap N'$, dos de estas
construcciones locales coinciden en la intersección $(N\cap N')\times I$
de sus dominios (ambas son levantamientos de $H|_{(N\cap N')\times I}$
que coinciden con $\tilde f$ en $t=0$). Así, las construcciones locales se
pegan en una única función continua $\tilde H:Y\times I\to X$, que es el
levantamiento buscado; la unicidad global es de nuevo la
Proposición~\ref{prop:unicidad-levantamiento}, aplicada a $Y\times I$
conexo si $Y$ lo es, o directamente a cada $\{y\}\times I$ en general.

\emph{Tercer punto (caminos homótopos).} Sea $H:I\times I\to B$ una
homotopía relativa a $\{0,1\}$ de $\alpha$ a $\beta$, y sea
$\tilde\alpha\in\Lev(\alpha)$. Aplicando el primer punto con $Y=I$,
$f=\alpha$, $\tilde f=\tilde\alpha$, obtenemos un único levantamiento
$\tilde H:I\times I\to X$ con $\tilde H(\bullet,0)=\tilde\alpha$. Como
$H(0,t)=b$ para todo $t$ (relativa a $\{0,1\}$), el camino
$t\mapsto\tilde H(0,t)$ es un levantamiento del camino constante $c_b$
que en $t=0$ vale $\tilde\alpha(0)$; por unicidad del levantamiento de
caminos, $\tilde H(0,t)=\tilde\alpha(0)$ para todo $t$ (constante).
Análogamente, $t\mapsto\tilde H(1,t)$ es constante igual a
$\tilde\alpha(1)$. Así, $s\mapsto\tilde H(s,1)$ es un levantamiento de
$\beta$ con origen $\tilde H(0,1)=\tilde\alpha(0)$; por unicidad de
levantamiento de caminos, $\tilde H(\bullet,1)=\tilde\beta$ para el
levantamiento $\tilde\beta$ de $\beta$ con $\tilde\beta(0)=\tilde\alpha(0)$.
Así, $\tilde H$ es exactamente una homotopía relativa a $\{0,1\}$ (pues
$\tilde H(0,\bullet)$ y $\tilde H(1,\bullet)$ son constantes) entre
$\tilde\alpha=\tilde H(\bullet,0)$ y $\tilde\beta=\tilde H(\bullet,1)$. En
particular, $\tilde\beta(1)=\tilde H(1,1)=\tilde H(1,0)=\tilde\alpha(1)$,
ya que $t\mapsto\tilde H(1,t)$ es constante.

\emph{Cuarto punto (inyectividad de $p_*$).} Sea $[\alpha]\in\ker p_*$,
$\alpha\in\Omega(X,x)$. Entonces $p\circ\alpha\simeq c_b$ relativo a
$\{0,1\}$ en $B$. Por el tercer punto (aplicado a los caminos $p\circ\alpha$
y $c_b$, homótopos), el levantamiento de $p\circ\alpha$ de origen $x$
--- que es $\alpha$ mismo, por unicidad --- es homótopo (relativo a
$\{0,1\}$) al levantamiento de $c_b$ de origen $x$, que es $c_x$. Por lo
tanto $[\alpha]=1$ en $\pi_1(X,x)$, y $\ker p_*=\{1\}$.
\end{proof}

\begin{theorem}[Teorema de levantamiento]
\label{thm:teo-levantamiento}
Sean $p:X\to B$ un cubriente, $Y$ un espacio conexo y localmente conexo por
arcos, y $f:Y\to B$ una aplicación continua. Sean $y\in Y$, $b=f(y)$ y
$x\in p^{-1}(b)$. Entonces, existe un levantamiento $\tilde f:Y\to X$ de $f$
tal que $\tilde f(y)=x$ si, y solamente si,
\[
f_*\pi_1(Y,y)\subset p_*\pi_1(X,x).
\]
\end{theorem}

\begin{proof}
($\Rightarrow$) Si $\tilde f$ existe con $p\circ\tilde f=f$ y
$\tilde f(y)=x$, entonces por funtorialidad
$f_*=p_*\circ\tilde f_*:\pi_1(Y,y)\to\pi_1(B,b)$, luego
$f_*\pi_1(Y,y)=p_*(\tilde f_*\pi_1(Y,y))\subset p_*\pi_1(X,x)$.

($\Leftarrow$) Supongamos $f_*\pi_1(Y,y)\subset p_*\pi_1(X,x)$. Para
$y'\in Y$, elijamos (usando que $Y$ es conexo y localmente conexo por
arcos, luego conexo por arcos) un camino $\gamma$ de $y$ a $y'$; sea
$\tilde\gamma$ el levantamiento de $f\circ\gamma$ con origen $x$
(Teorema~\ref{thm:levantamiento-caminos}), y definamos
$\tilde f(y'):=\tilde\gamma(1)$.

\emph{Buena definición.} Si $\gamma'$ es otro camino de $y$ a $y'$,
entonces $\gamma\overline{\gamma'}$ es un lazo en $y$, así
$[f\circ(\gamma\overline{\gamma'})]=f_*[\gamma\overline{\gamma'}]\in f_*\pi_1(Y,y)\subset p_*\pi_1(X,x)$.
Por lo tanto, el levantamiento de $f\circ\gamma\circ\overline{f\circ\gamma'}$
con origen $x$ es un lazo en $x$ (por definición de $p_*\pi_1(X,x)$, el
levantamiento de un representante de una clase en la imagen de $p_*$
regresa a $x$); esto fuerza a que los levantamientos de $f\circ\gamma$ y
$f\circ\gamma'$ con origen $x$ terminen en el mismo punto (concatenando
los levantamientos y usando unicidad). Así $\tilde f(y')$ no depende del
camino elegido.

\emph{Es levantamiento.} Por construcción, $p(\tilde f(y'))=p(\tilde\gamma(1))=(f\circ\gamma)(1)=f(y')$.

\emph{Continuidad.} Sea $y'\in Y$ y $U\ni f(y')$ evenly covered, con hoja
$V\ni\tilde f(y')$ tal que $p|_V:V\xrightarrow{\cong}U$. Por continuidad de
$f$ y conexidad local por arcos de $Y$, existe una vecindad $N$ de $y'$,
conexa por arcos, con $f(N)\subset U$. Para $y''\in N$, tomando un camino
de $y$ a $y'$ seguido de uno de $y'$ a $y''$ dentro de $N$, la buena
definición da $\tilde f(y'')=(p|_V)^{-1}(f(y''))$ (el levantamiento, tras
llegar a $\tilde f(y')\in V$, permanece en $V$ porque $f(N)\subset U$ y
$V$ es la única hoja sobre $U$ que contiene a $\tilde f(y')$). Por lo
tanto $\tilde f|_N=(p|_V)^{-1}\circ f|_N$ es continua.

La unicidad de $\tilde f$ es la Proposición~\ref{prop:unicidad-levantamiento}
($Y$ conexo).
\end{proof}

\begin{corollary}
\label{cor:levantamiento-simplemente-conexo}
Si $p:X\to B$ es un cubriente, para toda aplicación continua $f:Y\to B$ tal
que $Y$ sea localmente conexo por arcos y simplemente conexo, existe un
único levantamiento $\tilde f:Y\to X$ tal que $\tilde f(y)=x$.
\end{corollary}

\begin{proof}
$Y$ simplemente conexo (en particular conexo) da $\pi_1(Y,y)=1$, así
$f_*\pi_1(Y,y)=1\subset p_*\pi_1(X,x)$ trivialmente. El
Teorema~\ref{thm:teo-levantamiento} da la existencia, y la
Proposición~\ref{prop:unicidad-levantamiento} la unicidad.
\end{proof}

\subsection{$G$-cubrientes y $\pi_1$}

\begin{theorem}
\label{thm:pi1-G-cubriente}
Sean $G$ un grupo discreto actuando propia y libremente sobre un espacio
topológico $X$ separado y conexo por arcos, $p:X\to X/G$ la proyección
canónica, $b\in X/G$ y $x\in p^{-1}(b)$. Para todo $[\alpha]\in\pi_1(X/G,b)$,
existe un único $g_\alpha\in G$ tal que $x[\alpha]=g_\alpha x$. La aplicación
\[
\pi_1(X/G,b)\to G, \qquad [\alpha]\mapsto g_\alpha
\]
es un morfismo de grupos, sobreyectivo, de núcleo $p_*\pi_1(X,x)$. En
particular, si $X$ es simplemente conexo, entonces $\pi_1(X/G,b)\cong G$.
\end{theorem}

\begin{proof}
Por definición, $p^{-1}(b)$ es la órbita de $x$ por $G$, esto muestra la
existencia de $g_\alpha$. La unicidad se sigue debido a que la acción es
libre. Es claro que $g_e=e$.

Sea $g\in G$ y $[\alpha]\in\pi_1(X/G,b)$. Entonces, $x[\alpha]$ es la
extremidad del levantamiento $\tilde\alpha$ de $\alpha$ de origen $x$ y
$g(x[\alpha])$ es la extremidad del levantamiento $g\tilde\alpha$ de origen
$gx$. Por unicidad, $g(x[\alpha])=(gx)[\alpha]$. En consecuencia, para todo
$[\alpha]$ y $[\alpha']$ en $\pi_1(X/G,b)$,
\[
g_{\alpha'\alpha}=x([\alpha'\alpha])=(g_{\alpha'}x)[\alpha]=g_{\alpha'}(g_\alpha x)=(g_{\alpha'}g_\alpha)x
\]
lo que muestra que $\alpha\mapsto g_\alpha$ es un morfismo de grupos.

La sobreyectividad se sigue del hecho de que $\pi_1(X/G,b)$ actúa
transitivamente sobre $p^{-1}(b)$. Además,
\[
\ker(\alpha\mapsto g_\alpha)=\Stab_{\pi_1(X/G,b)}(x)\,p_*\pi_1(X,x).\qedhere
\]
\end{proof}

\subsection{Cubriente universal}

\begin{definition}[Cubriente universal]
\label{def:cubriente-universal}
Sea $B$ conexo y localmente conexo por arcos. Un \emph{cubriente universal}
de $B$ es un cubriente $\pi:\tilde B\to B$ de espacio total $\tilde B$
conexo tal que para todo cubriente $p:X\to B$ de espacio $X$ conexo, para
todo $\tilde b\in\tilde B$ y $x\in X$ tales que $\pi(\tilde b)=p(x)$, existe
un morfismo de cubrientes $\phi:\tilde B\to X\in\Mor(\pi,p)$ tal que
$\phi(\tilde b)=x$.
\end{definition}

\begin{theorem}
\label{thm:cubriente-universal-galois}
Un cubriente universal de $B$ es de Galois.
\end{theorem}

\begin{proof}
Sea $\pi:\tilde B\to B$ universal, $b\in B$ y $x_1,x_2\in\pi^{-1}(b)$. Por
el Ejercicio 10a de esta sección, basta ver que $\Aut(\pi)$ actúa
transitivamente sobre $\pi^{-1}(b)$. Aplicando la propiedad universal de
$\pi$ (Definición~\ref{def:cubriente-universal}) al cubriente $X=\tilde B$
consigo mismo, con $\tilde b=x_1$ y $x=x_2$ (ambos sobre $b$), obtenemos
$\varphi\in\Mor(\pi,\pi)$ con $\varphi(x_1)=x_2$; simétricamente, hay
$\psi\in\Mor(\pi,\pi)$ con $\psi(x_2)=x_1$. Entonces $\psi\circ\varphi$ y
$\mathrm{id}_{\tilde B}$ son ambos levantamientos de $\pi$ a través de sí
mismo que coinciden en $x_1$; por la
Proposición~\ref{prop:unicidad-levantamiento} ($\tilde B$ conexo),
$\psi\circ\varphi=\mathrm{id}_{\tilde B}$, y simétricamente
$\varphi\circ\psi=\mathrm{id}_{\tilde B}$. Así $\varphi\in\Aut(\pi)$ y
$\varphi(x_1)=x_2$: la acción es transitiva.
\end{proof}

\begin{theorem}
\label{thm:simplemente-conexo-es-universal}
Si $p:X\to B$ es un cubriente, de espacio total $X$ simplemente conexo,
entonces $p$ es un cubriente universal.
\end{theorem}

\begin{proof}
Sea $p':X'\to B$ un cubriente con $X'$ conexo, y sean $\tilde x\in X$,
$x'\in X'$ con $p(\tilde x)=p'(x')$. Como $X$ es simplemente conexo,
$p_*\pi_1(X,\tilde x)=1\subset p'_*\pi_1(X',x')$ trivialmente; además $X$
es localmente conexo por arcos (todo cubriente de un espacio localmente
conexo por arcos lo es, vía las hojas evenly covered, que son
homeomorfas a abiertos de $B$). Por el
Teorema~\ref{thm:teo-levantamiento} (aplicado con $f=p:X\to B$ en el
papel de la aplicación a levantar, y $p'$ en el papel del cubriente),
existe $\phi:X\to X'$ con $p'\circ\phi=p$ y $\phi(\tilde x)=x'$, es decir,
$\phi\in\Mor(p,p')$. Esto es exactamente la propiedad universal de
$p$.
\end{proof}

\begin{theorem}
\label{thm:existencia-cubriente-universal}
Sea $B$ un espacio topológico separado, conexo y localmente conexo por
arcos. Entonces, $B$ admite un cubriente simplemente conexo si, y
solamente si, $B$ es semilocalmente simplemente conexo.

En particular, todo espacio separado, conexo, localmente contraíble admite
un cubriente universal.
\end{theorem}

\begin{proof}
($\Rightarrow$) Si $p:\tilde B\to B$ es un cubriente con $\tilde B$
simplemente conexo, todo $b\in B$ tiene una vecindad evenly covered $U$;
la inclusión $i:U\to B$ satisface $i_*\pi_1(U,b)\subset p_*\pi_1(\tilde B,\cdot)=1$
por el Ejercicio 1j (con $r$ la retracción de una hoja sobre $U$ vía
$p^{-1}$), así $i_*$ es el morfismo nulo: $B$ es semilocalmente
simplemente conexo.

($\Leftarrow$) Fijemos $b_0\in B$. Sea
\[
\tilde B:=\{[\gamma]\mid\gamma:I\to B\text{ camino con }\gamma(0)=b_0\}\big/(\text{homotopía rel }\{0,1\}),
\]
y $\pi:\tilde B\to B$, $\pi([\gamma])=\gamma(1)$. Para $U\subset B$
abierto, conexo por arcos, con $i_*\pi_1(U)\to\pi_1(B)$ nulo (un
\emph{abierto admisible}; estos existen alrededor de cada punto por
hipótesis, y se pueden tomar conexos por arcos ya que $B$ es localmente
conexo por arcos), y para $[\gamma]\in\tilde B$ con $\gamma(1)\in U$,
definamos
\[
U_{[\gamma]}:=\{[\gamma\cdot\eta]\mid\eta:I\to U\text{ camino con }\eta(0)=\gamma(1)\}.
\]
La familia $\{U_{[\gamma]}\}$, variando $U$ admisible y $[\gamma]$, forma
la base de una topología en $\tilde B$: dados dos tales conjuntos y un
punto en su intersección, el abierto admisible más pequeño contenido en
ambos $U$'s da un tercer básico contenido en la intersección.

Con esta topología, $\pi|_{U_{[\gamma]}}:U_{[\gamma]}\to U$ es una
biyección (sobreyectiva porque $U$ es conexo por arcos; inyectiva porque
si $\gamma\cdot\eta\simeq\gamma\cdot\eta'$ rel $\{0,1\}$ con extremos
$\eta(1)=\eta'(1)$, entonces, ya que $i_*\pi_1(U)=1$ en $\pi_1(B)$, dos
caminos en $U$ con los mismos extremos son homótopos rel extremos en $B$
--- de hecho ya en $U$, pues basta con que la clase de
$\eta\overline{\eta'}$ muera en $\pi_1(B)$ para concluir
$[\gamma\eta]=[\gamma\eta']$) y de hecho homeomorfismo (es abierta y
continua, verificándose directamente sobre la base). Como
$\pi^{-1}(U)=\bigsqcup_{[\gamma]}U_{[\gamma]}$ (unión sobre las distintas
clases $[\gamma]$ con $\gamma(1)\in U$, disjuntas porque
$U_{[\gamma]}\cap U_{[\gamma']}\ne\varnothing$ obliga, vía el mismo
argumento de $\pi_1(U)$ nulo, a $U_{[\gamma]}=U_{[\gamma']}$), $\pi$ es un
cubriente.

$\tilde B$ es conexo por arcos: dado $[\gamma]\in\tilde B$, el camino
$t\mapsto[\gamma_t]$ (donde $\gamma_t(s)=\gamma(ts)$) conecta
$[c_{b_0}]$ con $[\gamma]$ en $\tilde B$ (es continuo pues su composición
con $\pi$ es $\gamma$, y localmente coincide con la inversa de $\pi$ en
cartas $U_{[\gamma_s]}$).

Es simplemente conexo: sea $\tilde b_0:=[c_{b_0}]$ y sea
$[\delta]\in\pi_1(\tilde B,\tilde b_0)$, con $\delta$ un lazo. Sea
$\gamma:=\pi\circ\delta$, un lazo en $B$ basado en $b_0$. El camino
$t\mapsto[\gamma_t]$ construido arriba es, por unicidad de levantamientos
de caminos (Teorema~\ref{thm:levantamiento-caminos}), el único
levantamiento de $\gamma$ con origen $\tilde b_0$; como $\delta$ también
lo es, $\delta(t)=[\gamma_t]$ para todo $t$. Ahora, $\delta$ es un lazo,
así que $\delta(1)=\tilde b_0$, es decir $[\gamma]=[c_{b_0}]$
\emph{como puntos de $\tilde B$} --- lo cual, por la definición misma de
$\tilde B$ como cociente por homotopía relativa a $\{0,1\}$, significa
$\gamma\simeq c_{b_0}$ en $B$. Por el tercer punto del
Teorema~\ref{thm:levantamiento-caminos} (caminos homótopos levantan a
caminos homótopos con los mismos extremos), el levantamiento de $\gamma$
con origen $\tilde b_0$ --- que es $\delta$ --- es homótopo relativo a
$\{0,1\}$ al levantamiento de $c_{b_0}$ con origen $\tilde b_0$, que es
$c_{\tilde b_0}$. Por lo tanto $[\delta]=1$ en $\pi_1(\tilde B,\tilde b_0)$,
y $\pi_1(\tilde B,\tilde b_0)=1$.

La última afirmación (`en particular') se sigue porque todo espacio
localmente contraíble es semilocalmente simplemente conexo (cada punto
tiene una base de vecindades contráctiles, en particular simplemente
conexas, cuya inclusión induce el morfismo nulo en $\pi_1$).
\end{proof}

\begin{example}[Cubrientes universales]
\label{ex:cubrientes-universales}
\begin{itemize}
\item $\RR\to S^1: t\mapsto e^{2\pi it}$.
\item $\RR^n\to\mathbb T^n=\RR^n/\ZZ^n$, para $n\ge2$.
\item $S^n\to\RR P^n$, para $n\ge2$.
\item $S^{2n-1}\to L_m(\ell_1,\dots,\ell_n)$.
\end{itemize}
\end{example}

\subsection{Clasificación de cubrientes}

\begin{theorem}[Clasificación de cubrientes]
\label{thm:clasificacion-cubrientes}
Se tiene una biyección entre
\[
\{\text{subgrupos de } \Aut(\tilde\pi)\}/\!\cong \ \leftrightarrow\ \{\text{cubrientes conexos de } B\}/\!\cong.
\]
La biyección induce una biyección entre
\[
\{\text{subgrupos distinguidos de } \Aut(\tilde\pi)\}/\!\cong \ \leftrightarrow\ \{\text{cubrientes de Galois de } B\}/\!\cong.
\]
A su vez, para cada $n\in\NN$, se induce una biyección entre
\[
\{\text{subgrupos de índice } n \text{ de } \Aut(\tilde\pi)\}/\!\cong \ \leftrightarrow\ \{\text{cubrientes con } n \text{ hojas de } B\}/\!\cong.
\]
\end{theorem}

\begin{proof}
Fijemos el cubriente universal $\pi:\tilde B\to B$
(Teorema~\ref{thm:existencia-cubriente-universal}) y $\tilde b_0\in\tilde B$
sobre $b_0\in B$. Por el Teorema~\ref{thm:cubriente-universal-galois},
$\pi$ es de Galois, y como $\pi_*\pi_1(\tilde B,\tilde b_0)=1$, el
Ejercicio 10b de esta sección da un isomorfismo
$\Aut(\pi)\cong\pi_1(B,b_0)/1=\pi_1(B,b_0)$. En lo que sigue identificamos
$\Aut(\tilde\pi)$ con $\pi_1(B,b_0)$ vía este isomorfismo.

\emph{De subgrupos a cubrientes.} Sea $H\le\Aut(\pi)$. Como $\Aut(\pi)$
actúa libremente sobre $\tilde B$ (un automorfismo no trivial de un
cubriente conexo no fija ningún punto, por la
Proposición~\ref{prop:unicidad-levantamiento}: si fijara un punto,
coincidiría con $\mathrm{id}$) y propiamente (las hojas de $\pi$ separan
las órbitas de $\Aut(\pi)\subset\Aut(\pi)$, ya que $\Aut(\pi)$ permuta las
hojas de cada abierto evenly covered de $\pi$), también $H$ actúa libre y
propiamente sobre $\tilde B$; por el
Teorema~\ref{thm:accion-propia-cubriente}, $q_H:\tilde B\to\tilde B/H$ es
un cubriente, y $\tilde B/H\to B$ es un cubriente conexo de $B$ (composición
de cubrientes con espacio total conexo). Por el
Teorema~\ref{thm:pi1-G-cubriente} aplicado a la acción de $H$ sobre
$\tilde B$ (simplemente conexo), $\pi_1(\tilde B/H,q_H(\tilde b_0))\cong H$,
y bajo esta identificación el subgrupo $q_{H*}\pi_1(\tilde B/H,\cdot)$ de
$\pi_1(B,b_0)$ (visto vía la proyección $\tilde B/H\to B$) es exactamente
$H\le\pi_1(B,b_0)=\Aut(\pi)$.

\emph{De cubrientes a subgrupos.} Sea $p:X\to B$ un cubriente conexo,
$x_0\in p^{-1}(b_0)$. Como $\tilde B$ es simplemente conexo, el
Corolario~\ref{cor:levantamiento-simplemente-conexo} da un levantamiento
$q:\tilde B\to X$ de $\pi$ con $q(\tilde b_0)=x_0$; $q$ es a su vez un
cubriente (composición: $\pi=p\circ q$ es cubriente y $p$ es localmente
homeomorfismo, de donde $q$ hereda ser cubriente sobre cada hoja de $p$).
Sea $H:=\Stab_{\Aut(\pi)}(q)\!:=\{\varphi\in\Aut(\pi)\mid q\circ\varphi=q\}$.
Como $q$ es cubriente con $\tilde B$ simplemente conexo, la fibra
$q^{-1}(x_0)$ es exactamente la órbita de $\tilde b_0$ bajo $\Aut(\pi)$
(Teorema~\ref{thm:pi1-G-cubriente} identifica $\Aut(\pi)\cong\pi_1(B,b_0)$
actuando transitivamente sobre $\pi^{-1}(b_0)$, y $q$ colapsa esta órbita
según las fibras de $q$); se verifica que $q$ identifica $X$ con
$\tilde B/H$ (dos puntos de $\tilde B$ tienen la misma imagen bajo $q$ si y
solo si están en la misma órbita de $H$, por la propiedad universal de
$\pi$ aplicada para comparar levantamientos). Bajo la identificación de
$\Aut(\pi)$ con $\pi_1(B,b_0)$, $H$ corresponde a $p_*\pi_1(X,x_0)$ (la
correspondencia del párrafo anterior, en sentido inverso).

Las dos construcciones son inversas una de la otra: partiendo de $H$, el
cubriente $\tilde B/H\to B$ da de vuelta el subgrupo $H$; partiendo de un
cubriente $X$, el subgrupo construido y luego cocientado da de vuelta un
cubriente isomorfo a $X$ (ambos tienen el mismo subgrupo asociado
$p_*\pi_1(X,x_0)\le\pi_1(B,b_0)$, y dos cubrientes conexos con el mismo
subgrupo asociado en el mismo punto base son isomorfos por el
Teorema~\ref{thm:teo-levantamiento} aplicado en ambas direcciones, más la
Proposición~\ref{prop:unicidad-levantamiento} para la unicidad del
isomorfismo). Al variar el punto base $x_0$ dentro de la misma componente
conexa de $X$, el subgrupo asociado varía dentro de una clase de
conjugación en $\pi_1(B,b_0)$ (Ejercicio 9 de esta sección); por lo tanto
la construcción desciende a una biyección entre clases de isomorfismo de
cubrientes conexos de $B$ y clases de conjugación de subgrupos de
$\Aut(\tilde\pi)\cong\pi_1(B,b_0)$, es decir --- ya que estamos
enumerando subgrupos ``hasta $\cong$'', lo cual para subgrupos de un grupo
fijo significa hasta conjugación --- la biyección enunciada.

Los subgrupos distinguidos corresponden a los cubrientes de Galois: por
el Ejercicio 10a, $p:X\to B$ es regular (de Galois) si y solo si
$p_*\pi_1(X,x_0)$ es normal en $\pi_1(B,b_0)$, que es la condición de que
el subgrupo asociado $H$ sea distinguido en $\Aut(\tilde\pi)$.

Los subgrupos de índice $n$ corresponden a cubrientes con $n$ hojas: por
el Ejercicio 8d, $p$ tiene $n$ hojas si y solo si
$[\pi_1(B,b_0):p_*\pi_1(X,x_0)]=n$, que es exactamente
$[\Aut(\tilde\pi):H]=n$.
\end{proof}

\subsection*{Ejercicios}

\begin{enumerate}[label=\arabic*.,leftmargin=*]

\item Cualquier aplicación cubriente $p:X\to B$ es abierta.

\noindent\textbf{Solución.} Sea $U\subset X$ un abierto. Sea $b\in p(U)$ y
$x\in p^{-1}(b)\cap U$. Sea $V\ni b$, vecindad de $b$ tal que
$x\in p^{-1}(V)=\bigsqcup_{i\in I}U_i$. Entonces, existe algún $U_i$ en $X$
tal que $x\in U_i$. También, $U_i\cap U$ es abierto en $U_i$ y $p|_{U_i}$ es
un homeomorfismo entre $U_i$ y $V$. Obsérvese que $p(U_i\cap U)$ es abierto
de $V$. Ya que $b\in p(U_i\cap U)\subset p(U)$ se sigue que $p(U)$ es
abierto en $B$.

\item (\textbf{Cubrientes})
\begin{enumerate}[label=\alph*)]
\item Sean $p:X\to B$ y $p':X'\to B'$ dos cubrientes. Mostrar que la
aplicación $X\times X'\to B\times B'$ definida por
$(x,x')\mapsto(p(x),p'(x'))$ es un cubriente.
\item Sea $p:X\to B$ un cubriente y sea $f:X'\to B$ una aplicación
continua. Considérese el espacio
\[
X\times_BX'=\{(x,x')\in X\times X'\mid p(x)=f(x')\},
\]
y la aplicación $\overline p:X\times_BX'\to X'$ dada por
$\overline p(x,x')=x'$. Probar que $\overline p$ es cubriente.
\end{enumerate}

\noindent\textbf{Solución.}
\begin{enumerate}[label=\alph*)]
\item Llamemos $p\times p':X\times X'\to B\times B'$. Sean $\{U_\alpha\}$
una cubierta abierta de $B$ y $\{U'_\beta\}$ una cubierta abierta de $B'$
tales que, para cada $\alpha$ y $\beta$,
\[
p^{-1}(U_\alpha)=\bigsqcup_iV_{i,\alpha} \qquad\text{y}\qquad
p'^{-1}(U'_\beta)=\bigsqcup_jV'_{j,\beta},
\]
en donde $V_{i,\alpha}$ y $V'_{j,\beta}$ son abiertos de $X$ y $X'$,
respectivamente, y tales que $p|_{V_{i,\alpha}}:V_{i,\alpha}\cong U_\alpha$
y $p|_{V'_{j,\beta}}:V'_{j,\beta}\cong U'_\beta$.

Entonces, $\{U_\alpha\times U'_\beta\}$ es una cubierta abierta de
$B\times B'$. Además,
\[
(p\times p')^{-1}(U_\alpha\times U'_\beta)=\bigsqcup_{i,j}(V_{i,\alpha}\times V'_{j,\beta}).
\]
Se tiene que la familia $V_{i,\alpha}\times V'_{j,\beta}$ es ajena, abierta
y, además, $p\times p'|_{V_{i,\alpha}\times V'_{j,\beta}}:V_{i,\alpha}\times V'_{j,\beta}\cong U_\alpha\times U'_\beta$.
\item Consideremos el siguiente diagrama conmutativo:
\[
\begin{tikzcd}
X\times_BX' \arrow[r,"\overline f"] \arrow[d,swap,"\overline p"] & X \arrow[d,"p"] \\
X' \arrow[r,swap,"f"] & B
\end{tikzcd}
\]
Sea $\{U_\alpha\}$ una cubierta abierta de $B$ tal que
$p^{-1}(U_\alpha)=\bigsqcup_iV_{i,\alpha}$, con $V_{i,\alpha}$ abiertos de
$X$ y $p:V_{i,\alpha}\cong U_\alpha$. Así, $f^{-1}(U_\alpha)$ forma una
cubierta abierta de $X'$ y además
\[
(\overline p)^{-1}(f^{-1}(U_\alpha))=(\overline f)^{-1}(p^{-1}(U_\alpha))=\bigsqcup_i(\overline f)^{-1}(V_{i,\alpha}).
\]
Ya que $\overline f$ es continua, se tiene que $(\overline f)^{-1}(V_{i,\alpha})$
son abiertos de $X\times_BX'$. Veamos que
$\overline p:(\overline f)^{-1}(V_{i,\alpha})\cong f^{-1}(U_\alpha)$. La
restricción $\overline p:(\overline f)^{-1}(V_{i,\alpha})\to f^{-1}(U_\alpha)$
es el pullback de $p:V_{i,\alpha}\to U_\alpha$ a lo largo de
$f:f^{-1}(U_\alpha)\to U_\alpha$. Por lo tanto $\overline p$ es
homeomorfismo.
\end{enumerate}

\item (\textbf{Acciones de grupo y cubrientes}) Sea $G$ un grupo actuando
sobre un espacio $X$. Supongamos que la acción de $G$ en $X$ es propiamente
discontinua\footnote{La acción de $G$ sobre $X$ es propiamente discontinua
si para cada $x\in X$ existe una vecindad $U=U_x$ de $x$ tal que
$gU\cap g'U=\varnothing$ para toda $g,g'\in G$ tales que $g\ne g'$.}.
Defínase la aplicación canónica $p:X\to X/G$.
\begin{enumerate}[label=\alph*)]
\item Probar que $p$ es una aplicación abierta.
\item Deducir que $p:X\to X/G$ es un cubriente.
\item Defínase $h:G\to\Aut(p)$ como $h(g)(x)=gx$. Probar que si $X$ es
conexo, entonces $h$ es un isomorfismo.
\end{enumerate}

\noindent\textbf{Solución.}
\begin{enumerate}[label=\alph*)]
\item Sea $U\subset X$ abierto. Entonces
\[
p^{-1}(p(U))=\{x\in X\mid p(x)\in p(U)\}
=\{x\in X\mid Gx=Gy \text{ para algún } y\in U\}
=\{x\in X\mid x=gy \text{ para algún } y\in U,\ g\in G\}
=\bigcup_{g\in G}gU.
\]
La acción de cada $g\in G$ sobre $X$ define un homeomorfismo
$\varphi_g:X\to X$, $x\mapsto gx$. Así, $\varphi_g(U)=gU$; es decir, $gU$ es
abierto para cada $g\in G$. Se sigue que $p^{-1}(p(U))$ es abierto y por lo
tanto $p(U)$ es abierto en $X/G$.
\item Sea $x\in X$ y $U=U_x$ una vecindad de $x$ que satisface la condición
de la discontinuidad propia. Entonces, por el inciso anterior $p(U)$ es
abierta y es una vecindad abierta de $p(x)=Gx$. Además,
$p^{-1}(p(U))=\bigcup_{g\in G}gU$ en donde $\{gU\mid g\in G\}$ son
subconjuntos abiertos disjuntos de $X$ gracias a que la acción es
propiamente discontinua. Finalmente, $p|_{gU}:gU\to p(U)$ es continua,
biyectiva (si $x,x'\in gU$ verifican que $p(x)=p(x')$, entonces existe
$g'\in G$ tal que $x=g'x'$; el elemento $x$ pertenece entonces a
$gU\cap(g'gU)$. Por lo tanto $g'=e$ y $x=x'$) y abierta; y por lo tanto un
homeomorfismo.
\item Veamos que $h:G\to\Aut(p)$ es un homomorfismo inyectivo. En efecto,
supongamos que $g\in\ker h$. Entonces $gx=x$. Si $g\ne e$, entonces podemos
encontrar una vecindad abierta de $x$, $U=U_x$, tal que $gU\cap U\ne\varnothing$,
lo cual es imposible ya que $x\in gU\cap U$.

Luego, veamos que $h:G\to\Aut(p)$ es sobreyectiva. Por demostrar que, para
toda $\phi\in\Aut(p)$, existe $g\in G$ tal que $h(g)=\phi$. Para observar
esto, dado que por hipótesis $X$ es conexo, usaremos la
Proposición~\ref{prop:unicidad-levantamiento}. Fijemos cualquier punto
$x_0\in X$. Dado $\phi\in\Aut(p)$, ya que $x_0$ y $\phi(x_0)$ se encuentran
en la misma órbita, existe $g\in G$ tal que $gx_0=\phi(x_0)$. Por lo tanto,
ya que $\phi$ y $h(g)$ coinciden en un punto de $X$, siendo éste conexo, la
Proposición~\ref{prop:unicidad-levantamiento} nos da lo requerido.
\end{enumerate}

\item (\textbf{Grupo fundamental de $(S^1)^n$ vía cubrientes}) Probar lo
siguiente:
\begin{enumerate}[label=\alph*)]
\item $p:\RR\to\RR/\ZZ$ es un cubriente.
\item $\pi_1(\RR/\ZZ)\cong\ZZ$.
\item $\pi_1(\RR^n/\ZZ^n)\cong\ZZ^n$.
\end{enumerate}

\noindent\textbf{Solución.}
\begin{enumerate}[label=\alph*)]
\item El grupo aditivo $\ZZ$ actúa sobre $\RR$ por medio de la acción
$\ZZ\times\RR\to\RR$, $(n,x)\mapsto x+n$. Notemos que la acción es
propiamente discontinua ya que, si $x\in\RR$ y $0<\epsilon<1/2$, entonces
$U=U_x=(x-\epsilon,x+\epsilon)$ es una vecindad abierta de $x$ que cumple la
condición de la discontinuidad propia. Por el ejercicio anterior, se tiene
el resultado.
\item Ya que $\RR$ es contraíble, es simplemente conexo. El
Teorema~\ref{thm:pi1-G-cubriente} prueba lo requerido.
\item Este inciso puede ser resuelto de las siguientes maneras.
\begin{enumerate}[label=\arabic*)]
\item Es posible definir una acción propiamente discontinua de $\ZZ^n$
sobre $\RR^n$ de manera análoga como en el inciso anterior y, usando que
$\RR^n$ es simplemente conexo, el Teorema~\ref{thm:pi1-G-cubriente} termina
el resultado.
\item Por el primer inciso, sea $p_i:\RR\to\RR/\ZZ$ un cubriente para
$i=1,\dots,n$. Por el ejercicio 2, se tiene que
\[
p:=\prod_{i=1}^np_i:\RR^n=\RR\times\cdots\times\RR\to\RR/\ZZ\times\cdots\times\RR/\ZZ=\RR^n/\ZZ^n
\]
es aplicación cubriente. Usando el Teorema~\ref{thm:pi1-G-cubriente}
concluimos.
\item Usando el ejercicio 1 inciso d),
\[
\pi_1(\RR^n/\ZZ^n)=\pi_1\left(\prod_{i=1}^n\RR/\ZZ\right)=\prod_{i=1}^n\pi_1(\RR/\ZZ)=\prod_{i=1}^n\ZZ=\ZZ^n.
\]
\end{enumerate}
\end{enumerate}

\item (\textbf{Grupo fundamental de $\RR P^n$ y de $L_m(\ell_1,\dots,\ell_n)$})
Probar lo siguiente:
\begin{enumerate}[label=\alph*)]
\item Si $G$ es un grupo finito actuando de manera libre sobre un espacio
Hausdorff $X$, entonces la acción es propiamente discontinua.
\item $\pi_1(\RR P^n)\cong\ZZ/2\ZZ$, si $n\ge2$. Cuando $n=1$,
$\pi_1(\RR P^1)\cong\ZZ$.
\item $\pi_1(L_m(\ell_1,\dots,\ell_n))\cong\ZZ/m\ZZ$.
\end{enumerate}

\noindent\textbf{Solución.}
\begin{enumerate}[label=\alph*)]
\item Sea $G=\{1=g_0,g_1,\dots,g_n\}$. Ya que $X$ es Hausdorff, existen
vecindades $U_0,U_1,\dots,U_n$ de $g_0x,g_1x,\dots,g_nx$ respectivamente con
$U_0\cap U_j=\varnothing$ para $j=1,\dots,n$. Sea $U=\bigcap_{j=0}^ng_j^{-1}U_j$
que es una vecindad abierta de $x$. Entonces,
\[
g_iU\cap g_jU=g_j((g_j^{-1}g_iU)\cap U)=g_j(g_kU\cap U)=\varnothing,
\]
para alguna $k$.
\item Sea $n\ge2$. Consideremos la siguiente acción de $\{\pm1\}$ sobre
$S^n$: $\{\pm1\}\times S^n\to S^n$, $(\pm1,x)\mapsto\pm x$. Veamos que la
acción es libre. Supongamos que $a\in\Stab_{\{\pm1\}}(x)$. Entonces, $ax=x$.
Por lo tanto $a=1$. De esta forma, por el
Teorema~\ref{thm:accion-propia-cubriente}, $S^n/\{\pm1\}:=\RR P^n$ es un
espacio topológico separado, y la aplicación canónica $S^n\to\RR P^n$ es un
cubriente con $2$ hojas (por el Corolario~\ref{cor:accion-finita-cubriente}).
Finalmente, ya que $S^n$ es simplemente conexo, por el
Teorema~\ref{thm:pi1-G-cubriente},
\[
\pi_1(\RR P^n)\cong\ZZ/2\ZZ.
\]
\item Recordemos la construcción del espacio lenticular. Sea $m\in\NN$, y
sean $\ell_1,\dots,\ell_n$ números enteros primos relativos a $m$. El grupo
$U_m$ de $m$-ésimas raíces de la unidad es
$U_m=\{e^{2\pi ik/m}\}_{k=0}^{m-1}$. Consideremos la siguiente acción de
$U_m$ sobre $S^{2n-1}=\{(z_1,\dots,z_n)\in\CC^n\mid|z_1|^2+\cdots+|z_n|^2=1\}$:
\[
U_m\times S^{2n-1}\to S^{2n-1}, \qquad (\lambda,(z_1,\dots,z_n))\mapsto(\lambda^{\ell_1}z_1,\dots,\lambda^{\ell_n}z_n).
\]
Veamos que la acción es libre. Sea $z\in S^{2n-1}$ y supongamos que
$\lambda\in U_m$ pertenece a $\Stab_{U_m}(z)$. Entonces,
$(\lambda^{\ell_1}z_1,\dots,\lambda^{\ell_n}z_n)=(z_1,\dots,z_n)$. Entonces,
$\lambda^{\ell_j}z_j=e^{2\pi ik\ell_j/m}z_j=z_j$, para $j\in\{1,\dots,n\}$
tal que $z_j\ne0$. De la ecuación, se obtiene que $e^{2\pi ik\ell_j/m}=1$.
Por tanto, existe un entero $r\in\ZZ$ tal que $r=k\ell_j/m$. Así, $m\mid k\ell_j$.
Ya que $(\ell_j,m)=(m,\ell_j)=1$, se sigue que $m\mid k$. Esto último solo
puede ocurrir cuando $k=0$. Entonces, debe cumplirse que $\lambda=1$ y, de
esta forma, la acción es libre.

De esta manera, obtenemos un espacio topológico $L_m(\ell_1,\dots,\ell_n):=S^{2n-1}/U_m$
separado. Luego, por el Teorema~\ref{thm:accion-propia-cubriente}, se tiene
el cubriente de $m$ hojas (Corolario~\ref{cor:accion-finita-cubriente}):
$S^{2n-1}\to L_m(\ell_1,\dots,\ell_n)$. Ya que $S^{2n-1}$ es simplemente
conexo, por el Teorema~\ref{thm:pi1-G-cubriente},
\[
\pi_1(L_m(\ell_1,\dots,\ell_n))\cong U_m\cong\ZZ/m\ZZ.
\]
\end{enumerate}

\item (\textbf{Simplemente conexo y levantamientos}) Sea $p:X\to B$ un
cubriente, de espacio total $X$ conexo por arcos, $b\in B$ y
$x\in p^{-1}(b)$. Entonces, $X$ es simplemente conexo si, y solamente si,
dos lazos $\alpha,\alpha'\in\Omega(B,b)$ son homótopos si sus
levantamientos de origen $x$ tienen los mismos extremos.

\noindent\textbf{Solución.} Si $X$ es simplemente conexo, dos caminos en
$X$ que tienen el mismo origen y el mismo extremo son homótopos, y por lo
tanto sus imágenes bajo $p$ también.

Recíprocamente, si $\alpha$ y $\alpha'$ son dos caminos en $X$ con el mismo
origen $x$ y mismo extremo, entonces los caminos $p\circ\alpha$ y
$p\circ\alpha'$ en $B$ tienen los mismos levantamientos en $X$ con mismo
origen y mismo extremo, y por tanto homótopos, y por levantamiento de
homotopías, $\alpha$ y $\alpha'$ son homótopos.

\item (\textbf{Acción sobre la fibra}) Sea $p:X\to B$ una aplicación
cubriente y $x\in p^{-1}(b)$.
\begin{enumerate}[label=\alph*)]
\item La aplicación
\[
p^{-1}(b)\times\pi_1(B,b)\to p^{-1}(b), \qquad (x,[\alpha])\mapsto x[\alpha]:=\tilde\alpha(1)
\]
define una acción de $\pi_1(B,b)$ sobre $p^{-1}(b)$.
\item $\Stab_{\pi_1(B,b)}(x)=p_*\pi_1(X,x)$.
\end{enumerate}

\noindent\textbf{Solución.}
\begin{enumerate}[label=\alph*)]
\item Sean $\alpha:I\to B$ y $\alpha':I\to B$ dos lazos de base $b$. El
levantamiento $\widetilde{\alpha\alpha'}\in\Lev(\alpha\alpha')$ tal que
$\widetilde{\alpha\alpha'}(0)=x$ es el producto del levantamiento
$\tilde\alpha\in\Lev(\alpha)$ tal que $\tilde\alpha(0)=x$ y del
levantamiento $\tilde\alpha'\in\Lev(\alpha')$ tal que
$\tilde\alpha'(0)=\tilde\alpha(1)$. Por consiguiente,
\[
x([\alpha][\alpha'])=(x[\alpha])[\alpha'].
\]
Además, el levantamiento de origen $x$ del lazo constante $c_b$ es el lazo
constante $c_x$. Así, $xe=x$ para todo $x\in p^{-1}(b)$ (aquí, $e$ es el
elemento neutro de $\pi_1(B,b)$).
\item Sea $g=[\alpha]\in\pi_1(B,b)$ un elemento tal que $xg=x$. Sea
$\tilde\alpha:I\to X\in\Lev(\alpha)$ con $\tilde\alpha(0)=x$. Ya que $xg=x$,
entonces $\tilde\alpha(1)=x$. Por lo tanto,
$\tilde\alpha\in\Omega(X,x)$ y $g=[\alpha]\in p_*\pi_1(X,x)$.
Recíprocamente, sea $g=[\beta]\in\pi_1(X,x)$. Entonces, por unicidad,
$\beta\in\Lev(p\circ\beta)$ con $\beta(0)=x$ y $\beta(1)=x$. Por lo tanto,
$xp_*(g)=x$.
\end{enumerate}

\item (\textbf{Fibra de un cubriente}) Sea $p:X\to B$ un cubriente.
\begin{enumerate}[label=\alph*)]
\item Supongamos que $B$ es conexo por arcos. La aplicación
$p^{-1}(b)\to\pi_0(X)$ que a todo punto $x\in p^{-1}(b)$ asocia la
componente conexa que contiene a $x$ induce una biyección
\[
p^{-1}(b)/\pi_1(B,b)\cong\pi_0(X)
\]
entre $\pi_0(X)$ y el conjunto de órbitas de $\pi_1(B,b)$ en $p^{-1}(b)$.
\item Supongamos que $B$ es conexo por arcos. El espacio $X$ es conexo por
arcos si y solo si la acción de $\pi_1(B,b)$ sobre $p^{-1}(b)$ es
transitiva.
\item Supongamos que $X$ es conexo por arcos. Entonces, para todo
$x\in p^{-1}(b)$, la aplicación de $\pi_1(B,b)$ en $p^{-1}(b)$ definida por
$[\alpha]\mapsto x[\alpha]$ induce una biyección
\[
\pi_1(B,b)/p_*\pi_1(X,x)\cong p^{-1}(b).
\]
\item Supongamos que $X$ es conexo por arcos. Entonces, $p$ es un cubriente
con $n$ hojas si y solo si, para todo punto $x\in p^{-1}(b)$, el subgrupo
$p_*\pi_1(X,x)$ es de índice $n$ en $\pi_1(B,b)$.
\item Si $X$ es conexo por arcos, entonces $p:X\to B$ es un homeomorfismo
si, y solamente si, el morfismo $p_*:\pi_1(X,x)\to\pi_1(B,b)$ es
isomorfismo de grupos.
\item Un cubriente conexo por arcos de un espacio simplemente conexo es un
homeomorfismo.
\end{enumerate}

\noindent\textbf{Solución.}
\begin{enumerate}[label=\alph*)]
\item La aplicación $p^{-1}(b)\to\pi_0(X)$ es sobreyectiva. En efecto, si
$y\in X$, es posible conectar $p(y)$ y $b$ por un camino $\alpha:I\to B$.
Sea $\tilde\alpha:I\to X\in\Lev(\alpha)$ con $\tilde\alpha(0)=y$ y
$\tilde\alpha(1)\in p^{-1}(b)$. De aquí que, toda componente conexa por
arcos de $X$ contiene al menos un punto de $p^{-1}(b)$.

Ahora, consideremos dos puntos $x$ y $x'$ de $p^{-1}(b)$. Si existe un
camino $\alpha:I\to X$ que conecte a $x$ con $x'$, tenemos que
$x'=x[p\circ\alpha]$. Recíprocamente, si existe un elemento
$[\beta]\in\pi_1(B,b)$ tal que $x'=x[\beta]$, el levantamiento de origen $x$
de un representante cualquiera $[\alpha]$ forma un camino entre $x$ y $x'$.
\item Supongamos que $X$ es conexo por arcos. Queremos ver que, para todo
$x$ y $x'$ en $p^{-1}(b)$, existe un lazo $\beta\in\Omega(B,b)$ tal que
$x[\beta]=x'$. Pero ya que $X$ es conexo por arcos, siempre existe
$\alpha:I\to X$ con $\alpha(0)=x$ y $\alpha(1)=x'$. Tomando $\beta=p\circ\alpha$
se sigue que la acción es transitiva.

Recíprocamente, si la acción es transitiva, el último punto de la prueba en
el inciso anterior concluye el resultado.
\item Sea $x\in p^{-1}(b)$. Verificamos que la aplicación
$\varphi_x:\pi_1(B,b)\to p^{-1}(b)$, $[\alpha]\mapsto x[\alpha]$, induce una
aplicación $\overline\varphi_x:\pi_1(B,b)/\Stab_{\pi_1(B,b)}(x)\to p^{-1}(b)$.
En efecto, dos elementos $[\alpha],[\alpha']$ de $\pi_1(B,b)$ son
equivalentes si, y solamente si, $[\alpha]^{-1}[\alpha']\in\Stab_{\pi_1(B,b)}(x)$,
que equivale a $x[\alpha]=x[\alpha']$. La aplicación $\overline\varphi_x$ es
inyectiva por construcción. Ya que $X$ es conexo por arcos, se tiene que la
acción es transitiva y, por lo tanto, $\overline\varphi_x$ es sobreyectiva.
Finalmente, ya que $\Stab_{\pi_1(B,b)}(x)=p_*\pi_1(X,x)$, se sigue lo que se
quería probar.
\item Ya que $X$ es conexo por arcos, gracias al ejercicio anterior se
tiene que
\[
n=\#p^{-1}(b)=[\pi_1(B,b):p_*\pi_1(X,x)].
\]
\item Si $p_*$ es un isomorfismo, entonces $p_*\pi_1(X,x)=\pi_1(B,b)$.
Entonces, $\#p^{-1}(b)=1$. Por lo tanto, $p$ es inyectiva y así, un
homeomorfismo (ya que todo cubriente es continuo, abierto y sobre).
\item Si $p:X\to B$ es tal cubriente, entonces $p_*$ es inyectiva, y de
imagen nula, de donde un isomorfismo. Por el inciso anterior, $p$ es un
homeomorfismo.
\end{enumerate}

\item (\textbf{Cubrientes y conjugación}) Sea $p:X\to B$ un cubriente.
Sean $x_0$ y $x_1$ puntos de $X$ que pertenecen a una componente conexa por
arcos de $X$ tales que $p(x_0)=p(x_1)$. Sea $b=p(x_0)$. Entonces, los
subgrupos $p_*\pi_1(X,x_0)$ y $p_*\pi_1(X,x_1)$ son conjugados en
$\pi_1(B,b)$: existe $[\alpha]\in\pi_1(B,b)$ tal que
\[
p_*\pi_1(X,x_1)=[\alpha]\,p_*\pi_1(X,x_0)\,[\alpha]^{-1}.
\]
Recíprocamente, para todo elemento $[\alpha']\in\pi_1(B,b)$, existe un
punto $x_1'\in p^{-1}(b)$ tal que
\[
p_*\pi_1(X,x_1')=[\alpha']\,p_*\pi_1(X,x_0)\,[\alpha']^{-1}.
\]

\noindent\textbf{Solución.} Sea $\alpha\in\Omega(X,x_0,x_1)$. Sea
$\alpha_*:\pi_1(X,x_0)\to\pi_1(X,x_1)$ el isomorfismo de cambio de punto de
base a lo largo de $\alpha$.

Sea $\beta=p\circ\alpha\in\Omega(B,b)$. El elemento $[\beta]\in\pi_1(B,b)$
define un automorfismo $\beta_*:\pi_1(B,b)\to\pi_1(B,b)$,
$[\gamma]\mapsto[\beta\gamma\overline\beta]$. Se tiene el diagrama
conmutativo siguiente:
\[
\begin{tikzcd}
\pi_1(X,x_0) \arrow[r,"\alpha_*"] \arrow[d,swap,hook,"p_*"] & \pi_1(X,x_1) \arrow[d,hook,"p_*"] \\
\pi_1(B,b) \arrow[r,swap,"\beta_*"] & \pi_1(B,b)
\end{tikzcd}
\]
En efecto, sea $\delta:I\to X\in\Omega(X,x_0)$. Entonces,
\[
(p_*\circ\alpha_*)([\delta])=p_*([\alpha\delta\overline\alpha])
=[p\circ\alpha][p\circ\delta][p\circ\overline\alpha]
=[\beta]p_*([\delta])[\beta]^{-1}
=(\beta_*\circ p_*)([\delta]).
\]
Ya que $\alpha_*$ es un isomorfismo,
$\beta_*(p_*\pi_1(X,x_0))=p_*\pi_1(X,x_1)$.

Recíprocamente, escojamos un elemento $[\alpha']\in\pi_1(B,b)$. Sea
$\overline\alpha':I\to B\in\Omega(B,b)$ un lazo que representa la clase
$\alpha'$ y sea $\beta':I\to X\in\Lev(\overline\alpha')$ tal que
$\beta'(0)=x_0$. Pongamos $x_1'=\beta'(1)$. Los argumentos precedentes
muestran que
\[
p_*\pi_1(X,x_1')=[\alpha']\,p_*\pi_1(X,x_0)\,[\alpha']^{-1}.
\]

\item (\textbf{Cubriente de Galois}) Sea $p:X\to B$ una aplicación cubriente
con espacio total $X$ conexo y localmente conexo por arcos. Sea $x\in X$ y
$b=p(x)$.
\begin{enumerate}[label=\alph*)]
\item La acción de $\Aut(p)$ sobre $p^{-1}(b)$ es transitiva si, y
solamente si, $p_*\pi_1(X,x)$ es normal en $\pi_1(B,b)$; esto es, $p$ es
regular.
\item Sea $p$ un cubriente regular, entonces
$\Aut(p)\cong\pi_1(B,b)/p_*\pi_1(X,x)$.
\end{enumerate}

\noindent\textbf{Solución.}
\begin{enumerate}[label=\alph*)]
\item Supongamos que la acción de $\Aut(p)$ sobre $p^{-1}(b)$ es
transitiva. Consideremos $\alpha\in\Omega(X,x)$. Veamos que, para todo
$\sigma\in\Omega(B,b)$,
\[
[\sigma]^{-1}[p\circ\alpha][\sigma]=[p\circ\sigma']
\]
para cierto $\sigma'\in\Omega(X,y)$. Sea $\tilde\sigma\in\Lev(\sigma)$ con
$\tilde\sigma(0)=x$ y $\tilde\sigma(1)=y$. Entonces, por hipótesis, $y$ es
la imagen de $x$ por un elemento $h\in\Aut(p)$. Tenemos,
\[
[\sigma]^{-1}[p\circ\alpha][\sigma]=[\overline\sigma(p\circ\alpha)\sigma],
\]
y $\overline\sigma(p\circ\alpha)\sigma$ se levanta en
$\overline{\tilde\sigma}\alpha\tilde\sigma$, un lazo basado en $y=h(x)$.
Así,
\[
[\overline\sigma(p\circ\alpha)\sigma]\in p_*\pi_1(X,y)=p_*\pi_1(X,h(x))=p_*h_*\pi_1(X,x)=(p\circ h)_*\pi_1(X,x)=p_*\pi_1(X,x).
\]
Recíprocamente, supongamos que $p_*\pi_1(X,x)$ es normal en $\pi_1(B,b)$.
Comencemos por construir un homomorfismo de grupos
$h:\pi_1(B,b)\to\Aut(p)$.

Sea $\alpha\in\Omega(B,b)$ y sea $\tilde\alpha\in\Lev(\alpha)$ con
$\tilde\alpha(0)=x$ y $\tilde\alpha(1)=y$. Por hipótesis,
$p_*\pi_1(X,x)=p_*\pi_1(X,y)$. Entonces, por el
Teorema~\ref{thm:teo-levantamiento}, existe un automorfismo
$h_\alpha=h([\alpha])\in\Aut(p)$ que envía $x$ sobre $y$.
\item Consideremos nuevamente $\alpha\in\Omega(B,b)$ y
$\tilde\alpha_x\in\Lev(\alpha)$ a su único levantamiento de origen
$x\in p^{-1}(b)$. Entonces, por definición de $h$,
$h([\alpha])(x)=\tilde\alpha_x(1)$. Además, para todo automorfismo $g$ de
cubrientes, se tiene $g\circ\tilde\alpha_x=\tilde\alpha_{g(x)}$ ya que
$g\circ\tilde\alpha_x\in\Lev(\alpha)$ y
$(g\circ\tilde\alpha_x)(1)=g(x)$.

Mostremos que $h:\pi_1(B,b)\to\Aut(p)$ es un homomorfismo de grupos.
Levantamos a $\alpha$ de $x$ con $h([\alpha])(x)=y$, luego $\alpha'$ de $y$
a $h([\alpha'])(y)=z$; en otros términos,
$\tilde\alpha_x(1)=y$, $\tilde\alpha'_y(1)=z$, $\widetilde{\alpha\alpha'}_x(1)=z$.
El automorfismo $h([\alpha\alpha'])$ envía todo punto $x\in p^{-1}(b)$ en la
extremidad $z$ de $\widetilde{\alpha\alpha'}_x$. Se tiene,
\[
h([\alpha])\circ h([\alpha'])(x)=h([\alpha])(\tilde\alpha'_x(1))=\tilde\alpha'_{h([\alpha])(x)}(1)=\tilde\alpha'_y(1)=z.
\]
Por lo tanto, $h:\pi_1(B,b)\to\Aut(p)$ es un homomorfismo de grupos. Ahora
bien,
\[
\ker(h)=\{[\alpha]\in\pi_1(B,b)\mid h([\alpha])(x)=x\ \forall x\in p^{-1}(b)\}=p_*\pi_1(X,x).
\]
\end{enumerate}

\item (\textbf{Ejemplos de cubrientes de Galois})
\begin{enumerate}[label=\alph*)]
\item Determinar todos los automorfismos del cubriente $S^1\to S^1$,
$z\mapsto z^n$ ($n\ge1$). ¿Es un cubriente de Galois?
\item Calcular $\Aut(S^n\to\RR P^n)$. ¿Es de Galois?
\end{enumerate}

\noindent\textbf{Solución.}
\begin{enumerate}[label=\alph*)]
\item El cubriente $z\mapsto z^n$ es de Galois debido a que el grupo
fundamental de $S^1$ es abeliano, $\ZZ$. Dicho esto, notemos que se tienen
$n$ automorfismos distintos de $z\mapsto z^n$ dados por
$\phi_k(z)=z\exp\left(\frac{2i\pi k}{n}\right)$ donde $k\in\{0,1,\dots,n-1\}$.
Como tiene $n$ hojas y son morfismos que actúan transitivamente sobre las
fibras, obtenemos que el cubriente es de Galois y que no pueden existir
otros morfismos de cubrientes. Además, $\Aut_{S^1}(S^1\to S^1)\cong\ZZ/n\ZZ$.
\item Tenemos que $\Aut(S^n\to\RR P^n)\cong\ZZ/2\ZZ$.
\end{enumerate}

\item (\textbf{Cubrientes que no son Galois}) ¿Es posible obtener un
cubriente de Galois de $\RR P^2\vee\RR P^2$ en donde el grupo de
automorfismos sea $\ZZ/4\ZZ$?

\noindent\textbf{Solución.} La respuesta es no, ya que si se tuviera tal
cubriente $p:X\to\RR P^2\vee\RR P^2$, entonces
\[
\ZZ/4\ZZ=\Aut(p)\cong\pi_1(\RR P^2\vee\RR P^2)/p_*\pi_1(X)\cong(\ZZ/2\ZZ*\ZZ/2\ZZ)/p_*\pi_1(X),
\]
de donde cualquier cociente conmutativo de tal grupo no contiene elementos
de orden $4$.

\item (\textbf{Aplicación sobre la esfera nulhomotópica}) Demuestra que si
$Y$ es conexo por arcos, localmente conexo por arcos y tiene grupo
fundamental finito, entonces cada aplicación $Y\to S^1$ es nulhomotópica.

\noindent\textbf{Solución.} Usaremos el
Teorema~\ref{thm:teo-levantamiento}. Sea $f:Y\to S^1$ y considérese el
cubriente $p:\RR\to S^1$. Ya que $\pi_1(Y)$ es finito, entonces
$f_*\pi_1(Y)$ es finito. Por otra parte, $p_*\pi_1(\RR)=\ZZ$. Así,
$f_*\pi_1(Y)\subset\ZZ$ si, y sólo si, $f_*\pi_1(Y)=\{1\}$. Por el
Teorema~\ref{thm:teo-levantamiento}, existe un levantamiento
$\tilde f:Y\to\RR$. Ya que $\RR$ es contraíble, existe una homotopía
$F:\RR\times[0,1]\to\RR$, $(x,t)\mapsto(1-t)x$. Sea
$F':Y\times[0,1]\to\RR$ la aplicación continua dada por
$F'(y,t)=(1-t)\tilde f(y)$. Entonces, $F':\tilde f\simeq c_{x=\tilde f(y)}$.
Finalmente, tomando $H=p\circ F'$: $f\simeq c_x$.

\item (\textbf{Espacio conexo y localmente conexo que no tiene cubriente
universal}) Sea $B=\prod_{i=1}^\infty S^1$. Probar que $B$ no tiene
cubriente universal.

\noindent\textbf{Solución.} Consideremos la proyección $r_n:B\to S^1$ dada
por $(z_1,z_2,z_3,\dots)\mapsto z_1$. El espacio $B$ satisface la propiedad
de que, para cualquier espacio $Z$, $f:Z\to B$ es continua si, y solamente
si, $r_n\circ f:Z\to S^1$ es continua. Entonces un lazo $\phi:I\to B$ es
equivalente a lazos $(\phi_1,\phi_2,\dots)$ de $S^1$ con la propiedad
$\phi_i=r_i\circ\phi$ para cada $i$. Así, $\phi\simeq\psi\ (\mathrm{rel}\,\partial I)$
si, y solo si, $\phi_i\simeq\psi_i\ (\mathrm{rel}\,\partial I)$ para todo
$i$. Concluimos que
\[
\pi_1(B)\cong\prod_{i=1}^\infty\pi_1(S^1)\cong\prod_{i=1}^\infty\ZZ.
\]

El espacio $B$ es separado, conexo y localmente conexo por arcos.
Supongamos que $B$ admite un cubriente universal. Entonces, por el
Teorema~\ref{thm:existencia-cubriente-universal}, $B$ tiene que ser
semilocalmente simplemente conexo. De esta manera, el punto
$\mathbf 1=(1,1,1,\dots)$ admite una vecindad $U\ni\mathbf 1$ tal que, si
$i:U\to B$ es la inclusión, el morfismo $i_*:\pi_1(U,\mathbf 1)\to\pi_1(B,\mathbf 1)$
es el morfismo nulo.

Ahora bien, para todo $U$, existe $n\in\NN$ tal que $U$ contiene al punto
$(1,1,\dots,1,z_{n+1},z_{n+2},\dots)$ para $z_{n+1},z_{n+2},\dots\in S^1$.
Consideremos el siguiente lazo
\[
\phi:I\to U, \qquad t\mapsto(1,\dots,1,e^{2\pi it},1,1,\dots).
\]
Entonces $[\phi]\ne1$ en $\pi_1(B)$; en contradicción con la hipótesis
sobre $U$.

\item (\textbf{Clasificación de cubrientes}) Describir, hasta isomorfismos
de espacios cubrientes, todos los espacios cubrientes conexos de los
siguientes espacios:
\begin{enumerate}[label=\alph*)]
\item $S^1$ como el círculo unitario en el plano complejo.
\item $\RR P^n$.
\end{enumerate}

\noindent\textbf{Solución.} Aplicaremos el
Teorema~\ref{thm:clasificacion-cubrientes}.
\begin{enumerate}[label=\alph*)]
\item Se tiene que $\pi_1(S^1)=\ZZ$. Los únicos subgrupos de $\ZZ$ no
triviales son $n\ZZ$. Por lo tanto, los cubrientes conexos de $S^1$ son de
la forma $\RR/n\ZZ\to S^1:t\ \mathrm{m\acute od}\ n\mapsto e^{2\pi it}$. La
multiplicación por $n$ induce un homeomorfismo entre $\RR/n\ZZ$ y
$\RR/\ZZ\cong S^1$. Así, hasta isomorfismos, los cubrientes conexos de
$S^1$ son $S^1\to S^1:z\mapsto z^n$, de fibra $\ZZ/n\ZZ$, y su cubriente
universal $\exp:\RR\to S^1:t\mapsto e^{2\pi it}$, que corresponde al
subgrupo trivial de $\ZZ$. El estabilizador es $n\ZZ$ y el grupo de
automorfismos es el cociente $\ZZ/n\ZZ$; sus elementos son
$z\mapsto ze^{2i\pi k/n}$.
\item Sabemos que $\pi_1(\RR P^n)=\ZZ/2\ZZ$ para $n\ge2$. Como los únicos
subgrupos de $\ZZ/2\ZZ$ son $\ZZ/2\ZZ$ y $\{0\}$, para $n\ge2$, los
cubrientes conexos del espacio proyectivo real $\RR P^n$ son, hasta
isomorfismos, su cubriente universal $S^n\to\RR P^n$ y la identidad
$\RR P^n\to\RR P^n$.
\end{enumerate}

\item (\textbf{Grupos de homotopía superiores y cubrientes})
\begin{enumerate}[label=\alph*)]
\item Sea $p:(X,x)\to(B,b)$ un cubriente en donde $b=p(x)$. Mostrar que
para $q\ge2$, se tiene un isomorfismo de grupos
$p_*:\pi_q(X,x)\to\pi_q(B,b)$.
\item Calcular $\pi_q(S^1,*)$ para $q\ge2$.
\item Probar que
\[
\pi_q(\mathbb T^2)=\begin{cases}\ZZ\oplus\ZZ & q=1\\ 0 & q>1.\end{cases}
\]
\item Mostrar que $\pi_q(S^n)\cong\pi_q(\RR P^n)$ para $q>1$.
\item Mostrar que $\pi_q(S^n)=0$ si $q<n$.
\item Supongamos que $q\le n$. Calcular $\pi_q(\RR P^n)$.
\item Si $q<n$, ¿existe una retracción de $\RR P^n$ a $\RR P^q$?
\end{enumerate}

\noindent\textbf{Solución.}
\begin{enumerate}[label=\alph*)]
\item Veamos que $p_*$ es sobreyectiva. Esto es, que para toda
$[f]\in\pi_q(B,b)$, existe $[g]\in\pi_q(X,x)$ tal que $p_*([g])=[f]$. Sea
$f:(S^q,s)\to(B,b)$ una aplicación continua. Ya que $q\ge2$,
$\pi_1(S^q,s)=0$. Entonces, se satisface que
$f_*\pi_1(S^q,s)=0\subset p_*\pi_1(X,x)$.

Así, por el Teorema~\ref{thm:teo-levantamiento}, existe un levantamiento
$\tilde f:(S^q,s)\to(X,x)\in\Lev(f)$ tal que el diagrama
\[
\begin{tikzcd}
& X \arrow[d,"p"] \\
Y \arrow[ur,dashed,"\tilde f"] \arrow[r,swap,"f"] & B
\end{tikzcd}
\]
conmuta: $p\circ\tilde f=f$. Tomando $g=\tilde f$ se cumple que
$[f]=[p\circ\tilde f]=p_*([\tilde f])$.

Ahora veamos que $p_*$ es inyectivo. Supongamos que $[\tilde f]\in\ker p_*$.
Entonces, $p_*([\tilde f])=[p\circ\tilde f]=0$. Por lo tanto, existe una
homotopía $H_t:f=p\circ\tilde f\simeq c_b$ vía
$H:(S^q,s)\times I\to(B,b)$ tal que $H_1=f$ y $H_0=c_b$. Por el
Teorema~\ref{thm:levantamiento-caminos}, existe un único levantamiento
$\tilde H_t:(S^q,s)\to(X,x)$ tal que $p\circ\tilde H_t=H_t$. Así,
$p\circ\tilde H_1=H_1=f$ y $p\circ\tilde H_0=H_0=c_b$. Por unicidad,
$\tilde H_1=\tilde f$ y $\tilde H_0=c_x$. Esto implica que
$\tilde H_t:\tilde f\simeq c_x$ y $[\tilde f]=0$.
\item Ya que $\RR\to S^1$ es cubriente y $\pi_q(\RR,*)=0$ se tiene, por el
inciso anterior, que $\pi_q(S^1,*)=\pi_q(\RR,*)=0$ para $q\ge2$.
\item Se tiene que $\pi_1(S^1,*)\cong\ZZ$. Entonces,
\[
\pi_1(\mathbb T^2)=\pi_1(S^1\times S^1)\cong\pi_1(S^1)\oplus\pi_1(S^1)=\ZZ\oplus\ZZ.
\]
Por otra parte, se tiene que $\RR^2\to S^1\times S^1$ es cubriente.
Entonces, ya que $\RR^n$ es contraíble entonces, para $q\ge2$,
$\pi_q(\mathbb T^2)\cong\pi_q(\RR^2)=0$.

También, es posible usar el ejercicio 4 como sigue:
\[
\pi_q(\mathbb T^2)=\pi_q(S^1)\oplus\pi_q(S^1)=\begin{cases}\ZZ\oplus\ZZ & q=1\\0 & q\ge0.\end{cases}
\]
\item Ya que $S^n\to\RR P^n$ es cubriente, entonces
$\pi_q(S^n)\cong\pi_q(\RR P^n)$ para $q\ge2$.
\item Sea $[g]\in\pi_q(S^n)$. Entonces, existe\footnote{Sean $M$ y $N$
variedades suaves y $g:M\to N$ una función continua, entonces existe
$f:M\to N$ suave tal que $g\simeq f$.} una función suave $f:S^q\to S^n$
tal que $[f]=[g]$. Por el Teorema de Brown--Sard\footnote{\textbf{Brown--Sard}.
Sea $f:M\to N$ una función suave, entonces el conjunto de valores regulares
de $f$ es denso en $N$.}, existe un valor regular $y_0\in S^n$ de $f$. Así,
$f^{-1}(y_0)$ es una subvariedad regular de $S^q$ de dimensión $q-n<0$, lo
que implica que $f^{-1}(y_0)=\varnothing$. Entonces, $f$ se puede
factorizar como
\[
\begin{tikzcd}
& S^n-\{y_0\}\cong\RR^n \arrow[d,hook,"i"] \\
S^q \arrow[ur,"\tilde f"] \arrow[r,swap,"f"] & S^n
\end{tikzcd}
\]
Entonces, $\tilde f\simeq c$ pues $\RR^n$ es contraíble. Por lo tanto,
$f\simeq i\circ c\simeq c$; esto es, $[g]=[f]=0$.
\item Por los puntos anteriores, $\pi_n(\RR P^n)\cong\pi_n(S^n)$ y
$\pi_q(\RR P^n)\cong0$.

Ahora bien, el Teorema de Hurewicz afirma que, para cualquier espacio $X$ y
para cualquier $q\ge1$, existe un homomorfismo de grupos
$h:\pi_q(X)\to H_q(X)$ que, para $q\le n$ con $n\ge2$, es un isomorfismo si
$X$ es simplemente conexo. En particular,
\[
\pi_n(\RR P^n)\cong\pi_n(S^n)\cong H_n(S^n)\cong\ZZ.
\]
\item Sea $q<n$ y supongamos que existe una retracción $r:\RR P^n\to\RR P^q$.
Sea $i:\RR P^q\to\RR P^n$ la inclusión. La aplicación $r\circ i$ induce la
identidad sobre los grupos de homotopía, y por tanto sobre
$\pi_q(\RR P^q)\cong\ZZ$, lo cual es imposible porque se factoriza por
$\{0\}$.
\end{enumerate}

\end{enumerate}

\section{Homología de espacios}
\label{sec:homologia-espacios}

\subsection{Definición de homología}

Sea $q\in\NN$, el \emph{$q$-simplejo estándar} (ordenado) $\Delta_q$ es la
envolvente convexa afín en $\RR^{q+1}$ de vectores en la base canónica
$(e_0,e_1,\dots,e_q)$:
\[
\Delta_q=\left\{(t_0,\dots,t_q)\in\RR^{q+1} \,\middle|\, t_i\ge0,\ \sum_it_i=1\right\}.
\]
Sea $X$ un espacio. Definimos el conjunto de \emph{$q$-simplejos singulares}
de $X$ como
\[
S_q(X)=R\big(\{\sigma:\Delta_q\to X\mid\sigma\text{ es continua}\}\big).
\]
Así, los $0$-simplejos se identifican con los puntos de $X$, y los
$1$-simplejos singulares con los caminos en $X$ (utilizando la única
aplicación afín que manda $[0,1]$ sobre $\Delta_1$ enviando $0\in[0,1]$
sobre $e_0=(1,0)\in\Delta_1$, la cual es $t\mapsto(1-t,t)$). Se define,
\[
\partial_q:S_q(X)\to S_{q-1}(X), \qquad \sigma\mapsto\sum_{i=0}^q(-1)^i\sigma\circ F_q^i,
\]
en donde $F_q^i:\Delta_{q-1}\to\Delta_q$ es la función afín dada por
\[
F_q^i(e_j)=\begin{cases}e_j & j<i\\ e_{j+1} & j\ge i.\end{cases}
\]
Se cumple que $\partial_q\circ\partial_{q+1}=0$.

\begin{definition}[Módulo de homología singular]
\label{def:homologia-singular}
Al complejo de cadenas $\{S_\bullet(X),\partial_\bullet\}$ se le llama
\emph{el complejo de cadenas singulares} de $X$. Se define el $q$-ésimo
módulo de homología singular como
\[
H_q(X):=H_q(X;R)=Z_q(X)/B_q(X),
\]
en donde $Z_q(X)=\ker\partial_q$ son los $q$-ciclos singulares y
$B_q(X)=\im\partial_q$ son las $q$-fronteras singulares.

Sea $f:X\to Y$ una aplicación continua. Definimos
\[
f_\#:S_q(X)\to S_q(Y), \qquad \sigma\mapsto f\circ\sigma.
\]
Entonces, $f_\#$ es morfismo de cadenas y, por lo tanto, define una
aplicación $f_*:H_\bullet(X)\to H_\bullet(Y)$.

Sea $A\subset X$ un subespacio. Definimos,
\[
Z_q(X,A)=\{c\in S_q(X)\mid\partial_q(c)\in S_{q-1}(A)\}
\]
y
\[
B_q(X,A)=\{c\in S_q(X)\mid c=\partial_{q+1}(c')+d,\ c'\in S_{q+1}(X),\ d\in S_q(A)\}.
\]
\end{definition}

\begin{definition}[Módulo de homología singular para una pareja]
\label{def:homologia-pareja}
Sea $(X,A)$ una pareja de espacios. Se define el $q$-ésimo módulo de
homología para $(X,A)$ como
\[
H_q(X,A)=Z_q(X,A)/B_q(X,A).
\]
Nótese que, si $A=\varnothing$, entonces $H_q(X,A)=H_q(X)$ pues
$S_q(\varnothing)=0$ para toda $q$.

Sea $f:(X,A)\to(Y,B)$ una aplicación de parejas. Se tiene que
$f_\#(Z_q(X,A))\subset f_\#(Z_q(Y,B))$.

En efecto, sea $c\in Z_q(X,A)$. Entonces, $\partial_q(c)\in S_{q-1}(A)$.
Sea $c\in B_q(X,A)$. Entonces, $c=\partial_{q+1}(c')+d$. Luego,
\[
f_\#(c)=f_\#\partial_{q+1}(c')+f_\#(d)=\partial_{q+1}f_\#(c')+f_\#(d).
\]
Como $d\in S_q(A)$, $f_\#(d)\in S_q(B)$; es decir $f_\#(c)\in B_q(Y,B)$.

Por lo tanto, $f_\#$ pasa al cociente y define
\[
f_*:H_\bullet(X,A)\to H_\bullet(Y,B), \qquad [z]\mapsto[f_\#(z)].
\]
\end{definition}

\begin{definition}[Homología reducida]
\label{def:homologia-reducida}
Sea $X\ne\varnothing$. El \emph{complejo aumentado} de $X$ es
$S_\bullet(X)$ con un término adicional $S_{-1}(X):=R$ y
$\partial_0:S_0(X)\to R$ dado por $\partial_0\left(\sum_xr_xx\right)=\sum_xr_x$
(el morfismo $\epsilon$ del Ejercicio 1b). Como $\partial_0\circ\partial_1=0$
(consecuencia directa del Ejercicio 1b, donde se muestra
$\mathrm{im}\,\partial_1=\ker\epsilon$), el complejo aumentado es un
complejo de cadenas. Se define
\[
\tilde H_q(X):=H_q\big(S_\bullet(X)\oplus_{S_0(X)}R\big)=\begin{cases}
\ker\partial_q/\mathrm{im}\,\partial_{q+1} & q>0\\
\ker\epsilon/\mathrm{im}\,\partial_1 & q=0.
\end{cases}
\]
Por el Ejercicio 1b, $\epsilon$ es sobreyectiva, así que la sucesión corta
$0\to\ker\epsilon\to S_0(X)\xrightarrow{\epsilon}R\to0$ se escinde (imagen
libre), dando
\[
H_0(X)\cong\tilde H_0(X)\oplus R,
\]
y $\tilde H_q(X)=H_q(X)$ para $q>0$.
\end{definition}

\subsection{Teoremas}

\begin{theorem}[Sucesión larga de homología]
\label{thm:sucesion-larga}
Sea $(X,A)$ una pareja de espacios. Si $i:A\hookrightarrow X$ y
$j:X\hookrightarrow(X,A)$ son inclusiones, entonces existe una sucesión
exacta larga
\[
\cdots\to H_q(A)\xrightarrow{i_*}H_q(X)\xrightarrow{j_*}H_q(X,A)\xrightarrow{\hat\partial_q}H_{q-1}(A)\to\cdots
\]
Además, la sucesión se extiende a la homología reducida
\[
\cdots\to\tilde H_q(A)\to\tilde H_q(X)\to H_q(X,A)\xrightarrow{\hat\partial_q}\tilde H_{q-1}(A)\to\cdots
\]
\end{theorem}

\begin{proof}
Para cada $q$, $S_q(A)\subset S_q(X)$ es un submódulo, y $\partial_q$ lo
preserva; así, $\{S_q(A)\}$ es un subcomplejo de $\{S_q(X)\}$, y se tiene
la sucesión corta exacta de complejos de cadenas
\[
0\to S_\bullet(A)\xrightarrow{i_\#}S_\bullet(X)\xrightarrow{j_\#}S_\bullet(X)/S_\bullet(A)\to0.
\]
Bajo la identificación $H_q(S_\bullet(X)/S_\bullet(A))=H_q(X,A)$ (un
$q$-ciclo de $S_\bullet(X)/S_\bullet(A)$ es la clase de un
$c\in S_q(X)$ con $\partial_q(c)\in S_{q-1}(A)$, es decir $c\in Z_q(X,A)$;
dos representantes $c,c'$ dan el mismo ciclo del cociente si
$c-c'\in S_q(A)+\partial_{q+1}(S_{q+1}(X))=B_q(X,A)$, exactamente la
Definición~\ref{def:homologia-pareja}), esto es una sucesión exacta corta
de complejos de cadenas.

El \emph{lema del zig-zag} de álgebra homológica afirma que toda sucesión
corta exacta de complejos de cadenas $0\to A_\bullet\to B_\bullet\to C_\bullet\to0$
induce una sucesión exacta larga en homología, con conectante
$\hat\partial_q:H_q(C)\to H_{q-1}(A)$ definido por: dado $[c]\in H_q(C)$,
levantar $c$ a $b\in B_q$ (posible pues $B_q\to C_q$ es sobre), notar que
$\partial b\in A_{q-1}$ (pues su imagen en $C_{q-1}$ es $\partial c=0$), y
poner $\hat\partial_q[c]:=[\partial b]\in H_{q-1}(A)$ (bien definido:
independiente del levantamiento y de un representante homólogo de $c$, y
$\partial b$ es un ciclo de $A_{q-1}$ porque $\partial\partial b=0$ en
$B_{q-2}$ y $A_{q-1}\to B_{q-1}$ es inyectiva). La exactitud en cada uno de
los tres puntos ($H_q(A)$, $H_q(B)$, $H_q(C)$) se verifica por una
persecución de diagramas directa a partir de las definiciones. Aplicado a
nuestra sucesión, con $\hat\partial_q[z]=[\partial_q(z)]$ para
$z\in Z_q(X,A)$ (tomando $b=z\in S_q(X)$ como levantamiento), esto da
exactamente la sucesión enunciada.

Para la versión reducida: si $A\ne\varnothing$, los complejos aumentados
de $A$ y $X$ dan la misma sucesión corta exacta
$0\to S_\bullet(A)\to S_\bullet(X)\to S_\bullet(X)/S_\bullet(A)\to0$ en
grados $\ge0$, pero ahora con el término común $S_{-1}(A)=S_{-1}(X)=R$ en
grado $-1$ (que se cancela en el cociente, pues
$S_{-1}(X)/S_{-1}(A)=0$); el mismo lema del zig-zag aplicado al complejo
aumentado da la sucesión exacta larga reducida, que coincide con la no
reducida salvo en $q=0,-1$, en donde el término $H_{-1}(S_\bullet(X)/S_\bullet(A))=0$
recupera $H_q(X,A)$ sin modificar (ya que $S_\bullet(X)/S_\bullet(A)$ no
tiene el término de aumentación).
\end{proof}

\begin{theorem}[Invarianza homotópica]
\label{thm:invarianza-homotopica}
Sean $f,g:(X,A)\to(Y,B)$ aplicaciones continuas tales que $f\simeq g$.
Entonces $f_*=g_*:H_q(X,A)\to H_q(Y,B)$ para toda $q$.

En particular, si $X$ y $Y$ son dos espacios con el mismo tipo de
homotopía, entonces $H_q(X)\cong H_q(Y)$ para toda $q$.

Por el teorema anterior, si $X$ es contraíble,
\[
H_q(X)=\begin{cases}R & q=0\\0 & q\ne0.\end{cases}
\]
\end{theorem}

\begin{proof}
Sea $F:(X\times I,A\times I)\to(Y,B)$ una homotopía de $f$ a $g$. Basta
construir un \emph{operador prisma}
$P:S_q(X)\to S_{q+1}(Y)$, para cada $q$, tal que
\[
\partial P+P\partial=g_\#-f_\#:S_q(X)\to S_q(Y)
\]
(``$P$ es una homotopía de cadenas entre $f_\#$ y $g_\#$''), y tal que
$P(S_q(A))\subset S_{q+1}(B)$: en efecto, dado esto, para
$z\in Z_q(X,A)$ (es decir $\partial z\in S_{q-1}(A)$),
\[
g_\#(z)-f_\#(z)=\partial(P(z))+P(\partial z),
\]
y $P(\partial z)\in S_q(B)$ ya que $\partial z\in S_{q-1}(A)$; así, módulo
$S_q(B)$, $g_\#(z)-f_\#(z)\equiv\partial(P(z))\in B_q(Y,B)$. Por lo tanto
$g_\#(z)$ y $f_\#(z)$ definen la misma clase en $H_q(Y,B)$, es decir
$f_*=g_*$ en $H_q(X,A)$.

\emph{Construcción de $P$.} Piénsese en $\Delta^q\times I$ como un prisma
con vértices inferiores $v_0,\dots,v_q$ (copia de $\Delta^q\times\{0\}$) y
superiores $w_0,\dots,w_q$ (copia de $\Delta^q\times\{1\}$). El prisma se
triangula en los $(q+1)$-símplices
$[v_0,\dots,v_i,w_i,\dots,w_q]$, $i=0,\dots,q$. Para
$\sigma:\Delta^q\to X$, definimos
\[
P(\sigma):=\sum_{i=0}^q(-1)^i\big(F\circ(\sigma\times\mathrm{id}_I)\big)\big|_{[v_0,\dots,v_i,w_i,\dots,w_q]}\ \in S_{q+1}(Y),
\]
extendido $R$-linealmente a $P:S_q(X)\to S_{q+1}(Y)$. Si $\sigma$ tiene
imagen en $A$, entonces $\sigma\times\mathrm{id}_I$ tiene imagen en
$A\times I$, y $F\circ(\sigma\times\mathrm{id}_I)$ tiene imagen en $B$ (ya
que $F(A\times I)\subset B$); así $P(S_q(A))\subset S_{q+1}(B)$, como se
requería.

\emph{Verificación de $\partial P+P\partial=g_\#-f_\#$.} Al calcular
$\partial$ de cada símplice $[v_0,\dots,v_i,w_i,\dots,w_q]$ en la suma que
define $P(\sigma)$ y sumar sobre $i$ con los signos $(-1)^i$, los términos
correspondientes a omitir un vértice interior de cada símplice se cancelan
en parejas consecutivas ($i$ e $i+1$, cada uno contribuye una vez con
signo $+$ y otra con signo $-$ al mismo símplice de codimensión $1$ con
$v_i=w_i$ colapsado), salvo dos: el símplice inferior
$[v_0,\dots,v_q]=\Delta^q\times\{0\}$, que aporta
$(F\circ(\sigma\times\mathrm{id}))|_{\Delta^q\times\{0\}}=f\circ\sigma$ con
signo $+$ (proviene del término $i=0$ de $P(\sigma)$, al omitir el último
vértice $w_0$ de $[v_0,w_0,\dots,w_q]$), y el símplice superior
$[w_0,\dots,w_q]=\Delta^q\times\{1\}$, que aporta $-g\circ\sigma$
(proviene del término $i=q$ de $P(\sigma)$, con signo $(-1)^q$, al omitir
el primer vértice $v_q$ de $[v_0,\dots,v_q,w_q]$, multiplicado por el
signo global $(-1)^{q+1}$ de esa cara al tomar $\partial$). Siguiendo
cuidadosamente los signos restantes (como en Hatcher, \emph{Algebraic
Topology}, Lema 2.10), se obtiene exactamente
\[
\partial(P(\sigma))=g\circ\sigma-f\circ\sigma-P(\partial\sigma),
\]
es decir $\partial P+P\partial=g_\#-f_\#$ sobre $S_q(X)$, que es la
identidad buscada.
\end{proof}

\begin{theorem}[Escisión]
\label{thm:escision}
Sean $X$ un espacio topológico y $A\subset X$ un subespacio. Sea $U\subset A$
tal que $\overline U\subset\mathring A$. La inclusión
$j:(X-U,A-U)\hookrightarrow(X,A)$, induce un isomorfismo
\[
H_q(X-U,A-U)\overset{\cong}{\to}H_q(X,A).
\]
\end{theorem}

\begin{proof}
Sea $B:=X\setminus U$ y $\mathcal U=\{A,B\}$. Veamos que
$\mathring A\cup\mathring B=X$: si $x\notin\mathring A$, entonces, como
$\overline U\subset\mathring A$, se tiene $x\notin\overline U$ (si
$x\in\overline U$ entonces $x\in\mathring A$), así
$x\in X\setminus\overline U$; pero $X\setminus\overline U$ es abierto
(complemento de un cerrado) y está contenido en $B=X\setminus U$, luego
$x\in\mathring B$. Sea $S_\bullet^{\mathcal U}(X)\subset S_\bullet(X)$ el
subcomplejo generado por los símplices singulares $\sigma:\Delta^q\to X$
cuya imagen está contenida en $A$ o en $B$ (los símplices
``$\mathcal U$-pequeños''); nótese
$S_\bullet^{\mathcal U}(X)=S_\bullet(A)+S_\bullet(B)$.

\emph{Lema de subdivisión baricéntrica (small simplices theorem).} Si
$\mathring A\cup\mathring B=X$, la inclusión
$\iota:S_\bullet^{\mathcal U}(X)\hookrightarrow S_\bullet(X)$ es una
equivalencia de homotopía de cadenas. La construcción usa un operador de
subdivisión baricéntrica $\mathrm{SD}:S_q(X)\to S_q(X)$ (un morfismo de
cadenas, $\mathrm{SD}\,\partial=\partial\,\mathrm{SD}$, homotópico a la
identidad vía un operador explícito $T$ con
$\partial T+T\partial=\mathrm{id}-\mathrm{SD}$, construido subdividiendo
cada símplice singular en símplices más pequeños por baricentros
iterados) junto con el siguiente hecho geométrico: para todo
$\sigma:\Delta^q\to X$, como $\{\sigma^{-1}(\mathring A),\sigma^{-1}(\mathring B)\}$
es una cubierta abierta del compacto $\Delta^q$ y el diámetro de las
piezas de la subdivisión baricéntrica iterada de $\Delta^q$ tiende a $0$,
existe $m=m(\sigma)$ tal que cada sumando de $\mathrm{SD}^m(\sigma)$ tiene
imagen en $A$ o en $B$. Con esto se construye, de manera estándar (véase
Hatcher, \emph{Algebraic Topology}, Proposición 2.21, cuya verificación
es puramente combinatoria y no usa hipótesis adicionales sobre $X$), un
operador $D:S_q(X)\to S_{q+1}(X)$ con
$\partial D+D\partial=\mathrm{id}-\rho\circ\iota$, donde
$\rho:S_\bullet(X)\to S_\bullet^{\mathcal U}(X)$ envía cada $\sigma$ a
$\mathrm{SD}^{m(\sigma)}(\sigma)\in S_\bullet^{\mathcal U}(X)$. Esto
exhibe a $\iota$ como equivalencia de homotopía de cadenas, con inversa
homotópica $\rho$.

Aceptando el lema (aplicable aquí, pues acabamos de verificar
$\mathring A\cup\mathring B=X$), la inclusión
$S_\bullet(A)+S_\bullet(B)=S_\bullet^{\mathcal U}(X)\hookrightarrow S_\bullet(X)$
induce isomorfismo en homología para toda $q$; ésta es precisamente la
condición de que la pareja $\{A,B\}$ sea \emph{escisiva} (Ejercicio 3 de
esta sección). Como $A\cup B=A\cup(X\setminus U)=X$ (pues $U\subset A$),
el Ejercicio 3 da que la inclusión $(B,A\cap B)\hookrightarrow(X,A)$
induce isomorfismo en homología para toda $q$. Pero $B=X\setminus U$ y
$A\cap B=A\setminus U$, así que esta inclusión es exactamente
$j:(X-U,A-U)\hookrightarrow(X,A)$, lo que concluye la demostración.
\end{proof}

\begin{definition}[Pareja buena]
\label{def:pareja-buena}
Sea $X$ un espacio y $A$ un subespacio cerrado de $X$ no vacío. Decimos que
la pareja $(X,A)$ es una \emph{buena pareja} si $A$ es un retracto por
deformación de alguna vecindad en $X$.
\end{definition}

\begin{example}
\label{ex:cw-buena-pareja}
Si $X$ es un complejo CW y $A$ es un subcomplejo no vacío, entonces $(X,A)$
es una buena pareja.
\end{example}

\begin{theorem}
\label{thm:sucesion-pareja-buena}
Si $(X,A)$ es una buena pareja, entonces existe una sucesión exacta
\[
\cdots\to\tilde H_q(A)\to\tilde H_q(X)\to\tilde H_q(X/A)\xrightarrow{\hat\partial_q}\tilde H_{q-1}(A)\to\cdots
\]
\end{theorem}

\begin{note}
\label{nota:cociente-pareja-buena}
El Teorema~\ref{thm:sucesion-pareja-buena} se sigue del
Teorema~\ref{thm:sucesion-larga} gracias al siguiente resultado.

Si $(X,A)$ es una buena pareja, entonces el cociente $q:(X,A)\to(X/A,A/A)$
induce isomorfismos
\[
H_q(X,A)\cong H_q(X/A,A/A)\cong\tilde H_q(X/A)
\]
para toda $q$.

En efecto, sea $V$ una vecindad de $A$ en $X$ (con $A\subset\mathring V$)
sobre la cual $A$ es un retracto por deformación (existe por hipótesis de
buena pareja); esto significa que la inclusión $A\hookrightarrow V$ es una
equivalencia de homotopía, con inversa homotópica la retracción, y ambas
compatibles con la pareja $(A,A)$, así que por el
Teorema~\ref{thm:invarianza-homotopica} aplicado a la pareja $(V,A)\simeq(A,A)$,
\[
H_q(V,A)\cong H_q(A,A)=0\qquad\text{para toda }q.
\]
Aplicando esto a la sucesión exacta larga de la tripleta $A\subset V\subset X$
(Ejercicio 4 de esta sección),
\[
\cdots\to H_{q+1}(X,V)\to H_q(V,A)\to H_q(X,A)\to H_q(X,V)\to H_{q-1}(V,A)\to\cdots,
\]
con $H_q(V,A)=0$ para toda $q$, se obtiene
$H_q(X,A)\cong H_q(X,V)$ para toda $q$.

Por otro lado, como $A$ es cerrado (parte de la definición de pareja
buena) y $A\subset\mathring V$, el Teorema de
escisión~\ref{thm:escision} (con $U=A$) da
\[
H_q(X-A,V-A)\cong H_q(X,V)
\]
para toda $q$.

La aplicación cociente $q:X\to X/A$ restringe a un homeomorfismo
$X-A\xrightarrow{\cong}(X/A)-(A/A)$ (es biyectiva, pues $q$ solo identifica
puntos de $A$, y es abierta/cerrada por ser $q$ una identificación),
llevando $V-A$ homeomórficamente sobre $(V/A)-(A/A)$; esto da
\[
H_q(X-A,V-A)\cong H_q\big(X/A-A/A,\,V/A-A/A\big).
\]

Repitiendo los dos primeros pasos con $(X,A)$ reemplazado por
$(X/A,A/A)$ ---la deformación de $V$ sobre $A$ desciende, vía $q$, a una
deformación de $V/A$ sobre $A/A$, y $A/A$ es cerrado en $X/A$, con
$A/A\subset\mathring{(V/A)}$--- se obtienen igualmente
\[
H_q(X/A,A/A)\cong H_q(X/A,V/A) \qquad\text{y}\qquad H_q\big(X/A-A/A,\,V/A-A/A\big)\cong H_q(X/A,V/A).
\]
Encadenando las cinco identificaciones,
\[
H_q(X,A)\cong H_q(X,V)\cong H_q(X-A,V-A)\cong H_q\big(X/A-A/A,V/A-A/A\big)\cong H_q(X/A,V/A)\cong H_q(X/A,A/A),
\]
y un chequeo directo sobre representantes de ciclos muestra que el
isomorfismo compuesto es precisamente $q_*$. Finalmente,
$H_q(X/A,A/A)\cong\tilde H_q(X/A)$ ya que $A/A=\{\ast\}$ es un punto: de
la sucesión exacta larga de la pareja $(X/A,\{\ast\})$,
$H_q(\{\ast\})=0$ para $q>0$ da $H_q(X/A,\{\ast\})\cong H_q(X/A)=\tilde H_q(X/A)$
para $q>0$, y para $q=0$ ambos lados son $0$ (pues $H_0(X/A,\{\ast\})=0$
por el Ejercicio 1e, y $\tilde H_0(X/A)$ mide las componentes conexas
además de la del punto base, coincidiendo con la definición de
$H_0(X/A,\{\ast\})$ vía el Ejercicio 1e aplicado en general).
\end{note}

\begin{proof}[Demostración del Teorema~\ref{thm:sucesion-pareja-buena}]
Basta sustituir, en la sucesión exacta larga reducida de la pareja
$(X,A)$ (Teorema~\ref{thm:sucesion-larga}), el término $H_q(X,A)$ por
$\tilde H_q(X/A)$ vía los isomorfismos de la Nota~\ref{nota:cociente-pareja-buena}
(naturales, por lo que conmutan con los morfismos de conexión y las
inclusiones de la sucesión).
\end{proof}

\begin{corollary}
\label{cor:union-subcomplejos}
Si $X$ es un complejo CW que es unión de subcomplejos $A$ y $B$, entonces
la inclusión $(B,A\cap B)\hookrightarrow(X,A)$ induce isomorfismos
\[
H_q(B,A\cap B)\cong H_q(X,A)
\]
para toda $q$.
\end{corollary}

\begin{proof}
Sea $U$ un entorno abierto de $A\cap B$ en $B$ que se retrae por
deformación sobre $A\cap B$ dentro de $B$ (existe pues $(B,A\cap B)$ es
buena pareja, Ejemplo~\ref{ex:cw-buena-pareja}, ya que $A\cap B$ es
subcomplejo de $B$). Extendiendo esta retracción por la identidad en
$A$, se obtiene una vecindad $A\cup U$ de $A$ en $X$ que se retrae por
deformación sobre $A$. Como $(X,A)$ y $(B,A\cap B)$ son ambas buenas
parejas (Ejemplo~\ref{ex:cw-buena-pareja}), y la composición
$B\hookrightarrow X\to X/A$ tiene fibra exactamente $A\cap B$ sobre el
punto base (pues $X=A\cup B$, así que un punto de $B$ colapsa al punto
base si y solo si pertenece a $A$, es decir a $A\cap B$) y es biyectiva
fuera de esa fibra, se obtiene un homeomorfismo
$B/(A\cap B)\xrightarrow{\cong}X/A$. Así, la
Nota~\ref{nota:cociente-pareja-buena} aplicada a ambas parejas da
\[
H_q(B,A\cap B)\cong\tilde H_q\big(B/(A\cap B)\big)\cong\tilde H_q(X/A)\cong H_q(X,A),
\]
y se verifica directamente que la composición de estos isomorfismos es la
inducida por la inclusión $(B,A\cap B)\hookrightarrow(X,A)$.
\end{proof}

\begin{corollary}
\label{cor:homologia-bouquet}
Sea $(X_\alpha,x_\alpha)$ una pareja buena para cada $\alpha$. La inclusión
$i_\alpha:X_\alpha\hookrightarrow\bigvee_\alpha X_\alpha$ induce un
isomorfismo
\[
\bigoplus_\alpha i_{\alpha*}:\bigoplus_\alpha\tilde H_q(X_\alpha)\to\tilde H_q\left(\bigvee_\alpha X_\alpha\right).
\]
\end{corollary}

\begin{proof}
Se sigue de la Nota~\ref{nota:cociente-pareja-buena} y del ejercicio 1
tomando $(X,A)=\big(\bigsqcup_\alpha X_\alpha,\bigsqcup_\alpha x_\alpha\big)$.
\end{proof}

\begin{theorem}[Mayer--Vietoris]
\label{thm:mayer-vietoris}
Sea $X=A\cup B$. Sean $A_1\subset A$ y $B_1\subset B$. Supongamos que los
pares $\{A,B\}$ y $\{A_1,B_1\}$ son escisivos. Entonces, se tiene la
sucesión exacta de homología relativa
\[
\cdots\to H_{q+1}(X,A_1\cup B_1)\to H_q(A\cap B,A_1\cap B_1)\xrightarrow{(i_*,-j_*)}H_q(A,A_1)\oplus H_q(B,B_1)\xrightarrow{k_*+l_*}H_q(X,A_1\cup B_1)\to H_{q-1}(A\cap B,A_1\cap B_1)\to\cdots
\]
en donde $i:(A\cap B,A_1\cap B_1)\hookrightarrow(A,A_1)$,
$j:(A\cap B,A_1\cap B_1)\hookrightarrow(B,B_1)$,
$k:(A,A_1)\hookrightarrow(X,A_1\cup B_1)$ y
$l:(B,B_1)\hookrightarrow(X,A_1\cup B_1)$.

En particular, si $A_1=B_1=\{*\}$, se obtiene la sucesión exacta
\[
\cdots\to H_{q+1}(X)\to H_q(A\cap B)\to H_q(A)\oplus H_q(B)\to H_q(X)\to H_{q-1}(A\cap B)\to\cdots
\]
\end{theorem}

\begin{proof}
En el nivel de cadenas, defínase
\[
\phi:S_\bullet(A\cap B)/S_\bullet(A_1\cap B_1)\to \big(S_\bullet(A)/S_\bullet(A_1)\big)\oplus\big(S_\bullet(B)/S_\bullet(B_1)\big), \qquad c\mapsto(i_\#c,\,-j_\#c),
\]
\[
\psi:\big(S_\bullet(A)/S_\bullet(A_1)\big)\oplus\big(S_\bullet(B)/S_\bullet(B_1)\big)\to\big(S_\bullet(A)+S_\bullet(B)\big)/\big(S_\bullet(A_1)+S_\bullet(B_1)\big), \qquad (a,b)\mapsto a+b.
\]
La sucesión $0\to\bullet\xrightarrow{\phi}\bullet\xrightarrow{\psi}\bullet\to0$
es exacta en cada grado: $\phi$ es inyectiva porque si
$i_\#c\in S_\bullet(A_1)$ y $j_\#c\in S_\bullet(B_1)$ entonces
$c\in S_\bullet(A_1)\cap S_\bullet(B_1)=S_\bullet(A_1\cap B_1)$ (una cadena
de $A\cap B$ está en $A_1$ y en $B_1$ si y solo si está en
$A_1\cap B_1$, comparando símplice a símplice); $\psi$ es sobreyectiva por
definición de $S_\bullet(A)+S_\bullet(B)$; y
$\mathrm{im}\,\phi=\ker\psi$ porque $\psi(i_\#c,-j_\#c)=i_\#c-j_\#c=0$
módulo $S_\bullet(A_1)+S_\bullet(B_1)$ (ya que $i_\#c=j_\#c$ como cadenas
de $X$), y recíprocamente si $\psi(a,b)=a+b\in S_\bullet(A_1)+S_\bullet(B_1)$,
escribiendo $a+b=a_1+b_1$ con $a_1\in S_\bullet(A_1)$, $b_1\in S_\bullet(B_1)$,
se tiene $a-a_1=b_1-b\in S_\bullet(A)\cap S_\bullet(B)=S_\bullet(A\cap B)$
(pues $a-a_1\in S_\bullet(A)$ y $b_1-b\in S_\bullet(B)$ son la misma
cadena), y esta cadena $c:=a-a_1$ satisface $\phi(c)=(a,-b)=(a,b)$ módulo
los subcomplejos correspondientes.

Por el lema del zig-zag (como en la demostración del
Teorema~\ref{thm:sucesion-larga}), esta sucesión corta exacta de
complejos induce una sucesión exacta larga
\[
\cdots\to H_{q+1}\Big(\tfrac{S_\bullet(A)+S_\bullet(B)}{S_\bullet(A_1)+S_\bullet(B_1)}\Big)\to H_q(A\cap B,A_1\cap B_1)\xrightarrow{(i_*,-j_*)}H_q(A,A_1)\oplus H_q(B,B_1)\xrightarrow{\psi_*}H_q\Big(\tfrac{S_\bullet(A)+S_\bullet(B)}{S_\bullet(A_1)+S_\bullet(B_1)}\Big)\to\cdots
\]
Como $\{A,B\}$ y $\{A_1,B_1\}$ son escisivos, las inclusiones
$S_\bullet(A)+S_\bullet(B)\hookrightarrow S_\bullet(X)$ y
$S_\bullet(A_1)+S_\bullet(B_1)\hookrightarrow S_\bullet(A_1\cup B_1)$
inducen isomorfismos en homología; por el lema de los cinco aplicado a las
sucesiones exactas largas de las parejas
$\big(S_\bullet(A)+S_\bullet(B),\,S_\bullet(A_1)+S_\bullet(B_1)\big)$ y
$(X,A_1\cup B_1)$ (comparadas vía la inclusión natural, usando que
$S_\bullet(A_1)+S_\bullet(B_1)\subset S_\bullet(A_1\cup B_1)$ también
induce isomorfismo en homología por ser union de dos escisiones
sucesivas), se obtiene
\[
H_q\Big(\tfrac{S_\bullet(A)+S_\bullet(B)}{S_\bullet(A_1)+S_\bullet(B_1)}\Big)\cong H_q(X,A_1\cup B_1).
\]
Sustituyendo esta identificación (natural, por construcción) en la
sucesión exacta larga de arriba, y verificando que
$\psi_*$ corresponde a $k_*+l_*$ bajo ella (pues $\psi$ es, en cadenas,
$a+b\mapsto$ su clase en $S_\bullet(X)/S_\bullet(A_1\cup B_1)$, que es
exactamente $k_\#(a)+l_\#(b)$), se obtiene la sucesión del enunciado.

La sucesión particular con $A_1=B_1=\{*\}$ se sigue tomando ese caso
específico, usando $H_q(A,\{*\})=H_q(A)$ etc.\ vía la identificación
usual con la homología reducida cuando corresponde.
\end{proof}

\begin{note}
\label{nota:mayer-vietoris-abiertos}
Si $U,V\hookrightarrow X$ son abiertos y $X=U\cup V$, entonces la pareja
$\{U,V\}$ es escisiva. Entonces, se tiene la sucesión exacta:
\[
\cdots\to H_{q+1}(X)\to H_q(U\cap V)\to H_q(U)\oplus H_q(V)\to H_q(X)\to H_{q-1}(U\cap V)\to\cdots
\]
\end{note}

\begin{theorem}[Tres incisos]
\label{thm:tres-incisos}
Sean $n\ge1$ un entero, $Y$ un espacio topológico y $f:S^{n-1}\to Y$ una
aplicación continua. Sea $X=Y\cup_fD^n$. Entonces,
\begin{enumerate}[label=\arabic*.]
\item $\tilde H_q(X)\cong\tilde H_q(Y)$ para todo entero $q\ne n,n-1$;
\item la sucesión $0\to\im f_*\to\tilde H_{n-1}(Y)\to\tilde H_{n-1}(X)\to0$
es exacta;
\item la sucesión $0\to\tilde H_n(Y)\to\tilde H_n(X)\to\ker f_*\to0$ es
exacta
\end{enumerate}
en donde $f_*:\tilde H_{n-1}(S^{n-1})\to\tilde H_{n-1}(Y)$.
\end{theorem}

\begin{proof}
Sea $\Phi:D^n\to X$ la aplicación característica de la celda adjuntada
(la composición $D^n\hookrightarrow Y\sqcup D^n\twoheadrightarrow X$), que
restringe a $f$ en $S^{n-1}$. El par $(X,Y)$ es una buena pareja
--- $Y$ es cerrado, y una vecindad-collar de $Y$ (imagen del anillo
$\{x\in D^n:\|x\|\ge1/2\}$ bajo $\Phi$, unido a $Y$) se retrae por
deformación sobre $Y$ --- así que, por la
Nota~\ref{nota:cociente-pareja-buena},
\[
H_q(X,Y)\cong\tilde H_q(X/Y).
\]
Pero $X/Y\cong D^n/S^{n-1}\cong S^n$: colapsar $Y$ en $X=Y\cup_fD^n$ es,
por construcción de la adjunción, exactamente colapsar $S^{n-1}$ en
$D^n$. Así, por el Ejercicio 6 de esta sección,
\[
H_q(X,Y)\cong\tilde H_q(S^n)=\begin{cases}R & q=n\\0 & q\ne n.\end{cases}
\]
De la sucesión exacta larga reducida de la pareja $(X,Y)$
(Teorema~\ref{thm:sucesion-larga}),
\[
\cdots\to H_{q+1}(X,Y)\to\tilde H_q(Y)\to\tilde H_q(X)\to H_q(X,Y)\to\tilde H_{q-1}(Y)\to\cdots,
\]
se sigue de inmediato la afirmación 1: para $q\ne n,n-1$, tanto
$H_q(X,Y)$ como $H_{q+1}(X,Y)$ se anulan (pues $q\ne n$ y $q+1\ne n$), así
$\tilde H_q(Y)\cong\tilde H_q(X)$.

Para 2 y 3, aíslese el segmento de la sucesión en torno a $q=n-1,n$:
\[
0=H_{n+1}(X,Y)\to\tilde H_n(Y)\to\tilde H_n(X)\xrightarrow{\ \partial\ }H_n(X,Y)\xrightarrow{\ \partial\ }\tilde H_{n-1}(Y)\to\tilde H_{n-1}(X)\to H_{n-1}(X,Y)=0.
\]
Por naturalidad del morfismo conectante aplicada a la aplicación de
parejas $\Phi:(D^n,S^{n-1})\to(X,Y)$, el siguiente diagrama conmuta:
\[
\begin{tikzcd}
H_n(D^n,S^{n-1}) \arrow[r,"\delta","\cong"'] \arrow[d,"\Phi_*","\cong"'] & \tilde H_{n-1}(S^{n-1}) \arrow[d,"f_*"] \\
H_n(X,Y) \arrow[r,swap,"\partial"] & \tilde H_{n-1}(Y)
\end{tikzcd}
\]
en donde $\delta$ es el isomorfismo de la sucesión reducida de
$(D^n,S^{n-1})$ (análogo al usado arriba para $(X,Y)$: $(D^n,S^{n-1})$ es
también buena pareja, vía el mismo collar $\{1/2\le\|x\|\le1\}$, y
$D^n/S^{n-1}\cong S^n$) y $\Phi_*$ es isomorfismo porque, bajo las dos
identificaciones $H_n(D^n,S^{n-1})\cong\tilde H_n(D^n/S^{n-1})$ y
$H_n(X,Y)\cong\tilde H_n(X/Y)$ de la Nota~\ref{nota:cociente-pareja-buena},
$\Phi_*$ corresponde exactamente al isomorfismo
$\tilde H_n(D^n/S^{n-1})\cong\tilde H_n(S^n)\cong\tilde H_n(X/Y)$ inducido
por el homeomorfismo $D^n/S^{n-1}\xrightarrow{\cong}X/Y$ ya usado en la
identificación de $H_q(X,Y)$ al inicio de esta demostración (pues $\Phi$
desciende exactamente a ese homeomorfismo al pasar a los cocientes). Como
el cuadrado conmuta y las flechas verticales son
isomorfismos, $\partial=f_*\circ\Phi_*^{-1}\circ\delta$ tiene, bajo la
identificación $H_n(X,Y)\cong\tilde H_{n-1}(S^{n-1})$ dada por
$\delta\circ\Phi_*^{-1}$, exactamente la misma imagen y el mismo núcleo
que $f_*:\tilde H_{n-1}(S^{n-1})\to\tilde H_{n-1}(Y)$.

Para la afirmación 2 ($q=n-1$): de
$H_n(X,Y)\xrightarrow{\partial}\tilde H_{n-1}(Y)\to\tilde H_{n-1}(X)\to0$
exacta, se tiene $0\to\im\partial\to\tilde H_{n-1}(Y)\to\tilde H_{n-1}(X)\to0$
exacta, y $\im\partial=\im f_*$ por lo anterior.

Para la afirmación 3 ($q=n$): de
$0\to\tilde H_n(Y)\to\tilde H_n(X)\xrightarrow{\partial}H_n(X,Y)$ exacta,
se tiene que $\tilde H_n(X)/\tilde H_n(Y)\cong\im\partial=\ker(\tilde H_{n-1}(Y)\to\tilde H_{n-1}(X))=\ker f_*$
(la última igualdad por exactitud en $H_n(X,Y)$, ya que
$\im\partial=\ker(H_n(X,Y)\to\tilde H_{n-1}(Y))$, y bajo la identificación
de arriba este núcleo corresponde a $\ker f_*$). Esto da
$0\to\tilde H_n(Y)\to\tilde H_n(X)\to\ker f_*\to0$ exacta.
\end{proof}

\begin{theorem}
\label{thm:cw-celular}
Sea $X$ un complejo CW. Sea $\Gamma^n=\{n\text{-células de }X\}$. Entonces,
\begin{enumerate}[label=\arabic*.]
\item $H_q(X^n,X^{n-1})=0$ para $q\ne n$ y es abeliano libre para $q=n$,
con base en correspondencia uno a uno con $\Gamma^n$.
\item $H_q(X^n)=0$ para $q>n$. En particular, si $X$ es de dimensión
finita, $H_q(X)=0$ para $q>\dim X$.
\item La inclusión $i:X^n\hookrightarrow X$ induce un isomorfismo
$i_*:H_q(X^n)\to H_q(X)$ si $q<n$.
\end{enumerate}

Sea $X$ un complejo CW. Entonces,
\[
\cdots\to H_{n+1}(X^{n+1},X^n)\xrightarrow{d_{n+1}}H_n(X^n,X^{n-1})\xrightarrow{d_n}H_{n-1}(X^{n-1},X^{n-2})\to\cdots
\]
en donde $d_{n+1}=j_{n*}\partial_{n+1}$ y $d_n=j_{n-1*}\partial_n$ con
$\partial_{n+1}:H_{n+1}\to H_n(X^n)$ y
$j_{n*}:H_n(X^n)\to H_n(X^n,X^{n-1})$.

Ya que $d_{n+1}\circ d_n=0$, la sucesión anterior forma un complejo de
cadenas llamado \emph{complejo de cadenas celular} de $X$. Los grupos de
homología de este complejo de cadenas celular se llaman \emph{grupos de
homología celular} de $X$.
\end{theorem}

\begin{proof}
\emph{1.} $X^n$ se obtiene de $X^{n-1}$ adjuntando simultáneamente las
celdas de $\Gamma^n$ vía $\bigsqcup_\alpha S^{n-1}_\alpha\to X^{n-1}$; el
Teorema~\ref{thm:tres-incisos} se extiende sin cambios a este caso
(adjuntar varias celdas a la vez), dando --- exactamente como en su
demostración, con $\bigvee_\alpha S^n_\alpha$ en lugar de $S^n$ ---
$H_q(X^n,X^{n-1})\cong\tilde H_q\big(\bigvee_\alpha S^n_\alpha\big)$, que
por el Corolario~\ref{cor:homologia-bouquet} es
$\bigoplus_\alpha\tilde H_q(S^n_\alpha)$: libre con base $\Gamma^n$ si
$q=n$, y $0$ si $q\ne n$.

\emph{2.} Por inducción en $n$. Para $n=0$, $X^0$ es discreto y
$H_q(X^0)=\bigoplus_{x\in X^0}H_q(\{x\})=0$ para $q>0$ (Ejercicio 1c).
Supuesto el resultado para $X^{n-1}$ (es decir, $H_q(X^{n-1})=0$ para
$q>n-1$), sea $q>n$; en particular $q\ne n,n-1$ (pues $q>n>n-1$), así el
inciso 1 del Teorema~\ref{thm:tres-incisos} (extendido a varias celdas,
como arriba) da $\tilde H_q(X^n)\cong\tilde H_q(X^{n-1})=0$ (usando la
hipótesis de inducción, ya que $q>n>n-1$). Si $X$ tiene dimensión
finita $N$ (es decir $X=X^N$), entonces $H_q(X)=H_q(X^N)=0$ para $q>N$
por lo anterior.

\emph{3.} Sea $q<n$. Para cualquier $k\ge n$ (en particular $k\ge q+1$,
es decir $q<k$ estrictamente), el inciso 1 (extendido a varias celdas) da
$H_j(X^{k+1},X^k)=0$ para todo $j\ne k+1$; en particular, como $q<k$
implica $q\ne k+1$ y $q+1\le k\ne k+1$, se anulan tanto
$H_q(X^{k+1},X^k)$ como $H_{q+1}(X^{k+1},X^k)$. De la sucesión exacta
larga de $(X^{k+1},X^k)$,
\[
0=H_{q+1}(X^{k+1},X^k)\to H_q(X^k)\to H_q(X^{k+1})\to H_q(X^{k+1},X^k)=0,
\]
se obtiene que $H_q(X^k)\to H_q(X^{k+1})$ es isomorfismo. Iterando para
$k=n,n+1,n+2,\dots$, se obtiene que $H_q(X^n)\to H_q(X^m)$ es isomorfismo
para todo $m>n$. Falta pasar al límite $m\to\infty$ (o a $X$, si $X$ es de
dimensión infinita): como $\Delta^q$ es compacto, la imagen de cualquier
símplice singular
$\sigma:\Delta^q\to X$ es compacta, y toda parte compacta de un complejo
CW está contenida en un subcomplejo finito (por inducción en el esqueleto:
un compacto interseca solo finitas celdas abiertas, pues de lo contrario
tendría un punto en cada una de infinitas celdas, formando un conjunto
infinito sin puntos de acumulación en el compacto, ya que cada celda
abierta es un entorno de sus propios puntos disjunto de las demás). Por lo
tanto, toda cadena singular en $X$ tiene imagen en algún $X^m$ con $m$
finito, lo que muestra que
$S_\bullet(X)=\varinjlim_mS_\bullet(X^m)$ y, ya que la homología conmuta
con límites directos de complejos de cadenas (los ciclos y fronteras de
$X$ en grado $q$ son unión creciente de los de cada $X^m$),
$H_q(X)=\varinjlim_mH_q(X^m)$. Como $H_q(X^n)\to H_q(X^m)$ es isomorfismo
para todo $m\ge n$ (por lo mostrado, o trivialmente si $X$ es de
dimensión $\le n$), se sigue que $H_q(X^n)\to H_q(X)$ es isomorfismo.

\emph{Complejo celular y $d^2=0$.} Por definición,
$d_n\circ d_{n+1}=j_{n-1*}\partial_n\circ j_{n*}\partial_{n+1}$. Los
morfismos $j_{n*}:H_n(X^n)\to H_n(X^n,X^{n-1})$ y
$\partial_n:H_n(X^n,X^{n-1})\to H_{n-1}(X^{n-1})$ son dos morfismos
\emph{consecutivos} en la sucesión exacta larga de la pareja
$(X^n,X^{n-1})$ (Teorema~\ref{thm:sucesion-larga}); por exactitud,
$\im(j_{n*})=\ker(\partial_n)$, así $\partial_n\circ j_{n*}=0$
idénticamente. Por lo tanto $d_n\circ d_{n+1}=j_{n-1*}\circ0\circ\partial_{n+1}=0$.
\end{proof}

\begin{theorem}
\label{thm:cw-iso-singular}
$H_n^{\CWq}(X)\cong H_n(X)$.
\end{theorem}

\begin{proof}
Escribamos $C_k:=H_k(X^k,X^{k-1})$, así que el complejo celular es
$(\mathcal C_\bullet,d_\bullet)$ con $d_k=j_{k-1*}\partial_k$.

\emph{Paso 1: $H_n(X^n)\cong\ker(d_n)$.} Como $H_n(X^{n-1})=0$
(Teorema~\ref{thm:cw-celular}.2, ya que $n>n-1$), la sucesión exacta larga
de $(X^n,X^{n-1})$ da que $j_{n*}:H_n(X^n)\to C_n$ es \emph{inyectiva}
(su núcleo es la imagen de $H_n(X^{n-1})=0$). Además,
\[
\ker(d_n)=\ker(j_{n-1*}\partial_n)=\ker(\partial_n)
\]
ya que $j_{n-1*}:H_{n-1}(X^{n-1})\to C_{n-1}$ es también inyectiva (mismo
argumento, con $n-1$ en el papel de $n$: $H_{n-1}(X^{n-2})=0$). Y
$\ker(\partial_n)=\im(j_{n*})$ por exactitud en $H_n(X^n,X^{n-1})$ de la
sucesión larga de $(X^n,X^{n-1})$. Por lo tanto
$\ker(d_n)=\im(j_{n*})$, y como $j_{n*}$ es inyectiva,
\[
j_{n*}:H_n(X^n)\xrightarrow{\cong}\ker(d_n).
\]

\emph{Paso 2: $H_n(X^{n+1})\cong H_n(X^n)/\im(\partial_{n+1})$.} De la
sucesión exacta larga de $(X^{n+1},X^n)$,
\[
H_{n+1}(X^{n+1},X^n)\xrightarrow{\partial_{n+1}}H_n(X^n)\to H_n(X^{n+1})\to H_n(X^{n+1},X^n)=0
\]
(el último término se anula por el Teorema~\ref{thm:cw-celular}.1, ya que
$n\ne n+1$), de donde $H_n(X^n)\to H_n(X^{n+1})$ es sobreyectiva con
núcleo $\im(\partial_{n+1})$, dando el isomorfismo del paso 2. Bajo el
isomorfismo $j_{n*}$ del Paso 1 (inyectivo con imagen $\ker(d_n)$), la
imagen de $\partial_{n+1}$ dentro de $H_n(X^n)$ corresponde exactamente a
$j_{n*}(\im\partial_{n+1})=\im(j_{n*}\partial_{n+1})=\im(d_{n+1})$ dentro
de $\ker(d_n)$ (la igualdad usa que $j_{n*}$ es inyectiva, así que
conjuga imágenes fielmente). Por lo tanto
\[
H_n(X^{n+1})\cong H_n(X^n)/\im(\partial_{n+1})\cong\ker(d_n)/\im(d_{n+1})=H_n^{\CWq}(X).
\]

\emph{Paso 3: conclusión.} Por el Teorema~\ref{thm:cw-celular}.3 (con
$q=n<n+1$), $H_n(X^{n+1})\cong H_n(X)$. Combinando con el Paso 2,
\[
H_n(X)\cong H_n(X^{n+1})\cong H_n^{\CWq}(X).\qedhere
\]
\end{proof}

\begin{remark}
\label{obs:cw-observaciones}
Sea $X$ un complejo CW.
\begin{enumerate}[label=\arabic*.]
\item Si $\Gamma^n=\varnothing$, entonces $H_n(X)=0$.
\item Si $X$ es conexo y $\#\Gamma^0=1$, entonces $d_1:\mathcal C_1\to\mathcal C_0$
es el mapeo cero.
\item Si $X$ no tiene dos de sus células en dimensiones adyacentes,
entonces $d_n=0$ para toda $n$ y $H_n(X)\cong\ZZ^{\#\Gamma^n}$.
\end{enumerate}
\end{remark}

\begin{theorem}
\label{thm:grado-pegado}
Sea $d_{\alpha\beta}=\deg(S^{n-1}_\alpha\to X^{n-1}\to S^{n-1}_\beta)$ que es
el grado de la composición de la aplicación de pegado de $e^n_\alpha$ con la
aplicación cociente que colapsa $X^{n-1}-e^{n-1}_\beta$ a un punto.
Entonces,
\[
d_n(e^n_\alpha)=\sum_\beta d_{\alpha\beta}e_\beta^{n-1}.
\]
\end{theorem}

\begin{proof}
Sean $\Phi_\alpha:D^n_\alpha\to X^n$ la aplicación característica de
$e^n_\alpha$ y $\varphi_\alpha:=\Phi_\alpha|_{S^{n-1}_\alpha}:S^{n-1}_\alpha\to X^{n-1}$
su aplicación de pegado. Por construcción de la base del
Teorema~\ref{thm:cw-celular}.1, la clase $e^n_\alpha\in H_n(X^n,X^{n-1})$
es $(\Phi_\alpha)_*(g_\alpha)$, donde $g_\alpha$ es el generador de
$H_n(D^n_\alpha,S^{n-1}_\alpha)\cong R$ que corresponde, bajo el
isomorfismo $\delta$ de la sucesión reducida de $(D^n_\alpha,S^{n-1}_\alpha)$
(Ejercicio 7c), al generador fundamental de
$\tilde H_{n-1}(S^{n-1}_\alpha)$.

Por naturalidad del morfismo conectante $\partial$ aplicada a
$\Phi_\alpha:(D^n_\alpha,S^{n-1}_\alpha)\to(X^n,X^{n-1})$ (que restringe a
$\varphi_\alpha$ en el borde),
\[
\partial_n(e^n_\alpha)=\partial_n\big((\Phi_\alpha)_*(g_\alpha)\big)=(\varphi_\alpha)_*\big(\delta(g_\alpha)\big)=(\varphi_\alpha)_*[S^{n-1}_\alpha]\ \in\ H_{n-1}(X^{n-1}),
\]
en donde $[S^{n-1}_\alpha]$ denota el generador fundamental de
$\tilde H_{n-1}(S^{n-1}_\alpha)$.

Falta calcular $j_{n-1*}\big((\varphi_\alpha)_*[S^{n-1}_\alpha]\big)\in H_{n-1}(X^{n-1},X^{n-2})$
en la base $\{e^{n-1}_\beta\}_\beta$. Por la Nota~\ref{nota:cociente-pareja-buena}
y el Teorema~\ref{thm:cw-celular}.1, la identificación
$H_{n-1}(X^{n-1},X^{n-2})\cong\tilde H_{n-1}(X^{n-1}/X^{n-2})$ manda
$e^{n-1}_\beta$ al generador fundamental de la esfera $S^{n-1}_\beta$ en
el bouquet $X^{n-1}/X^{n-2}\cong\bigvee_\beta S^{n-1}_\beta$
(Corolario~\ref{cor:homologia-bouquet}), y bajo esta identificación, la
$\beta$-ésima coordenada de la clase de un elemento
$y\in H_{n-1}(X^{n-1})$ en $H_{n-1}(X^{n-1},X^{n-2})=\bigoplus_\beta R\,e^{n-1}_\beta$
es exactamente el grado de la composición
$S^{n-1}_\gamma\to X^{n-1}\to X^{n-1}/X^{n-2}\to S^{n-1}_\beta$ (colapsar
todas las esferas del bouquet salvo la $\beta$-ésima) aplicada a un
representante esférico de $y$ --- que es, para $y=(\varphi_\alpha)_*[S^{n-1}_\alpha]$,
precisamente el grado de
\[
S^{n-1}_\alpha\xrightarrow{\ \varphi_\alpha\ }X^{n-1}\longrightarrow X^{n-1}/(X^{n-1}-e^{n-1}_\beta)\cong S^{n-1}_\beta,
\]
es decir $d_{\alpha\beta}$ (la aplicación cociente que colapsa
$X^{n-1}-e^{n-1}_\beta$ es exactamente, salvo homeomorfismo, la
composición de colapsar $X^{n-2}$ y luego todas las esferas del bouquet
salvo la $\beta$-ésima). Por lo tanto,
\[
j_{n-1*}\big((\varphi_\alpha)_*[S^{n-1}_\alpha]\big)=\sum_\beta d_{\alpha\beta}\,e^{n-1}_\beta,
\]
y así $d_n(e^n_\alpha)=j_{n-1*}\partial_n(e^n_\alpha)=\sum_\beta d_{\alpha\beta}e^{n-1}_\beta$.
\end{proof}

\subsection*{Ejercicios}

\begin{enumerate}[label=\arabic*.,leftmargin=*]

\item (\textbf{Propiedades de homología singular}) Sea $X$ un espacio.
\begin{enumerate}[label=\alph*)]
\item Si $X=\{*\}$, entonces
\[
H_q(\{*\})=\begin{cases}R & q=0\\0 & q\ne0.\end{cases}
\]
\item Si $X$ un espacio conexo por arcos no vacío. Entonces,
$H_0(X)\cong R$.
\item Sea $\{X_\alpha\}_{\alpha\in\Lambda}$ el conjunto de componentes por
arcos de $X$, entonces la inclusión $i_\alpha:X_\alpha\to X$ induce el
isomorfismo $\bigoplus_\alpha i_{\alpha*}:\bigoplus_\alpha H_q(X_\alpha)\to H_q(X)$.
\item $H_0(X)$ es isomorfo a una suma de copias de $R$, una por cada
componente por arcos de $X$.
\item Sea $X$ conexo por arcos y $A\subset X$ un subespacio no vacío.
Entonces, $H_0(X,A)=0$.
\item Sea $(X,A)$ una pareja y $\{X_\alpha\}_{\alpha\in\Lambda}$ la familia
de componentes por arcos de $X$. Sea $A_\alpha=X_\alpha\cap A$. Entonces,
las inclusiones $i_\alpha:(X_\alpha,A_\alpha)\to(X,A)$ inducen el
isomorfismo $\bigoplus_\alpha i_{\alpha*}:\bigoplus_\alpha H_q(X_\alpha,A_\alpha)\to H_q(X,A)$.
\item $H_0(X,A)$ es isomorfo a la suma directa de copias de $R$, una copia
por cada componente por arcos de $X$ tal que su intersección con $A$ es
vacía.
\end{enumerate}

\noindent\textbf{Solución.}
\begin{enumerate}[label=\alph*)]
\item Consideremos la sucesión
\[
\cdots\xrightarrow{\partial_{q+1}}S_q(\{*\})=R\xrightarrow{\partial_q}S_{q-1}(\{*\})=R\xrightarrow{\partial_{q-1}}\cdots
\]
Entonces,
\[
\partial_q(\sigma_q)=\sum_{i=0}^q(-1)^i\sigma_q\circ F_q^i=\sum_{i=0}^q(-1)^i\sigma_{q-1}
=\begin{cases}\sigma_{q-1} & q>0\text{ par}\\0 & q\text{ impar.}\end{cases}
\]
Por lo tanto,
\[
Z_q(\{*\})=\ker\partial_q=\begin{cases}0 & q>0\text{ par}\\R & q\text{ impar}\end{cases}
\qquad
B_q(\{*\})=\im\partial_q=\begin{cases}0 & q>0\text{ par}\\R & q\text{ impar.}\end{cases}
\]
Entonces, $H_q(\{*\})=0$, $q>0$. Además,
$H_0(\{*\})=\ker\partial_0/\im\partial_1=S_0(\{*\})/0=R$.
\item Dar $\sigma:\Delta_0=\{*\}\to X\in S_0(X)$ es equivalente a dar un
punto en $X$. Por lo tanto, $S_0(X)$ es el $R$-módulo libre generado por
los puntos de $X$. Así, cada elemento de $S_0(X)$ se puede escribir como
$\sum_{x\in X}r_xx$, $r_x\in R$. Defínase,
\[
\epsilon:S_0(X)\to R, \qquad \sum_{x\in X}r_xx\mapsto\sum_{x\in X}r_x,
\]
que es suprayectivo ya que si $r\in R$, $\epsilon(rx)=r$ para $x\in X$.
Queremos ver que $B_0(X)=\im\partial_1=\ker\epsilon$.

Sea $c\in B_0(X)$. Entonces, $c=\partial_1\left(\sum_\sigma r_\sigma\sigma\right)$,
donde $\sum_\sigma r_\sigma\sigma\in S_1(X)$ y $\sigma:\Delta_1\to X$. Así,
\[
c=\sum_\sigma r_\sigma\partial_1(\sigma)=\sum_\sigma r_\sigma(\sigma(e_1)-\sigma(e_0))
=\sum_\sigma r_\sigma\sigma(e_1)-\sum_\sigma r_\sigma\sigma(e_0).
\]
Por lo tanto $\epsilon(c)=0$ y $B_0(X)\subset\ker\epsilon$.

Ahora, sea $\sum r_xx\in\ker(\epsilon)$; entonces $\sum r_x=0$. Se quiere
probar que esta cadena está en $B_0(X)$. Sea $\sigma_x:\Delta_1\to X$
continua tal que $\sigma_x(e_0)=x_0$ y $\sigma_x(e_1)=x_1$. Consideremos
$c=\sum_xr_x\sigma_x\in S_1(X)$, entonces
\[
\partial_1(c)=\sum_xr_x\partial_1(\sigma_x)=\sum_xr_x(x-x_0)=\sum_xr_xx.
\]
Por lo tanto, $\sum r_xx\in B_1(X)$. De esta manera, se cumple que
\[
H_0(X)=Z_0(X)/B_0(X)=S_0(X)/\ker(\epsilon)\cong R.
\]
\item Una base para $\bigoplus_\alpha S_q(X_\alpha)$ está dada por
$\bigsqcup_\alpha s_q(X_\alpha)$ en donde
$s_q(X_\alpha)=\{\sigma:\Delta_q\to X_\alpha\mid\sigma\text{ es continua}\}$.
Una base para $S_q(X)$ es $s_q(X)$. Defínase,
\[
\bigsqcup_\alpha\iota_\alpha:\bigsqcup_\alpha s_q(X_\alpha)\to s_q(X)
\]
donde $\iota_\alpha(\sigma)=i_\alpha\circ\sigma$.

Se tiene que $\bigsqcup_\alpha\iota_\alpha$ es suprayectiva. En efecto, sea
$\tau:\Delta_q\to X\in s_q(X)$. Como $\Delta_q$ es conexo por arcos y $\tau$
es continua, $\tau(\Delta_q)\subset X_\alpha$ para alguna $\alpha$. Así,
existe $\tau'\in s_q(X_\alpha)$ tal que $\iota_\alpha(\tau')=\tau$.

Ahora, sean $\sigma_1,\sigma_2\in s_q(X_\alpha)$ tales que
$\iota_\alpha(\sigma_1)=\iota_\alpha(\sigma_2)$. Entonces,
$i_\alpha\circ\sigma_1=i_\alpha\circ\sigma_2$. Como $i_\alpha$ es inyectiva
$\sigma_1=\sigma_2$. Supongamos ahora que $\sigma_1\in s_q(X_\alpha)$ y
$\sigma_2\in s_q(X_\beta)$ con $\alpha\ne\beta$ cumplen que
$\iota_\alpha(\sigma_1)=\iota_\beta(\sigma_2)$. Entonces
$i_\alpha\circ\sigma_1=i_\beta\circ\sigma_2$. Entonces, $\sigma_1(t)=\sigma_2(t)$
para toda $t\in\Delta_q$. Entonces, $X_\alpha\cap X_\beta\ne\varnothing$ y
$\alpha=\beta$. Contradicción.

Por lo tanto, se tiene un isomorfismo
$\bigoplus_\alpha\iota_\alpha:\bigoplus_\alpha S_q(X_\alpha)\cong S_q(X)$.
Así,
\[
H_q(X)=H_q\left(\bigoplus_\alpha S_\bullet(X_\alpha)\right)
\cong\bigoplus_\alpha H_q(S_\bullet(X_\alpha))
=\bigoplus_\alpha H_q(X_\alpha).
\]
\item Por el inciso anterior, si $\{X_\alpha\}_\alpha$ es la familia de
componentes por arcos de $X$,
\[
H_0(X)\cong\bigoplus_\alpha H_0(X_\alpha)\cong\bigoplus_\alpha R.
\]
\item Se tiene que $H_0(X,A)=S_0(X)/B_0(X,A)$. Veamos que
$S_0(X)\subset B_0(X,A)$. Sea $c=\sum_xr_xx\in S_0(X)$ y sea $a_0\in A$.
Para cada $x\in X$, sea $\sigma_x\to X$ con $\sigma_x(e_0)=a_0$ y
$\sigma_x(e_1)=x$. Consideremos la $1$-cadena $\sum_xr_x\sigma_x\in S_1(X)$.
Entonces,
\[
\partial_1\left(\sum_xr_x\sigma_x\right)=\sum_xr_x(\sigma_x(e_1)-\sigma_x(e_0))=\sum_xr_xx-a_0\sum_xr_x.
\]
Entonces,
\[
c=\partial_1\left(\sum_xr_x\sigma_x\right)+\sum_xr_xa_0.
\]
Por lo tanto, $c\in B_0(X,A)$ y, así, $S_0(X)=B_0(X,A)$.
\item Tenemos que la inclusión $i_\alpha:X_\alpha\hookrightarrow X$ induce
$i_{\alpha\#}S_\bullet(X_\alpha)\hookrightarrow S_\bullet(X)$. Luego, se
tiene un morfismo
\[
\overline i_{\alpha\#}:S_\bullet(X_\alpha)/S_\bullet(A_\alpha)\to S_\bullet(X)/S_\bullet(A).
\]
Por la propiedad universal de la suma directa, se induce
\[
\oplus i_{\alpha*}:\bigoplus_\alpha S_\bullet(X_\alpha)/S_\bullet(A_\alpha)\to S_\bullet(X)/S_\bullet(A).
\]
Se tiene una biyección
\[
\bigsqcup_\alpha\{\sigma:\Delta_q\to X_\alpha\mid\im\sigma\not\subset A_\alpha\}\cong\{\sigma:\Delta_q\to X\mid\im\sigma\not\subset A\}.
\]
Así,
\[
H_q(X,A)=H_q(S_\bullet(X)/S_\bullet(A))
=H_q\left(\bigoplus_\alpha S_\bullet(X_\alpha)/S_\bullet(A_\alpha)\right)
\cong\bigoplus_\alpha H_q(S_\bullet(X_\alpha)/S_\bullet(A_\alpha))
=\bigoplus_\alpha H_q(X_\alpha,A_\alpha).
\]
\end{enumerate}

\item (\textbf{Homología y retractos}) Sea $X$ un espacio topológico y $A$
un subespacio de $X$ tal que $A$ es un retracto de $X$.
\begin{enumerate}[label=\alph*)]
\item Mostrar que $H_q(X)\cong H_q(A)\oplus H_q(X,A)$ para toda $q$.
\item Probar que existe un homomorfismo $\varphi:H_q(X,A)\to H_q(X)$ tal
que $j_*\circ\varphi=\mathrm{id}$, donde $j:X\to(X,A)$ es la inclusión.
\end{enumerate}

\noindent\textbf{Solución.}
\begin{enumerate}[label=\alph*)]
\item Dado $i:A\hookrightarrow X$ y $j:(X,\varnothing)\hookrightarrow(X,A)$,
y una retracción $r:X\to A$, entonces $ri=\mathrm{id}_A$. Por lo tanto,
$r_*i_*=\mathrm{id}_{H_q(A)}$ y $i_*$ es un monomorfismo de $H_q(A)$ en un
sumando directo de $H_q(X)$. El otro sumando es el núcleo de $r_*$. De la
sucesión exacta larga de la pareja $(X,A)$:
\[
\cdots\to H_q(A)\xrightarrow{i_*}H_q(X)\xrightarrow{j_*}H_q(X,A)\to\cdots
\]
Como $i_*$ es inyectivo, se sigue que
\[
0\to H_q(A)\xrightarrow{i_*}H_q(X)\xrightarrow{j_*}H_q(X,A)\to0
\]
Ya que $\ker j_*=\im i_*$, $j_*$ induce un isomorfismo de $\ker r_*$ en
$H_q(X,A)$.
\item Tomemos $\varphi:H_q(X,A)\to H_q(X)\cong H_q(A)\oplus H_q(X,A)$ como
la inclusión. Por lo anterior, se sigue el resultado.
\end{enumerate}

\item (\textbf{Escisión}) Sea $X$ un espacio y $A,B\subset X$ subespacios
tales que $X=A\cup B$. Probar que la pareja $\{A,B\}$ es escisiva si y sólo
si la inclusión $j:(A,A\cap B)\hookrightarrow(X,B)$ induce isomorfismos en
homología.

\noindent\textbf{Solución.} Consideremos el siguiente diagrama
\[
\begin{tikzcd}
\dfrac{S_*(A)}{S_*(A\cap B)} \arrow[rr,"\overline j"] \arrow[dr,swap,"\overline\imath"] & & \dfrac{S_*(X)}{S_*(B)} \\
& \dfrac{S_*(A)+S_*(B)}{S_*(B)} \arrow[ur,swap,"\overline\imath"] &
\end{tikzcd}
\]
También, consideremos el siguiente diagrama:
\[
\begin{tikzcd}
0 \arrow[r] & S_*(B) \arrow[r] \arrow[d,"="] & S_*(A)+S_*(B) \arrow[r] \arrow[d,"i"] & \dfrac{S_*(A)+S_*(B)}{S_*(B)} \arrow[r] \arrow[d,"\overline\imath"] & 0 \\
0 \arrow[r] & S_*(B) \arrow[r] & S_*(X) \arrow[r] & \dfrac{S_*(X)}{S_*(B)} \arrow[r] & 0
\end{tikzcd}
\]
Aplicando el teorema fundamental del álgebra homológica se obtiene que
\[
\begin{tikzcd}[column sep=small]
\cdots \arrow[r] & H_q(B) \arrow[r] \arrow[d,"="] & H_q(S_*(A)+S_*(B)) \arrow[r] \arrow[d,"i_*"] & H_q\left(\dfrac{S_*(A)+S_*(B)}{S_*(B)}\right) \arrow[r] \arrow[d,"\overline\imath_*"] & H_{q-1}(B) \arrow[r] \arrow[d,"="] & H_{q-1}(S_*(A)+S_*(B)) \arrow[r] \arrow[d] & \cdots \\
\cdots \arrow[r] & H_q(B) \arrow[r] & H_q(X) \arrow[r] & H_q(X,B) \arrow[r] & H_{q-1}(B) \arrow[r] & H_{q-1}(X) \arrow[r] & \cdots
\end{tikzcd}
\]
Supóngase que la pareja $\{A,B\}$ es escisiva; es decir
$i:S_*(A)+S_*(B)\to S_*(X)$ induce el isomorfismo $i_*$. Por el lema del $5$,
$\overline\imath_*$ es isomorfismo. Ya que $\overline\imath$ es isomorfismo
por Noether (y por lo tanto $\overline\imath_*$ es isomorfismo) se tiene
que $\overline\jmath_*$ es isomorfismo y, así, $j$ es escisión.

Recíprocamente, supongamos que $j$ es escisión. Entonces, $\overline\jmath_*$
es isomorfismo y, por lo tanto, $\overline\imath_*$ es isomorfismo. Por el
lema del $5$, $i_*$ es isomorfismo.

\item (\textbf{Sucesión exacta larga de una tripleta}) Si $B\subset A$, se
tiene la sucesión exacta de homología relativa
\[
\cdots\to H_{q+1}(X,A)\to H_q(A,B)\to H_q(X,B)\to H_q(X,A)\to H_{q-1}(A,B)\to\cdots
\]

\noindent\textbf{Solución.} Observemos que las parejas $\{X,A\}$ y
$\{B,A\}$ son escisivas pues $S_*(X)+S_*(A)=S_*(X)$ y
$S_*(B)+S_*(A)=S_*(A)$. La sucesión de Mayer--Vietoris
(Teorema~\ref{thm:mayer-vietoris}) en este caso es:
\[
\cdots\to H_q(X\cap A,A\cap B)\to H_q(X,B)\oplus H_q(A,A)\to H_q(X,A\cup B)\to H_{q-1}(X\cap A,A\cap B)\to\cdots
\]
Ya que $H_q(A,A)=0$, se tiene el resultado.

\item (\textbf{Homología del toro y de la botella de Klein con
Mayer--Vietoris}) Calcular:
\begin{enumerate}[label=\alph*)]
\item $\mathbb T^2=S^1\times S^1$.
\item $\RR P^2\#\RR P^2$.
\end{enumerate}

\noindent\textbf{Solución.}
\begin{enumerate}[label=\alph*)]
\item Consideremos dos cilindros $A$ y $B$. Entonces, $A,B\cong S^1$ y
$A\cap B\cong S^1\sqcup S^1$. Entonces,
\[
H_q(A)=H_q(B)\cong\begin{cases}\ZZ & q=0,1\\0 & \text{en otro caso}\end{cases}
\qquad
H_q(A\cap B)=\begin{cases}\ZZ\oplus\ZZ & q=0,1\\0 & \text{en otro caso.}\end{cases}
\]
Por lo tanto, en particular, $H_q(A)=H_q(B)=H_q(A\cap B)=0$ para $q>2$. Por
Mayer--Vietoris:
\[
\cdots\to H_q(A)\oplus H_q(B)\to H_q(\mathbb T^2)\to H_{q-1}(A\cap B)\to\cdots
\]
Por lo tanto, se sigue que $H_q(\mathbb T^2)=0$, $q>2$. Ya que $\mathbb T^2$
es conexo, $H_0(\mathbb T^2)=\ZZ$. Para $n=2$, la sucesión de
Mayer--Vietoris es
\[
\cdots\to H_2(A)\oplus H_2(B)\to H_2(\mathbb T^2)\xrightarrow{\delta}H_1(A\cap B)\xrightarrow{(i_*,j_*)}H_1(A)\oplus H_1(B)\to\cdots
\]
Entonces, se tiene:
\[
0\to H_2(\mathbb T^2)\xrightarrow{\delta}\ZZ\oplus\ZZ\xrightarrow{(i_*,j_*)}\ZZ\oplus\ZZ\to\cdots
\]
Ahora, sean $\alpha$ y $\beta$ las clases de homología que provienen de la
circunferencia ecuatorial de $A\cap B$. De hecho,
$(i_*,j_*):\ZZ[\alpha]\oplus\ZZ[\beta]\to\ZZ[\alpha]\oplus\ZZ[\beta]$.

Si $a:\Delta_1\to A\cap B$ y $b:\Delta_1\to A\cap B$ son los ciclos que
representan a las clases $\alpha$ y $\beta$,
\[
(i_*,j_*)(\alpha,0)=(i_*,j_*)[a]=([i\circ a],[j\circ a])=(\alpha,\beta)=(i_*,j_*)(0,\beta)
\]
ya que $\alpha=\beta$ en $H_1(A)$ y en $H_1(B)$. Por lo tanto, en términos
de estas bases, el morfismo $(i_*,j_*)$ puede representarse como:
\[
A=\begin{pmatrix}1&1\\1&1\end{pmatrix}:\ZZ\oplus\ZZ\to\ZZ\oplus\ZZ.
\]
Por lo tanto, $H_2(\mathbb T^2)=\ker A=\ZZ[\alpha-\beta]=\ZZ$.

Para $H_1(\mathbb T^2)$, concentrémonos en la parte de la sucesión de
Mayer--Vietoris:
\[
\cdots\to H_1(A\cap B)\xrightarrow{(i_*,j_*)}H_1(A)\oplus H_1(B)\xrightarrow{k_*-l_*}H_1(\mathbb T^2)\xrightarrow{\delta}H_0(A\cap B)\xrightarrow{(i_*,j_*)}H_0(A)\oplus H_0(B)\to\cdots
\]
en donde $k_*$ y $l_*$ son inducidos por las inclusiones $k:A\to\mathbb T^2$
y $l:B\to\mathbb T^2$. A continuación, cortamos la sucesión como
\[
0\to\ker\delta\to H_1(\mathbb T^2)\to\im\delta\to0
\]
en donde $\im\delta=\ker(i_*,j_*)\cong\ZZ$. Por otro lado,
\[
\ker\delta=\im(k_*-l_*)=\ZZ\oplus\ZZ/\ker(k_*-l_*)=\ZZ\oplus\ZZ/\ker(i_*,j_*)=\ZZ.
\]
Finalmente, $H_1(\mathbb T^2)=\ZZ\oplus\ZZ$. En resumen,
\[
H_q(\mathbb T^2)=\begin{cases}\ZZ & q=0,2\\ \ZZ\oplus\ZZ & q=1\\0 & \text{e.o.c.}\end{cases}
\]
\item Descomponemos $\RR P^2\#\RR P^2=A\cup B$ como la unión de dos bandas
de Möbius pegadas por sus círculos fronteras. Como la banda de Möbius se
retracta por deformación sobre su círculo central,
$H_2(A)\cong H_2(B)=H_2(S^1)=0$. Usando homología reducida, la parte
relevante de la sucesión de Mayer--Vietoris:
\[
0\to H_2(A\cup B)\to H_1(A\cap B)\xrightarrow{\phi}H_1(A)\oplus H_1(B)\to H_1(A\cup B)\to0
\]
Entonces,
\[
0\to H_2(\RR P^2\#\RR P^2)\to\ZZ\xrightarrow{\phi}\ZZ\oplus\ZZ\to H_1(\RR P^2\#\RR P^2)\to0
\]
Usando estas identificaciones, tenemos $\phi(n)=(2n,-2n)$. En particular,
$\phi$ es inyectiva y por lo tanto $H_2(\RR P^2\#\RR P^2)=0$. Por lo tanto,
\[
H_1(\RR P^2\#\RR P^2)\cong\ZZ\oplus\ZZ/\im\phi\cong\ZZ\oplus\ZZ/2\ZZ,
\]
ya que $\{(1,0),(1,-1)\}$ es una base para $\ZZ\oplus\ZZ$. También,
$H_0(\RR P^2\#\RR P^2)=\ZZ$ ya que $\RR P^2\#\RR P^2$ es conexo por arcos.
Por lo tanto,
\[
H_q(\RR P^2\#\RR P^2)=\begin{cases}\ZZ & q=0\\ \ZZ\oplus\ZZ/2\ZZ & q=1\\0 & q\ge2.\end{cases}
\]
\end{enumerate}

\item (\textbf{Homología de esferas con Mayer--Vietoris}) Se cumple que
\[
H_q(S^n)=\begin{cases}R\oplus R & q=n=0\\ R & 0=q<n\\ 0 & 0<q<n\\ R & 0<q=n\\ 0 & q>n.\end{cases}
\]

\noindent\textbf{Solución.} Si $n=0$, entonces $S^n=\{-1,+1\}$. Así,
\[
H_q(S^n)=H_q(\{-1\})\oplus H_q(\{1\})=\begin{cases}R\oplus R & q=0\\0 & \text{e.o.c.}\end{cases}
\]
Si $q=0$, ya que $S^n$ es conexo por arcos, $H_0(S^n)=R$.

Mostremos el resultado por inducción sobre $q$. Sean $U=S^n-\{N\}$ y
$V=S^n-\{S\}$. Entonces, $U$ y $V$ son abiertos contraíbles y por lo tanto
tienen el mismo tipo de homotopía que un punto. La inclusión,
$S^{n-1}\hookrightarrow U\cap V$ es una equivalencia de homotopía (debido a
que $U\cap V$ se retracta por deformación fuerte sobre $S^{n-1}$ a lo largo
de los grandes círculos que pasan por los polos).

Para $q=1$, la sucesión de Mayer--Vietoris aplicado a $S^n$ con la cubierta
abierta $\{U,V\}$ provee la sucesión exacta
\[
0=H_1(U)\oplus H_1(V)\to H_1(S^n)\to H_0(U\cap V)\xrightarrow{\phi}H_0(U)\oplus H_0(V)\cong R\oplus R
\]
Si $n>1$, $H_0(U\cap V)\cong H_0(S^{n-1})\cong R$ y la aplicación
$\phi:x\mapsto(x,-x)$. En particular, $\phi$ es inyectiva, y por
exactitud, $H_1(S^n)=0$.

Si $n=1$, entonces $H_0(U\cap V)\cong H_0(S^0)\cong R\oplus R$. La
aplicación $\phi:(x,y)\mapsto(x+y,-(x+y))$, ya que las dos componentes
conexas por arcos de $U\cap V$, que forman una base de $H_0(U\cap V)$, se
envían sobre la única componente conexa por arcos de $U$, que genera
$H_0(U)$, y de igual forma para $V$. En particular, $\ker\phi\cong R$, y
por exactitud, $H_1(S^1)\cong R$.

Para $q>1$, se tiene la sucesión de Mayer--Vietoris,
\[
0=H_q(U)\oplus H_q(V)\to H_q(S^n)\to H_{q-1}(U\cap V)\to H_{q-1}(U)\oplus H_{q-1}(V)=0
\]
Así, $H_q(S^n)\cong H_{q-1}(U\cap V)\cong H_{q-1}(S^{n-1})$. El resultado se
sigue por recurrencia.

\item (\textbf{Homología de esferas y aplicaciones})
\begin{enumerate}[label=\alph*)]
\item Sea $X$ un espacio topológico y sean $V\subset U\subset A\subset X$
subespacios tales que
\begin{enumerate}[label=\arabic*)]
\item $V$ puede escindirse de la pareja $(X,A)$
\item $j:(X-U,A-U)\hookrightarrow(X-V,A-V)$ es un retracto por deformación.
\end{enumerate}
Entonces, $U$ se puede escindir de $(X,A)$.
\item Sean $S^n_+$ y $S^n_-$ los hemisferios cerrados norte y sur de $S^n$,
respectivamente. Entonces, $j:(S^n_+,S^{n-1})\to(S^n,S^n_-)$ es una
escisión.
\item Probar que
\[
\tilde H_q(S^n)\cong\begin{cases}R & q=n\\0 & q\ne n.\end{cases}
\]
\item $S^{n-1}$ no es retracto de $D^n$ para $n\ge1$.
\item (\textbf{Teorema del punto fijo de Brouwer}) Toda función continua
$f:D^n\to D^n$ tiene un punto fijo.
\item Si $n\ne m$, entonces $S^n$ y $S^m$ no son del mismo tipo de
homotopía.
\item Si $n\ne m$, entonces $\RR^n$ y $\RR^m$ no son homeomorfos.
\item Sea $\rho:S^n\to S^n$ la reflexión $\rho(x_0,x_1,\dots,x_n)=(-x_0,x_1,\dots,x_n)$.
Entonces, $\rho_*:\tilde H_n(S^n)\to\tilde H_n(S^n)$ es multiplicar por
$-1$.
\item Sea $f:S^n\to S^n$ la restricción de una transformación ortogonal en
$\RR^{n+1}$. Entonces, $f_*:\tilde H_n(S^n)\to\tilde H_n(S^n)$ es
multiplicar por $\det f$.
\item Sea $a:S^n\to S^n$ la función antípoda $a(x)=-x$. Entonces,
$a_*:\tilde H_n(S^n)\to\tilde H_n(S^n)$ es multiplicar por $(-1)^{n+1}$.
\item $S^n$ tiene un campo vectorial tangente que no se anula si, y solo
si, $n$ es impar.
\end{enumerate}

\noindent\textbf{Solución.}
\begin{enumerate}[label=\alph*)]
\item Tenemos que existe $r:(X-V,A-V)\to(X-U,A-U)$ tal que $r\circ j=\mathrm{id}$
y $j\circ r\simeq\mathrm{id}$. Por invarianza homotópica, $j_*$ es
isomorfismo. Sean $\iota:(X-V,A-V)\hookrightarrow(X,A)$ e
$i:(X-U,A-U)\hookrightarrow(X,A)$ inclusiones. Por el primer punto,
$\iota_*$ es isomorfismo. Como $i=\iota\circ j$, entonces
$i_*=\iota_*\circ j_*$ y por lo tanto $i_*$ es isomorfismo.
\item Queremos escindir $U=$ hemisferio sur sin el ecuador. Sea
$V=\{x\in S^n\mid x_{n+1}<-\tfrac12\}$. Entonces, $V$ se puede escindir y
además $(S^n_+,S^{n-1})$ es retracto por deformación de
$(S^n-V,S^n_--V)$ haciendo la retracción a través de los círculos máximos
de $S^n$. Por el inciso anterior, $U$ se puede escindir.
\item Consideremos la sucesión de homología reducida de la pareja
$(D^n,S^{n-1})$:
\[
\to\tilde H_q(S^{n-1})\to\tilde H_q(D^n)\to H_q(D^n,S^{n-1})\xrightarrow{\cong}\tilde H_{q-1}(S^{n-1})\to\tilde H_{q-1}(D^n)\to
\]
Así, se tiene el isomorfismo
$\tilde H_{q-1}(S^{n-1})\cong H_q(D^n,S^{n-1})$.

Ahora, consideremos la sucesión de homología reducida para $(S^n,S^n_-)$:
\[
\cdots\to0=\tilde H_q(S^n_-)\to\tilde H_q(S^n)\xrightarrow{\cong}H_q(S^n,S^n_-)\xrightarrow{\hat\partial_q}\tilde H_{q-1}(S^n_-)=0\to\cdots
\]
Combinando, se tienen los isomorfismos:
\[
\tilde H_q(S^n)\xrightarrow{i_*}H_q(S^n,S^n_-)\xleftarrow{j_*}H_q(S^n_+,S^{n-1})\xrightarrow{\pi_*}H_q(D^n,S^{n-1})\xrightarrow{\hat\delta_q}\tilde H_{q-1}(S^{n-1})
\]
en donde $\pi_*$ es isomorfismo porque $\pi$ es homeomorfismo. Continuemos
por inducción sobre $n$. En primer lugar $S^0=\{a,b\}$, y por el Teorema de
escisión~\ref{thm:escision}, escindiendo a $U=\{b\}$,
\[
\tilde H_q(S^0)=H_q(S^0,\{b\})\cong H_q(\{a\})=\begin{cases}R & q=0\\0 & q\ne0.\end{cases}
\]
Supongamos que se tiene la homología de $S^{n-1}$. Usando el isomorfismo de
suspensión $\tilde H_q(S^n)\cong\tilde H_{q-1}(S^{n-1})$.
\item Sea $i:S^{n-1}\hookrightarrow D^n$. Supongamos que existe un retracto
$r:D^n\to S^{n-1}$. Se tiene
\[
\tilde H_{n-1}(S^{n-1})\xrightarrow{i_*}\tilde H_{n-1}(D^n)\xrightarrow{r_*}\tilde H_{n-1}(S^{n-1}).
\]
Ya que $r$ es retracto, se cumple que $(r\circ i)_*=\mathrm{id}$. Sin
embargo, $\tilde H_{n-1}(S^{n-1})\cong R$ y $\tilde H_{n-1}(D^n)=0$ y, por
lo tanto, $(r\circ i)_*=0$. Contradicción.
\item Supongamos que $f(x)\ne x$ para toda $x\in D^n$. Si $n=0$ no hay nada
que probar. Supongamos que $n>0$. Consideremos el semirrayo que va de
$f(x)$ a $x$ y sea $r(x)$ el punto en donde intersecta a $S^{n-1}$. Esto
define $r:D^n\to S^{n-1}$ que es una retracción pues los puntos de
$S^{n-1}$ no se mueven. Esto contradice el punto anterior.
\item Supongamos que existe una equivalencia homotópica $h:S^n\to S^m$.
Entonces $h_*:\tilde H_n(S^n)\overset{\cong}{\to}\tilde H_n(S^m)\cong0$.
Contradicción.
\item Supongamos que existe $h:\RR^n\to\RR^m$ un homeomorfismo. Componiendo
con una traslación, podemos suponer que $h(0)=0$. Sea $i:S^{n-1}\to\RR^n-\{0\}$.
Y que $i$ es un retracto por deformación, es una equivalencia homotópica.
Sea $\tilde h:\RR^n-\{0\}\to\RR^m-\{0\}$ y $\rho:\RR^m-\{0\}\to S^{m-1}$.
Entonces, se tiene una equivalencia homotópica $\rho\circ\tilde h\circ i:S^{n-1}\to S^{m-1}$.
Contradicción.
\item Se hará por inducción sobre $n$. Tomemos $S^0=\{a,b\}$. De la
sucesión,
\[
S_1(S^0)\xrightarrow{\delta_1}S_0(S^0)\xrightarrow{\epsilon}R
\]
se sigue que
\[
\tilde H_0(S^0)=\ker\epsilon/\im\delta_1=\{\alpha a+\beta b\mid\alpha+\beta=0,\ \alpha,\beta\in R\}.
\]
En este caso, $\rho:S^0\to S^0$ se da por $\rho(a)=b$ y $\rho(b)=a$. Así,
$\rho_*(\alpha a+\beta b)=\alpha\rho(a)+\beta\rho(b)=\alpha b+\beta a$. Sin
embargo, $\alpha=-\beta$ y por lo tanto $\alpha b+\beta a=-(\alpha a+\beta b)$.

El caso general se sigue del diagrama conmutativo
\[
\begin{tikzcd}
\tilde H_{n-1}(S^{n-1}) \arrow[r,"\rho_*"] \arrow[d,swap,"\sigma","\cong"'] & \tilde H_{n-1}(S^{n-1}) \arrow[d,"\cong","\sigma"'] \\
\tilde H_n(S^n) \arrow[r,swap,"\rho_*"] & \tilde H_n(S^n)
\end{tikzcd}
\]
en donde $\sigma=\hat\delta_q\circ\pi_*\circ j_*^{-1}\circ i_*$. Por
naturalidad de $\hat\delta_q$ se tiene el resultado.
\item Vamos a usar el hecho de que \emph{cualquier rotación de la esfera es
homotópica a la identidad}. Por ser $f$ ortogonal, $\det f=\pm1$. Si
$\det f=1$ entonces $f\in SO(n)$ y $f\simeq\mathrm{id}$, por lo que
$f_*=\mathrm{id}$. Si $\det f=-1$, componemos a $f$ con la reflexión
$\rho$. Entonces, $\rho\circ f$ es una rotación y por un resultado
$\rho\circ f\simeq\mathrm{id}$. Aplicando homología, $\rho_*\circ f_*=\mathrm{id}$.
Además, ya que $\rho$ y $f$ son homeomorfismos, $\rho_*$ y $f_*$ son
isomorfismos. Pero $\rho\circ\rho=\mathrm{id}$ implica que
$\rho_*\circ\rho_*=\mathrm{id}$ y así $\rho_*=f_*$; pero vimos que $\rho_*$
es multiplicar por $-1=\det f$.
\item Se sigue inmediatamente componiendo $n+1$ reflexiones.
\item Un campo vectorial tangente en $S^n$ puede verse como una función
continua $v:S^n\to\RR^{n+1}$ tal que $v\perp v(x)$. Si $n=2m+1$, definimos
\[
v(x_0,x_1,\dots,x_{2m+1})=(-x_1,x_0,-x_3,x_2,\dots,-x_{2m+1},x_{2m}).
\]
Supongamos ahora que existe un campo vectorial $v$ tal que $v(x)\ne0$ para
todo $x\in S^n$. Entonces, se define
\[
w:S^n\to S^n, \qquad x\mapsto\frac{v(x)}{\|v(x)\|}
\]
y $x\perp w(x)$. Sea
\[
F:S^n\times I\to S^n, \qquad (x,t)\mapsto x\cos t\pi+w(x)\sin t\pi.
\]
Entonces $F(x,0)=x$ y $F(x,1)=-x=a(x)$. Por lo tanto,
$\mathrm{id}\simeq a$ por el inciso anterior $n$ tiene que ser impar.
\end{enumerate}

\item (\textbf{Grado de una aplicación entre esferas}) Sea $f:S^n\to S^n$
una aplicación. Entonces,
\[
f_*:\tilde H_n(S^n;\ZZ)\to\tilde H_n(S^n;\ZZ), \qquad \alpha\mapsto d(f)\alpha,
\]
con $d(f)\in\ZZ$. Defínase $d(f)=\deg(f)$.
\begin{enumerate}[label=\alph*)]
\item $\deg(\mathbb 1)=1$.
\item $\deg(f)=0$ si $f$ no es sobreyectiva.
\item Si $f\simeq g$ entonces $\deg(f)=\deg(g)$. Hopf: si $\deg(f)=\deg(g)$
entonces $f\simeq g$.
\item $\deg(fg)=\deg(f)\deg(g)$. Como consecuencia, $\deg(f)=\pm1$ si $f$
es una equivalencia homotópica.
\item Si $f:S^n\to S^n$ es un homeomorfismo, $\deg(f)=\pm1$.
\item Si $f:S^n\to S^n$ no tiene puntos fijos, entonces
$\deg(f)=(-1)^{n+1}$.
\item Para toda aplicación continua, $f:S^{2n}\to S^{2n}$, sin puntos fijos,
existe $x_0\in S^n$ tal que $f(x_0)=-x_0$.
\item No existe una aplicación continua $f:S^{2n}\to S^{2n}$, tal que $x$
es ortogonal a $f(x)$, para todo $x\in S^{2n}$.
\item Sea $k$ un entero y $f:S^1\to S^1$, $z\mapsto z^k$. Entonces,
$\deg(f)=k$.
\item Para todo entero $k$, existe una aplicación continua $f:S^n\to S^n$
de grado $k$.
\end{enumerate}

\noindent\textbf{Solución.}
\begin{enumerate}[label=\alph*)]
\item Es claro ya que $\mathbb 1_*=\mathbb 1$.
\item Si $x_0\in S^n-f(S^n)$, entonces $f$ puede factorizarse a través de
la composición $S^n\to S^n-\{x_0\}\hookrightarrow S^n$ y
$H_q(S^n-x_0)=0$ ya que $S^n-\{x_0\}$ es contraíble. Por lo tanto, $f_*=0$.
\item Si $f\simeq g$, entonces $f_*=g_*$.
\item Se cumple que
$\deg(f\circ g)\alpha=(f\circ g)_*(\alpha)=f_*(\deg(g)\alpha)=\deg(f)\deg(g)\alpha$.
\item Se sigue directamente del punto anterior.
\item Si $f(x)\ne x$ entonces el segmento de línea de $f(x)$ a $-x$,
definida por $t\mapsto(1-t)f(x)-tx$, $0\le t\le1$, no pasa por el origen.
Por lo tanto, si $f$ no tiene puntos fijos, se tiene una homotopía
\[
f_t(x)=\frac{(1-t)f(x)-tx}{|(1-t)f(x)-tx|}
\]
entre $f$ y la aplicación antípoda $a$.
\item Si $f$ no tiene puntos fijos, entonces $f$ es homótopa a la
aplicación antípoda y $\deg(f)=-1$. Supongamos que $f(x)\ne-x$ para todo
$x\in S^{2n}$. Entonces, $a\circ f$ no tiene punto fijo, de modo que
$a\circ f$ es homótopa a la aplicación antípoda y por lo tanto
$\deg(f)=1$. Hemos así obtenido una contradicción.
\item Supongamos que tal aplicación $f$ existe. Entonces, $f$ no puede
tener punto fijo debido a que $f(x)=x$ implica que $f(x)\perp x$. Entonces,
existe $x_0$ tal que $f(x_0)=-x_0$. Pero esto último implica que
$f(x_0)\perp x_0=-1$ lo cual contradice el hecho de que $f(x)$ es ortogonal
a $x$ para todo $x\in S^{2n}$.
\item El caso $k=0$ es evidente. Si $|k|\ge1$, para $0\le j\le|k|-1$
entero, consideremos las cadenas singulares $c_j:I\to S^1$,
$t\mapsto e^{2i\pi(j+t)/|k|}$. El camino $c_0\cdots c_{|k|-1}$ es un
generador $g$ de $\pi_1(S^1,1)$ y $f_*(c_k)$ es un lazo que representa $g$
si $k>0$ y $g^{-1}$ si $k<0$. En homología, el generador $\alpha$
correspondiente a $g$ es entonces representado por el ciclo
$c_0+\cdots+c_{|k|-1}$ y, para todo $k$, $f_*(c_k)$ es un ciclo que
representa $\alpha$ si $k>0$, $-\alpha$ si $k<0$. Se deduce que
$f_*(\alpha)=k\alpha\in H_1(S^1)$.
\item Se demostrará por inducción sobre $n$. Para $n=1$ es el inciso
anterior. Supongamos que la aserción es verdadera para $n$ y sea
$f:S^n\to S^n$ con $\deg(f)=n$. La suspensión $\Sigma S^n$ es homeomorfa a
$S^{n+1}$ y la suspensión $\Sigma f$ de $f$ hace conmutar el diagrama
\[
\begin{tikzcd}
H_{n+1}(\Sigma S^n) \arrow[r,"(\Sigma f)_*"] & H_{n+1}(\Sigma S^n) \\
H_n(S^n) \arrow[u,"\cong"] \arrow[r,swap,"f_*"] & H_n(S^n) \arrow[u,swap,"\cong"]
\end{tikzcd}
\]
Esto muestra la aserción para $n+1$.
\end{enumerate}

\item (\textbf{Acción libre sobre $S^{2n}$}) Sea $G$ un grupo no trivial y
$G\times S^n\to S^n$ una acción continua y libre, con $n$ par. Pruebe que
$G=\ZZ/2$.

\noindent\textbf{Solución.} Recordemos que una acción de un grupo $G$
sobre un espacio $X$ es un homomorfismo $G\to\Homeo(X)$. La acción es libre
si el homeomorfismo correspondiente a cada elemento no trivial de $G$ no
tiene puntos fijos. Se tiene entonces un homomorfismo
\[
\eta:G\to\Homeo(S^n), \qquad g\mapsto(\eta(g):S^n\to S^n).
\]
El homeomorfismo $\eta(g)$ induce,
\[
(\eta(g))_*:\tilde H_\bullet(S^n;\ZZ)\to\tilde H_\bullet(S^n;\ZZ), \qquad \alpha\mapsto\deg(\eta(g))\alpha.
\]
Ya que $\eta(g)$ es homeomorfismo, por el ejercicio 8 se tiene que
$\deg(\eta(g))=\pm1$. Sea entonces,
\[
d:G\to\{\pm1\}, \qquad g\mapsto\deg(\eta(g)).
\]
Se tiene que $d$ es un homomorfismo de grupos por el ejercicio 8. Ya que la
acción es libre, $\eta(g):S^n\to S^n$ no tiene puntos fijos. Por lo tanto,
$d(g)=(-1)^{n+1}$. Ya que $n$ es par, $d(g)=-1$. De esta forma,
$\ker(d)=\{g\in G\mid d(g)=1\}=\{1\}$. Por lo tanto, $G\cong\ZZ/2\ZZ$.

\item (\textbf{Homología local de variedades})
\begin{enumerate}[label=\alph*)]
\item Sea $M$ una variedad topológica de dimensión $n$ y sea $x\in M$.
Probar que
\[
H_q(M,M-\{x\};R)\cong\begin{cases}R & q=n\\0 & q\ne n.\end{cases}
\]
\item Si $n\ne m$, entonces un abierto $U$ de $\RR^n$ no puede ser
homeomorfo a un abierto $V$ de $\RR^m$.
\end{enumerate}

\noindent\textbf{Solución.}
\begin{enumerate}[label=\alph*)]
\item Recordemos que $M$ es una variedad topológica de dimensión $n$ si
cada $x\in M$ tiene una vecindad abierta $U$ homeomorfa a una bola abierta
en $\RR^n$.

Ya que $\overline{M-U}=M-U\subset M-\{x\}=(M-\{x\})^\circ$, aplicando el
Teorema~\ref{thm:escision} obtenemos que, la inclusión
\[
j:(M-(M-U),M-\{x\}-(M-U))=(U,U-\{x\})\hookrightarrow(M,M-\{x\})
\]
induce el isomorfismo $H_q(U,U-\{x\})\cong H_q(M,M-\{x\})$.

Luego, consideremos la sucesión exacta de la pareja $(U,U-\{x\})$:
\[
\cdots\to\tilde H_q(U-\{x\})\to\tilde H_q(U)\to H_q(U,U-\{x\})\xrightarrow{\partial_q}\tilde H_{q-1}(U-\{x\})\to\cdots.
\]
Ya que $U$ es contraíble, $\tilde H_q(U)=0$ para toda $q\ge0$, dando así
$H_q(U,U-\{x\})\cong\tilde H_{q-1}(U-\{x\})$.

Sin embargo, $U-\{x\}$ es homotópicamente equivalente a $S^{n-1}$, de modo
que
\[
\tilde H_{q-1}(U-\{x\})\cong\tilde H_{q-1}(S^{n-1})=\begin{cases}R & q=n\\0 & q\ne n.\end{cases}
\]
\item Si $x$ es un punto de $U$, entonces $\{U,\RR^n-\{x\}\}$ es una
cubierta abierta de $\RR^n$. Ya que tanto $U$ como $\RR^n-\{x\}$ es
abierto, entonces la pareja $\{U,\RR^n-\{x\}\}$ es escisiva. Por el
ejercicio 3, se tiene que la inclusión
\[
j:(U,U\cap(\RR^n-\{x\}))=(U,U-\{x\})\hookrightarrow(\RR^n,\RR^n-\{x\})
\]
induce el isomorfismo $H_\bullet(U,U-\{x\})\cong H_\bullet(\RR^n,\RR^n-\{x\})$.
De igual forma, si $y\in V$, se tiene
$H_\bullet(V,V-\{y\})\cong H_\bullet(\RR^m,\RR^m-\{y\})$.

Supongamos que existe un homeomorfismo entre $U$ y $V$, entonces
\[
H_\bullet(U,U-\{x\})\cong H_\bullet(V,V-\{y\}).
\]
Por el inciso anterior, tiene que suceder que $n=m$.
\end{enumerate}

\item (\textbf{Invarianza del borde}) Si $n\ge1$, sea
$\mathcal H_n=[0,+\infty[\times\RR^{n-1}$ y
$\partial\mathcal H_n=\{0\}\times\RR^{n-1}\subset\RR^n$.
\begin{enumerate}[label=\alph*)]
\item Sea $x\in\mathcal H_n$. Demostrar que $x\in\partial\mathcal H_n$ si y
solo si $H_n(\mathcal H_n,\mathcal H_n-\{x\})=0$.
\item Sea $\phi:\mathcal H_n\to\mathcal H_n$ un homeomorfismo. La
restricción de $\phi$ a $\partial\mathcal H_n$ es un homeomorfismo.
\end{enumerate}

\noindent\textbf{Solución.}
\begin{enumerate}[label=\alph*)]
\item Si $x\in\partial\mathcal H_n$, entonces $\mathcal H_n-\{x\}$ es un
convexo y por lo tanto es contraíble. Así,
$\tilde H_n(\mathcal H_n-\{x\})=0$. La sucesión larga de la pareja
$(\mathcal H_n,\mathcal H_n-\{x\})$ es:
\[
\cdots\to\tilde H_n(\mathcal H_n-\{x\})\to\tilde H_n(\mathcal H_n)\to H_n(\mathcal H_n,\mathcal H_n-\{x\})\to\tilde H_{n-1}(\mathcal H_n-\{x\})\to\tilde H_{n-1}(\mathcal H_n)\to\cdots
\]
Por lo tanto, $H_n(\mathcal H_n,\mathcal H_n-\{x\})=0$.

Recíprocamente, sea $x\in\mathcal H_n-\partial\mathcal H_n$. Entonces,
existe un disco abierto $D$ centrado en $x$. Como $\mathcal H_n-D$ es un
cerrado contenido en el abierto $\mathcal H_n-\{x\}$, por el Teorema de
escisión~\ref{thm:escision}, se tiene que
$H_n(D,D-\{x\})\cong H_n(\mathcal H_n,\mathcal H_n-\{x\})$.

Finalmente, debido a que $h:D\to\RR^n$, $x\mapsto0$, es un homeomorfismo,
$H_n(D,D-\{x\})=\ZZ$ y de esta manera
$H_n(\mathcal H_n,\mathcal H_n-\{x\})=\ZZ$.
\item Si $x\in\mathcal H_n$, el homeomorfismo $\phi$ manda
$\mathcal H_n-\{x\}$ en $\mathcal H_n-\{\phi(x)\}$. Por lo tanto,
\[
H_n(\mathcal H_n,\mathcal H_n-\{x\})\cong H_n(\mathcal H_n,\mathcal H_n-\{\phi(x)\}).
\]
En particular, $\phi(\partial\mathcal H_n)=\partial\mathcal H_n$.
\end{enumerate}

\item (\textbf{Homología de un espacio de adjunción sobre un dip}) Sean $Y$
un espacio topológico, $f:S^{n-1}\to Y$ y $Z=Y\cup_fD^n$. Sea $R$ un
dominio de ideales principales. Probar que si $n>0$, entonces
\[
\tilde H_n(Z;R)\cong\tilde H_n(Y;R)\oplus\ker f_*.
\]

\noindent\textbf{Solución.} Por el Teorema~\ref{thm:tres-incisos} inciso
3., se tiene la sucesión exacta:
\[
0\to\tilde H_n(Y;R)\to\tilde H_n(Z;R)\to\ker f_*\to0,
\]
en donde $f_*:R=\tilde H_{n-1}(S^{n-1})\to\tilde H_{n-1}(Y)$. Ya que $R$ es
dominio de ideales principales, $\ker f_*$ está generado por un elemento no
cero que no es divisor de cero. Por lo tanto $\ker f_*$ es libre (y por lo
tanto proyectivo) y así, la sucesión exacta se escinde.

\item (\textbf{Homología de un CW sobre un anillo noetheriano}) Sea $X$ un
complejo CW de dimensión $n$ con un número finito de células y sea $R$ un
anillo noetheriano. Probar que $H_q(X;R)$ es finitamente generado para todo
$q$ y que $H_q(X;R)=0$ para $q>n$.

\noindent\textbf{Solución.} Sea $n\ge0$. Para $i\le n$, sea
$\Gamma^i=\{i\text{-células de }X\}$ y $\alpha_i<+\infty$. Supongamos que
$n=0$. Entonces,
\[
H_q(X)=\bigoplus_{x\in X}H_q(\{x\})=\begin{cases}R^{\alpha_0} & q=0\\0 & q>0.\end{cases}
\]
Por lo tanto, cuando $n=0$, se cumple que $H_q(X)=0$ para $q>0$ y
$H_q(X)=R^{\alpha_0}$. De este modo, el resultado se cumple para $n=0$.
Luego, supongamos que $n>0$ y que el resultado se cumple para complejos CW
de dimensión $n-1$. Sea $X$ un complejo CW de dimensión $n$; entonces
\[
X=X^{(n-1)}\cup_fD^n,
\]
en donde $f:S^{n-1}\to X^{(n-1)}$. Por el Teorema~\ref{thm:tres-incisos},
\begin{enumerate}[label=\arabic*)]
\item $\tilde H_q(X)\cong\tilde H_q(X^{(n-1)})$ para $q\ne n,n-1$
\item $0\to\im f_*\to\tilde H_{n-1}(X^{(n-1)})\to\tilde H_{n-1}(X)\to0$
\item $0\to\tilde H_n(X^{(n-1)})\to\tilde H_n(X)\to\ker f_*\to0$
\end{enumerate}
Supongamos que $q>n$. Entonces, ya que $q>n-1$ y por inducción,
$H_q(X)\cong H_q(X^{(n-1)})=0$.

Cuando $q<n-1$, se tiene por inducción que $H_q(X^{(n-1)})$ es finitamente
generado. Por lo tanto, $H_q(X)$ es finitamente generado para $q<n-1$.

El punto 2) afirma que $H_{n-1}(X)\cong H_{n-1}(X^{(n-1)})/\im f_*$ y por
ser cociente de un módulo finitamente generado es finitamente generado.

Ahora bien, ya que $R$ es noetheriano, $\ker f_*$ es finitamente generado.
Por lo tanto, debido a que $H_n(X^{(n-1)})$ es finitamente generado, por el
punto 3) se tiene que $H_n(X)$ es finitamente generado.

\item (\textbf{Generadores})
\begin{enumerate}[label=\alph*)]
\item Considérese la aplicación
\[
e:(I,\partial I)\to(S^1,*), \qquad t\mapsto e(t):=e^{2\pi it}.
\]
Entonces, $[e]\in H_1(S^1)$ es un generador.
\item Sea $e_0=t_0<t_1<\cdots<t_n=e_1$ una subdivisión de $\Delta_1$ y sea
la cadena
\[
c=\langle t_0,t_1\rangle+\cdots+\langle t_{n-1},t_n\rangle\in S_1(\Delta_1)
\]
en donde $\langle t_{i-1},t_i\rangle$ es la función afín de $\Delta_1$ en
$\Delta_1$ que manda $e_0$ en $t_{i-1}$ y $e_1$ en $t_i$. Sea
$\sigma:\Delta_1\to X$ un $1$-simplejo singular. Entonces
$\sigma-\sigma_\#(c)\in B_1(X)$ y por lo tanto $[\sigma-\sigma_\#(c)]=0$ en
$H_1(X)$.
\item Sea $\iota:\Delta_1\to\Delta_1$ un simplejo singular tal que
$\iota(e_0)=e_1$ y $\iota(e_1)=e_0$. Sea $\sigma:\Delta_1\to X$ continua.
Entonces $\sigma+\sigma_\#(\iota)\in B_1(X)$ y $[\sigma+\sigma_\#(\iota)]=0$
en $H_1(X)$.
\item Sea $z:\Delta_1\to X$ un $1$-ciclo. Sea $\overline z(t)=z(1-t)$.
Entonces, $[\overline z]=-[z]$ en $H_1(X)$.
\item Consideremos $n$ puntos distintos sobre $S^1$ y
$[\beta_1+\cdots+\beta_n]\in H_1(S^1)$, donde $\beta_i:I\to S^1$ es un
encaje tal que $\beta_i(0)=x_i$ y $\beta_i(1)=x_{i+1\ \mathrm{m\acute od}\ n}$.
Entonces $[\beta_1+\cdots+\beta_n]$ es un generador.
\end{enumerate}

\noindent\textbf{Solución.}
\begin{enumerate}[label=\alph*)]
\item Consideremos la sucesión de la pareja $(I,\partial I)$:
\[
\cdots\tilde H_1(\partial I)\to\tilde H_1(I)\to H_1(I,\partial I)\xrightarrow{\hat\partial_1}\tilde H_0(\partial I)\to\tilde H_0(I)
\]
Entonces, ya que $I$ es contraíble
$H_1(I,\partial I)\cong\tilde H_0(\partial I)=\ker\epsilon=\{\alpha a+\beta b\mid\alpha,\beta\in R,\ \alpha+\beta=0\}$.
Sea $[\mathrm{id}_{\Delta_1}]\in H_1(I,\partial I)$. Nótese que
$[\mathrm{id}_{\Delta_1}]$ es un $1$-ciclo pues
$\partial_1(\mathrm{id}_{\Delta_1})\in S_0(\partial I)$. Además,
$\hat\partial_1[\mathrm{id}_{\Delta_1}]=[\partial_1(\mathrm{id}_{\Delta_1})]=b-a$.
Por lo tanto $[\mathrm{id}_{\Delta_1}]\in H_1(I,\partial I)\cong R$ es un
generador. Consideremos el siguiente diagrama
\[
\begin{tikzcd}
(I,\partial I) \arrow[rr,"e"] \arrow[dr,swap,"q"] & & (S^1,*) \\
& (I/\partial I,*) \arrow[ur,swap,"\overline e"] &
\end{tikzcd}
\]
Entonces, se tiene que
\[
\begin{tikzcd}
H_1(I,\partial I) \arrow[rr,"e_*"] \arrow[dr,swap,"q_*"] & & H_1(S^1,*) \\
& H_1(I/\partial I,*) \arrow[ur,swap,"\overline e_*"] &
\end{tikzcd}
\]
Ya que $q_*$ es isomorfismo y $\overline e_*$ son isomorfismos, $e_*$ es
isomorfismo. Por lo tanto $e_*[\mathrm{id}_{\Delta_1}]=[e]\in H_1(S^1)$ es
un generador.
\item Consideremos $\mathrm{id}_{\Delta_1}:\Delta_1\to\Delta_1\in S_1(\Delta_1)$.
Obsérvese que $\partial_1(\mathrm{id}_{\Delta_1})=e_1-e_0$. También,
\[
\partial_1(c)=t_1-t_0+t_2-t_1+\cdots+t_n-t_{n-1}=e_1-e_0.
\]
Por lo tanto, $\partial_1(\mathrm{id}_{\Delta_1}-c)=0$ y
$[\mathrm{id}_{\Delta_1}-c]\in H_1(\Delta_1)=0$. De esta manera, existe una
$2$-cadena $d$ de $S_2(\Delta_1)$ tal que $\partial_2(d)=\mathrm{id}_{\Delta_1}-c$.
Así,
\[
\sigma-\sigma_\#(c)=\sigma_\#(\mathrm{id}_{\Delta_1})-\sigma_\#(c)=\sigma_\#(\mathrm{id}_{\Delta_1}-c)=\sigma_\#\partial_2(d)=\partial_2\sigma_\#(d).
\]
Por lo tanto, $\sigma-\sigma_\#(c)\in B_1(X)$ y $[\sigma-\sigma_\#(c)]=0\in H_1(X)$.
\item Se tiene que $\partial_1(\iota)=e_0-e_1$. Por tanto,
$\partial_1(\mathrm{id}_{\Delta_1})=-\partial_1(\iota)$ y
$\partial_1(\mathrm{id}_{\Delta_1}+\iota)=0$. Entonces,
$[\mathrm{id}_{\Delta_1}+\iota]\in H_1(\Delta_1)=0$. Así, existe
$c\in S_2(\Delta_1)$ tal que $\partial_2(c)=\mathrm{id}_{\Delta_1}+\iota$.
Se sigue que
\[
\sigma+\sigma_\#(\iota)=\sigma_\#(\mathrm{id}_{\Delta_1}+\iota)=\sigma_\#\partial_2(c)=\partial_2\sigma_\#(c)\in B_1(X)
\]
y $[\sigma+\sigma_\#(\iota)]=0$ en $H_1(X)$.
\item Sea $\overline c:\Delta_2\to X$ dado por $\overline c(t_0,t_1)=z(t_0)$.
Tenemos,
\[
\partial_2(\overline c)=\overline c\circ F_2^0-\overline c\circ F_2^1+\overline c\circ F_2^2
=\overline c(1-t,t)-\overline c(0,t)+\overline c(t,0)
=z(1-t)-z(0)+z(t).
\]
De esto último,
$\partial_2(\overline c)=\overline z-c_{z_0}+z$. Sea $\hat c:\Delta_2\to X$
la función constante con valor $z_0$. Por lo tanto, tomando $c=\overline c+\hat c$,
se sigue $\partial_2(c)=\overline z+z$.
\item Sea $X=S^1$, $B=\im\beta_1$ y $A=\im\beta_2\cup\cdots\cup\im\beta_n$.
Tenemos que $A\cap B=\{x_1,x_2\}$. Empleando la sucesión de
Mayer--Vietoris,
\[
0=\tilde H_1(A)\oplus\tilde H_1(B)\to H_1(S^1)\xrightarrow{\Delta_1}\tilde H_0(\{x_1,x_2\})\to0
\]
Tomemos $[\beta_1]\in H_1(B,A\cap B)$. Basta ver que
$[\beta_1+\cdots+\beta_n]=[\beta_1]\in H_1(S^1,A)$. Pero,
\[
\beta_1+\cdots+\beta_n-\beta_1=\beta_2+\cdots+\beta_n\in B_1(S^1,A)
\]
por lo que
\[
\Delta_1[\beta_1+\cdots+\beta_n]=\hat\partial_1[\beta_1]=\partial_1(\beta_1)=x_2-x_1
\]
que es un generador pues $\epsilon(x_2-x_1)=1-1=0$.
\end{enumerate}

\item (\textbf{Construcción CW})
\begin{enumerate}[label=\alph*)]
\item Sea $Z$ un espacio Hausdorff y $Y\subset Z$ un subespacio cerrado.
Supongamos que existe una función $F:(D^n,S^{n-1})\to(Z,Y)$ tal que
$F|_{D^n-S^{n-1}}:D^n-S^{n-1}\to Z-Y$ es homeomorfismo. Entonces,
\[
Z\cong Y\cup_{F|_{S^{n-1}}}D^n.
\]
\item $\RR P^n$ es un complejo CW con una celda de dimensión $0\le r\le n$.
\item $\CC P^n$ es un complejo CW con una celda de dimensión $0\le2r\le2n$.
\item Sea $A\subset X$ y supongamos que tenemos una identificación
$F:X\to X/\!\sim$ que solamente identifica puntos de $A$. Consideremos
$F|_A:A\to F(A)$ y el espacio de adjunción $F(A)\cup_{F|_A}X$. Supongamos
que, si $p:F(A)\sqcup X\to F(A)\cup_{F|_A}X$, $p|_X$ es identificación.
Entonces
\[
X/\!\sim\ \cong F(A)\cup_{F|_A}X.
\]
\end{enumerate}

\noindent\textbf{Solución.}
\begin{enumerate}[label=\alph*)]
\item Veamos que $Z$ es el pushout del siguiente diagrama
\[
\begin{tikzcd}
S^{n-1} \arrow[r,"F|_{S^{n-1}}"] \arrow[d,swap,"j"] & Y \arrow[d] \\
D^n \arrow[r,swap,"F"] & Z
\end{tikzcd}
\]
Ya que $Z$ es Hausdorff y $D^n$ es compacto, si $F$ es continua, entonces
el diagrama es pushout debido a que $F(D^n)\subset Z$ es cerrado.
\item Por inducción sobre $n$. Para $n=0$, $\RR P^0=\{*\}$ que es un
complejo CW con una única $0$-célula. Supongamos que $\RR P^{n-1}$ es un
complejo CW. Consideremos el diagrama
\[
\begin{tikzcd}
S^{n-1} \arrow[r,"F|_{S^{n-1}}"] \arrow[d,swap,"j"] & \RR P^{n-1} \arrow[d] \\
D^n \arrow[r,swap,"F"] & \RR P^n
\end{tikzcd}
\]
en donde estamos considerando a $\RR P^{n-1}\to\RR P^n$ como
$\RR P^{n-1}=\{[x_1,\dots,x_{n+1}]\mid x_{n+1}=0\}\subset\RR P^n$ y en
donde las flechas horizontales se obtienen como la composición
\[
(D^n,S^{n-1})\overset{\cong}{\to}(S^n_+,S^{n-1})\hookrightarrow(S^n,S^{n-1})\twoheadrightarrow(\RR P^n,\RR P^{n-1})
\]
y por lo tanto $F$ es de la forma
\[
F(x_1,\dots,x_n)=\left[x_1,\dots,x_n,\sqrt{1-\sum_{i=1}^nx_i^2}\right].
\]
Si consideramos $F|_{D^n-S^{n-1}}$, su inversa está dada por
\[
F':\RR P^n-\RR P^{n-1}\to D^n-S^{n-1}, \qquad
[x_1,\dots,x_{n+1}]\mapsto\frac{x_{n+1}}{\sqrt{\sum_{i=1}^n(x_ix_{n+1})^2}}(x_1,\dots,x_n).
\]
Se sigue que $\RR P^n\cong\RR P^{n-1}\cup_{F|_{S^{n-1}}}D^n$.
\item Por inducción sobre $n$. Para $n=0$, $\CC P^n=\{*\}$ es un complejo
CW con una sola $0$-célula. Consideremos el diagrama
\[
\begin{tikzcd}
S^{2n-1} \arrow[r,"F|_{S^{2n-1}}"] \arrow[d,swap,"j"] & \CC P^{n-1} \arrow[d] \\
D^{2n} \arrow[r,swap,"F"] & \CC P^n
\end{tikzcd}
\]
Defínase,
\[
F:(D^{2n},S^{2n-1})\to(\CC P^n,\CC P^{n-1}), \qquad
(x_1,\dots,x_{2n})\mapsto\left[x_1+ix_2,\dots,x_{2n-1}+ix_{2n},\sqrt{1-\sum_{i=1}^{2n}x_i^2}\right].
\]
Definimos la inversa de $F$ como
\[
F':\CC P^{n-1}-\CC P^n\to D^{2n}-S^{2n-1},
\]
\[
[z_1,\dots,z_{n+1}]\mapsto\frac{1}{a_{n+1}}\Big(\Re(z_1\overline{z_{n+1}}),\Im(z_1\overline{z_{n+1}}),\dots,\Re(z_n\overline{z_{n+1}}),\Im(z_n\overline{z_{n+1}})\Big)
\]
donde $a_{n+1}=|z_{n+1}|\sqrt{\sum_{i=1}^{n+1}|z_i|^2}$. Así,
$\CC P^n\cong\CC P^{n-1}\cup_{F|_{S^{2n-1}}}D^{2n}$.
\item Consideremos el diagrama
\[
\begin{tikzcd}
X \arrow[r] \arrow[d,swap,"F"] & F(A)\cup_{F|_A}X \\
X/\!\sim \arrow[ur] &
\end{tikzcd}
\]
Notemos que $p|_X$ es sobreyectivo ya que $p$ identifica $a\sim F(a)$ y
$a\in X$. Supongamos que, dados $a,a'\in X$, $p|_X(a)=p|_X(a')$; entonces
$a=a'$ y así $F(a)=F(a')$.

Supongamos que $F(a)=F(a')$. Como $pF(a)=p(a)$ y $pF(a')=p(a')$, entonces
$p(a)=p(a')$. Por lo tanto, las identificaciones $p|_X$ y $F$ son
compatibles (identifican el mismo subespacio).
\end{enumerate}

\item (\textbf{Misma homología distinto tipo de homotopía})
\begin{enumerate}[label=\alph*)]
\item Probar que $\tilde H_q\left(\bigvee_nS^1;R\right)=\bigoplus_n\tilde H_q(S^1;R)$.
\item Calcular $H_q(S^1\times S^1)$ para $q\ge0$.
\item Calcular $\pi_1(S^1\times S^1)$.
\item Calcular $H_q(S^1\vee S^1\vee S^2)$ para $q\ge0$.
\item Calcular $\pi_1(S^1\vee S^1\vee S^2)$.
\item Concluir que el toro $S^1\times S^1$ tiene los mismos grupos de
homología que el bouquet $S^1\vee S^1\vee S^2$ y que, sin embargo, no son
homotópicamente equivalentes.
\end{enumerate}

\noindent\textbf{Solución.}
\begin{enumerate}[label=\alph*)]
\item Se procede por inducción sobre $n$. Para $n=1$ se tiene trivialmente.
Supongamos que el resultado es válido para $n-1$. Sea $X=\bigvee_nS^1=A\cup B$,
con $A=\bigvee_{n-1}S^1$ y $B=S^1$. Para usar el Teorema de
Mayer--Vietoris~\ref{thm:mayer-vietoris}, es necesario establecer que
$\{A,B\}$ es escisiva. Recordemos del ejercicio 3 que lo anterior es
equivalente a mostrar que la inclusión $j:(A,A\cap B)\hookrightarrow(X,B)$
induce el isomorfismo $j_*$. Vamos a usar el primer inciso del ejercicio 7
como sigue. Sea $U=B-A$. Observe que $\overline U\not\subset\mathring B$.
Sea $V\subset U$ tal que $\overline V\subset\mathring B$. Entonces,
\[
(X-U,B-U)\hookrightarrow(X-V,B-V)
\]
es retracto por deformación. Entonces, $U$ se puede escindir de la pareja
$(X,B)$; esto es, $(X-U,B-U)\hookrightarrow(X,B)$ induce isomorfismo en
homología. Sin embargo, la inclusión anterior no es más que $j$. Por lo
tanto, $\{A,B\}$ es escisiva.

Se tiene la sucesión de Mayer--Vietoris:
\[
\cdots\to H_q(\{*\},\{*\})\to H_q(A,\{*\})\oplus H_q(S_1,\{*\})\to H_q(X,\{*\})\to H_{q-1}(\{*\},\{*\})\to\cdots
\]
Por lo tanto, se tiene el isomorfismo
$H_q\left(\bigvee_{n-1}S^1,\{*\}\right)\oplus H_q(S^1,\{*\})\cong H_q\left(\bigvee_nS^1,\{*\}\right)$.
De esta forma,
\[
\tilde H_q\left(\bigvee_nS^1\right)=\tilde H_q\left(\bigvee_{n-1}S^1\right)\oplus\tilde H_q(S_1)
=\bigoplus_{n-1}\tilde H_q(S^1)\oplus\tilde H_q(S^1)
=\bigoplus_n\tilde H_q(S^1).
\]
\item Homología del toro $\mathbb T^2=S^1\times S^1$.
\begin{itemize}
\item \textbf{Tres incisos.} Describamos a $S^1\times S^1$ como un complejo
CW
\[
\begin{tikzcd}
\partial(I\times I) \arrow[r,"F|_{\partial(I\times I)}"] \arrow[d,swap,"j"] & S^1\vee S^1 \arrow[d] \\
I\times I \arrow[r,swap,"F"] & S^1\times S^1
\end{tikzcd}
\]
en donde $F(s,t)=(e^{2\pi is},e^{2\pi it})$. Claramente,
$F|_{I\times I-\partial(I\times I)}$ es un homeomorfismo entre
$I\times I-\partial(I\times I)$ y $S^1\times S^1-S^1\vee S^1$. Así,
\[
S^1\times S^1\cong(S^1\vee S^1)\cup_{F|_{\partial(I\times I)}}I^2.
\]
Aplicando el Teorema~\ref{thm:tres-incisos} tenemos lo siguiente:
\begin{enumerate}[label=\arabic*)]
\item $\tilde H_q(S^1\times S^1)\cong\tilde H_q(S^1\vee S^1)=0$ para
$q\ne1,2$
\item $\tilde H_1(S^1\times S^1)\cong\tilde H_1(S^1\vee S^1)/\im f_*\cong\ZZ\oplus\ZZ/\im f_*$
en donde $f_*:\tilde H_1(S^1)\to\tilde H_1(S^1\vee S^1)$. Sean
\[
i_1:S^1\to S^1\vee S^1, \qquad z\mapsto(z,1), \qquad\qquad
i_2:S^1\to S^1\vee S^1, \qquad z\mapsto(1,z).
\]
Los generadores de $H_1(S^1\vee S^1)$ son $i_{1*}[e]$ y $i_{2*}$. Denotemos
$i_{1*}[e]=[\alpha]$ y $i_{2*}=[\beta]$ con $\alpha(t)=(e^{2\pi it},1)$ y
$\beta(t)=(1,e^{2\pi it})$. Por otra parte, por el ejercicio 14, tenemos
que el generador para $H_1(S^1)$ es $[a+b+c+d]\in H_1(S^1)$ donde
\[
a(t)=(t,0), \qquad b(t)=(1,t), \qquad c(t)=(1-t,1), \qquad d(t)=(0,1-t).
\]
Así las cosas,
\[
f_*[a+b+c+d]=[f\circ a]+[f\circ b]+[f\circ c]+[f\circ d]
=[\alpha]+[\beta]+[\overline\alpha]+[\overline\beta]
=[\alpha]+[\beta]-[\alpha]-[\beta]
=0,
\]
en donde usamos el ejercicio 14. Por lo tanto,
$H_1(S^1\times S^1)\cong\ZZ\oplus\ZZ$.
\item $0\to H_2(S^1\vee S^1)=0\to H_2(S^1\times S^1)\xrightarrow{\cong}\ker f_*\cong\ZZ\to0$,
esto es, $H_2(S^1\times S^1)\cong\ZZ$.
\end{enumerate}
En resumen,
\[
H_q(\mathbb T^2)=\begin{cases}\ZZ & q=0,2\\ \ZZ\oplus\ZZ & q=1\\0 & \text{e.o.c.}\end{cases}
\]
\item \textbf{Celular.} La descomposición celular de $\mathbb T^2$ consiste
de $\Gamma^0=\{e_1^0\}$ que se forma a través de los cuatro vértices del
rectángulo que se identifican a un punto, $\Gamma^1=\{e_1^1,e_2^1\}$ que se
forman identificando los segmentos paralelos, $\Gamma^2=\{e_1^2\}$.

Así, $\mathbb T^2=\bigcup_{i=0}^2\Gamma^i$. Por el
Teorema~\ref{thm:cw-celular}, se tiene que $H_q(\mathbb T^2)=0$ para $q>2$.
Por el mismo teorema, se tiene la sucesión exacta:
\[
0\to\ZZ\xrightarrow{d_2}\ZZ\oplus\ZZ\xrightarrow{d_1}\ZZ\to0
\]
Ya que $\#\Gamma^0=1$, entonces, por la Observación~\ref{obs:cw-observaciones}
inciso 3., se tiene que $d_1=0$. Ahora, por medio del
Teorema~\ref{thm:grado-pegado}, calculemos $d_2$.
\[
d_2(e_1^2)=e_1^1+e_2^1-e_1^1-e_2^1=0.
\]
Por lo tanto, $d_1$ y $d_2$ son los mapeos cero. Finalmente,
\[
H_0(\mathbb T^2)=\mathcal C_0/\im d_1=\ZZ; \qquad
H_1(\mathbb T^2)=\ker d_1/\im d_2\cong\mathcal C_1/\im d_2=\ZZ\oplus\ZZ; \qquad
H_2(\mathbb T^2)=\ker d_2/\im d_3=\ZZ.
\]
En resumen,
\[
H_q(\mathbb T^2)=\begin{cases}\ZZ & q=0,2\\ \ZZ\oplus\ZZ & q=1\\0 & \text{e.o.c.}\end{cases}
\]
\end{itemize}
\item Se tiene que
$\pi_1(\mathbb T^2,*)=\pi_1(S^1,*)\oplus\pi_1(S^1,*)=\ZZ\oplus\ZZ$.
\item Homología de $X=S^1\vee S^1\vee S^2$.
\begin{itemize}
\item Por el Corolario~\ref{cor:homologia-bouquet},
\[
\tilde H_q(S^1\vee S^1\vee S^2)=\tilde H_q(S^1)\oplus\tilde H_q(S^1)\oplus\tilde H_q(S^2)
=\begin{cases}\ZZ & q=2\\ \ZZ\oplus\ZZ & q=1\\0 & \text{e.o.c.}\end{cases}
\]
Usando la definición $H_0(X)=\tilde H_0(X)\oplus\ZZ$ se sigue.
\item \textbf{Celular.} El espacio $S^1\vee S^1\vee S^2$ tiene la
descomposición celular: $\Gamma^0=\{e^0\}$, $\Gamma^1=\{e_1^1,e_2^1\}$,
$\Gamma^2=\{e_1^2\}$. Así, $X=\bigcup_{i=0}^2\Gamma^i$. El mismo cálculo
hecho para el toro funciona en este caso.
\end{itemize}
\item Grupo fundamental de $X=S^1\vee S^1\vee S^2$. Por el Teorema de van
Kampen, $\pi_1(X,*)\cong\ZZ*\ZZ$.
\item Dos espacios con grupos fundamentales no isomorfos no pueden tener el
mismo tipo de homotopía.
\end{enumerate}

\item (\textbf{Números de Betti y característica de Euler}) Definimos al
\emph{número de Betti} de una pareja $(X,A)$ como
\[
\beta(X)=\sum_{q\ge0}(-1)^q\ran\big(H_q(X,A;\ZZ)\big).
\]
Sea $X$ un complejo CW con un número finito de células. Sea
$\alpha_q(X)=\#\Gamma^q(X)$. Se define la \emph{característica de Euler} de
$X$ como
\[
\chi(X)=\sum_{q\ge0}(-1)^q\alpha_q(X).
\]
\begin{enumerate}[label=\alph*)]
\item Consideremos $(X,A)$ una pareja tal que $\beta(X)$, $\beta(A)$ y
$\beta(X,A)$ están todos definidos. Entonces, $\beta(X,A)=\beta(X)-\beta(A)$.
\item Mostrar que $\beta(S^n)=1+(-1)^n$.
\item Sea $Y$ un espacio topológico tal que $\beta(Y)$ está definida. Sea
$f:S^{n-1}\to Y$ y $Z=Y\cup_fD^n$. Probar que $\beta(Z)=\beta(Y)+(-1)^n$.
\item Si $X$ es un complejo CW con un número finito de celdas. Entonces
$\chi(X)=\beta(X)$.
\item Sea $X$ un complejo CW con un número finito de celdas. Entonces
$\chi(X)$ es invariante topológico.
\item Sea $X$ un espacio triangulable. Entonces, $\chi(X)$ no depende de la
triangulación.
\item Calcular $\chi(\RR P^n)$.
\item Calcular $\chi(\CC P^n)$.
\item Calcular $\chi(T_g)$.
\item Calcular $\chi(P_g)$.
\end{enumerate}

\noindent\textbf{Solución.}
\begin{enumerate}[label=\alph*)]
\item Para este inciso, usaremos el siguiente resultado general:
consideremos la sucesión exacta de grupos abelianos finitamente generados
\[
0\to A_1\to A_2\to\cdots\to A_{r-1}\to A_r\to0
\]
Entonces $\sum_{i=1}^r(-1)^{i+1}\ran(A_i)=0$.

En nuestro caso, considérese la sucesión
\[
0\to\cdots\to H_q(A;\ZZ)\to H_q(X;\ZZ)\to H_q(X,A;\ZZ)\to H_{q-1}(A;\ZZ)\to\cdots\to H_0(X,A;\ZZ)\to0
\]
Por lo anterior, $\pm\beta(X)\pm\beta(A)\pm\beta(X,A)=0$. De aquí, se
sigue que $\beta(X,A)=\beta(X)-\beta(A)$.
\item Recordemos que,
\[
H_q(S^n;R)=\begin{cases}R\oplus R & q=n=0\\ R & 0=q<n\\ 0 & 0<q<n\\ R & 0<q=n\\ 0 & q>n.\end{cases}
\]
Entonces,
\[
\beta(S^n)=\sum_{q\ge0}(-1)^q\ran(H_q(S^n;\ZZ))
=\ran(H_0(S^n;\ZZ))+\sum_{q=1}^{n-1}(-1)^q\ran(H_q(S^n;\ZZ))+(-1)^n\ran(H_n(S^n;\ZZ))
=\begin{cases}1+(-1)^n & n>0\\2 & n=0\end{cases}
=1+(-1)^n.
\]
\item Sabemos que $F:(D^n,S^{n-1})\to(Z,Y)$ induce
$F_*:H_q(D^n,S^{n-1};\ZZ)\overset{\cong}{\to}H_q(Z,Y;\ZZ)\ \forall q$. Por
el inciso anterior,
\[
\beta(D^n,S^{n-1})=\beta(D^n)-\beta(S^n)=1-\left(1+(-1)^{n-1}\right)=(-1)^n.
\]
Por lo tanto,
\[
\beta(Z)-\beta(Y)=\beta(Z,Y)=\beta(D^n,S^{n-1})=(-1)^n.
\]
\item Por inducción sobre el número de celdas de $X$. Si $X$ tiene una
celda, entonces $X=\{*\}$. Así, $\beta(X)=1=\chi(X)$. Supongamos el
resultado cierto para todo CW con $m$ celdas. Sea $X$ un CW con $(m+1)$
celdas. Entonces, $X\cong Y\cup_fD^n$ en donde $Y$ es CW con $m$ celdas.
Por el inciso anterior $\beta(X)=\beta(Y)+(-1)^n$. Por la hipótesis de
inducción, $\beta(Y)=\chi(Y)$. Entonces,
\[
\chi(X)=\chi(Y)+(-1)^n=\beta(Y)+(-1)^n=\beta(X).
\]
\item Se sigue por el inciso anterior ya que $\beta(X)$ es invariante
topológico.
\item Sean $K_1$ y $K_2$ dos triangulaciones de $X$. Entonces,
$|K_1|\cong X\cong|K_2|$. Por lo tanto,
$\chi(K_1)=\beta(|K_1|)=\beta(|K_2|)\cong\chi(K_2)$.
\item Tenemos que $\RR P^n=\RR P^{n-1}\cup_fD^n$ en donde
$f:S^{n-1}\to\RR P^{n-1}$. Ya que $\RR P^n$ es complejo CW finito,
\[
\chi(\RR P^n)=\beta(\RR P^n)=\beta(\RR P^{n-1})+(-1)^n
=\beta(\RR P^{n-2})+(-1)^{n-1}+(-1)^n
=\beta(\RR P^{n-3})+(-1)^{n-2}+(-1)^{n-1}+(-1)^n
=\cdots
=\beta(\RR P^{n-(n-1)})+(-1)^{n-(n-2)}+\cdots+(-1)^{n-1}+(-1)^n
=\beta(S^1)+\sum_{j=2}^n(-1)^j
=\sum_{j=2}^n(-1)^j
=1+\sum_{j=3}^{n-1}(-1)^j+(-1)^n.
\]
Obsérvese que si $n$ es par, entonces $n-1$ es impar. Por lo tanto, la suma
en la ecuación anterior contiene $n-3$ términos. Esto es, un número impar
de términos cuyo primer término es $-1$. Así, cuando $n$ es par,
$\chi(\RR P^n)=1-1+(-1)^n=1$.

Por otra parte, si $n$ es impar, entonces $n-1$ es par. De este modo, la
suma en la ecuación anterior contiene un número par de términos. Por lo
tanto, la suma es cero. Cuando $n$ es impar,
$\chi(\RR P^n)=1+0+(-1)^n=0$. En resumen,
\[
\chi(\RR P^n)=\begin{cases}1 & n\text{ par}\\0 & \text{e.o.c.}\end{cases}
\]
\item Tenemos que $\CC P^n=\CC P^{n-1}\cup_fD^{2n}$ en donde
$f:S^{2n-1}\to\CC P^{n-1}$. Ya que $\CC P^n$ es un complejo CW finito,
\[
\chi(\CC P^n)=\beta(\CC P^n)=\beta(\CC P^{n-1})+(-1)^{2n}
=\beta(\CC P^{n-2})+(-1)^{2(n-1)}+(-1)^{2n}
=\cdots
=\beta(\CC P^{n-(n-1)})+(-1)^{2(n-(n-2))}+\cdots+(-1)^{2n}
=\beta(S^2)+\sum_{j=2}^n(-1)^{2j}
=2+(n-1)
=n+1.
\]
\item Tenemos que $T_g=\bigvee_{2g}S^1\cup_fD^2$. Ya que $T_g$ es un
complejo CW finito, se tiene que
\[
\chi(T_g)=\beta(T_g)=\beta\left(\bigvee_{2g}S^1\right)+(-1)^2=1-2g+1=2-2g.
\]
\item Tenemos que $P_g=\bigvee_gS^1\cup_fD^2$. Ya que $P_g$ es un complejo
CW finito, se tiene que
\[
\chi(P_g)=\beta(P_g)=\beta\left(\bigvee_gS^1\right)+(-1)^2=1-g+1=2-g.
\]
\end{enumerate}

\item (\textbf{Homología y característica de Euler de espacios}) Calcular
\begin{enumerate}[label=\alph*)]
\item $H_q(\RR P^n)$ y $\chi(\RR P^n)$.
\item $H_q(\CC P^n)$ y $\chi(\CC P^n)$.
\item $H_q(T_g)$ y $\chi(T_g)$.
\item $H_q(P_g)$ y $\chi(P_g)$.
\end{enumerate}

\noindent\textbf{Solución.}
\begin{enumerate}[label=\alph*)]
\item \textbf{Cálculo de $H_q(\RR P^2)$.} Se tiene que
$\RR P^2\cong S^1\cup_fD^2$, en donde $f:S^1\to S^1$ está dada por
$z\mapsto z^2$. Usando el Teorema~\ref{thm:tres-incisos},
\begin{enumerate}[label=\arabic*)]
\item $\tilde H_q(\RR P^2)\cong\tilde H_q(S^1)=\begin{cases}R & q=1\\0 & q\ne1.\end{cases}$
\item $H_1(\RR P^2)\cong H_1(S^1)/\im f_*$, en donde
$f_*:H_1(S^1)\to H_1(S^1)$. Determinemos $\im f_*$. Sea $[e]\in H_1(S^1)$
un generador. Consideremos $c=\langle t_0,t_1\rangle+\langle t_1,t_2\rangle=\sigma_1+\sigma_2$
donde $\sigma_1(t)=t/2$ y $\sigma_2(t)=(t+1)/2$. Por el ejercicio 14,
$e-e_\#(c)\in B_1(S^1)$. Así, $[e]-[e_\#(c)]=0$ en $H_1(S^1)$ y
$[e]=[e_\#(c)]$. Por lo tanto,
\[
f_*[e]=f_*[e_\#(c)]=[f_\#e_\#(c)]=[f_\#e_\#(\sigma_1+\sigma_2)]=[f e\sigma_1+f e\sigma_2],
\]
donde
\[
fe\sigma_1(t)=fe(t/2)=f\exp(2\pi it/2)=e(t)
\]
y también
\[
fe\sigma_2(t)=fe((t+1)/2)=f\exp(2\pi i(t+1)/2)=e(t).
\]
De aquí, resulta que $f_*[e]=[e]+[e]=2[e]$. Por lo tanto, podemos concluir
que $\im f_*=2R$. Así,
\[
H_1(\RR P^2)\cong H_1(S^1)/2R\cong R/2R.
\]
\item $0\to H_2(S^1)=0\to H_2(\RR P^2)\xrightarrow{\cong}\ker f_*=R_2\to0$
En resumen,
\[
H_q(\RR P^2)=\begin{cases}R & q=0\\ R/2R & q=1\\ R_2 & q=2\\0 & \text{e.o.c.}\end{cases}
\]
\end{enumerate}

\textbf{Tres incisos.} Por inducción sobre $n$. Para $n=0$, $\RR P^0=\{*\}$.
Entonces,
\[
H_q(\RR P^0)\cong\begin{cases}R & q=0\\0 & q\ne0.\end{cases}
\]
Supongamos que tenemos la homología para $\RR P^{n-1}$. Se tiene que
$\RR P^n\cong\RR P^{n-1}\cup_fD^n$. Por el Teorema~\ref{thm:tres-incisos}:
\begin{enumerate}[label=\arabic*)]
\item $\tilde H_q(\RR P^n)\cong\tilde H_q(\RR P^{n-1})$, $q\ne n,n-1$.
\item $\tilde H_{n-1}(\RR P^n)\cong\tilde H_{n-1}(\RR P^{n-1})/\im f_*$, en
donde $f_*:\tilde H_{n-1}(S^{n-1})\to\tilde H_{n-1}(\RR P^{n-1})$.

\textbf{Hecho.} $f_*$ es $0$ si $n-1$ es par y es multiplicar por $2$ si
$n-1$ es impar. Por lo tanto,
\[
\tilde H_{n-1}(\RR P^n)\cong\begin{cases}R_2 & n-1\text{ par}\\ R/2R & n-1\text{ impar.}\end{cases}
\]
\item $0\to\tilde H_n(\RR P^{n-1})=0\to\tilde H_n(\RR P^n)\xrightarrow{\cong}\ker f_*\to0$
Por lo tanto,
\[
\tilde H_n(\RR P^n)\cong\ker f_*\cong\begin{cases}R & n\text{ impar}\\ R_2 & n\text{ par.}\end{cases}
\]
\end{enumerate}
En resumen,
\[
H_q(\RR P^n)=\begin{cases}R & q=0\text{ o }q=n\text{ impar}\\ R/2R & q\text{ impar con }1\le q\le n-1\\ R_2 & q\text{ par con }1\le q\le n\\0 & \text{e.o.c.}\end{cases}
\]
\begin{itemize}
\item \textbf{Celular.} Se tiene que $\RR P^n=\bigcup_{k=0}^n\Gamma^k$ en
donde $\Gamma^k=\{e^k\}$ para $0\le k\le n$; y el mapeo pegado para $e^k$
es la proyección $\varphi:S^{k-1}\to\RR P^{k-1}$. Ya que $H_q(\RR P^n)=0$
para $q>n$, y $\#\Gamma^k=1$, se tiene la sucesión exacta
\[
0\to R\xrightarrow{d_n}R\xrightarrow{d_{n-1}}R\xrightarrow{d_{n-2}}\cdots\xrightarrow{d_2}R\xrightarrow{d_1}R\to0
\]
Obsérvese que, por la Observación~\ref{obs:cw-observaciones}, $d_1=0$.
Para calcular cada $d_k$ necesitamos calcular el grado de la composición
\[
S^{k-1}\xrightarrow{\varphi}\RR P^{k-1}\xrightarrow{q}\RR P^{k-1}/\RR P^{k-2}=S^{k-2}.
\]
Pasando al cociente, se induce una aplicación continua
$r:S^{k-1}\vee S^{k-1}=S^{k-1}/S^{k-2}\to\RR P^{k-1}/\RR P^{k-2}=S^{k-2}$.
La aplicación $r$ es la aplicación que vale la identidad sobre la primer
esfera y vale la antípoda sobre la segunda. Por lo tanto, el grado de la
composición $q\varphi$ es entonces $1+(-1)^k$.

De esta manera,
\[
0\to R\xrightarrow{2}R\xrightarrow{0}R\xrightarrow{2}\cdots\xrightarrow{0}R\xrightarrow{2}R\xrightarrow{0}R\to0; \qquad
0\to R\xrightarrow{0}R\xrightarrow{2}R\xrightarrow{0}R\xrightarrow{2}\cdots\xrightarrow{0}R\xrightarrow{2}R\to0
\]
Obsérvese que,
\[
H_0(\RR P^n)=\ker d_0/\im d_1=R; \qquad
H_1(\RR P^n)=\ker d_1/\im d_2=R/2R; \qquad
H_2(\RR P^n)=\ker d_2/\im d_3=R_2; \qquad
H_3(\RR P^n)=\ker d_3/\im d_4=R/2R;
\]
\[
\vdots
\]
\[
H_{2j}(\RR P^{2j})=\ker d_{2j}/\im d_{2j+1}=R_2; \qquad
H_{2j+1}(\RR P^{2j+1})=\ker d_{2j+1}/\im d_{2j+2}=R.
\]
Por lo tanto,
\[
H_q(\RR P^n)=\begin{cases}R & q=0\text{ o }q=n\text{ impar}\\ R/2R & q\text{ impar con }1\le q\le n-1\\ R_2 & q\text{ par con }1\le q\le n\\0 & \text{e.o.c.}\end{cases}
\]
\item \textbf{Cálculo de $\chi(\RR P^n)$.} Nótese que $\ZZ_2=0$.
\[
\chi(\RR P^n)=\sum_{q\ge0}(-1)^q\ran\big(H_q(\RR P^n;\ZZ)\big)
=\ran(H_0(\RR P^n;\ZZ))+\sum_{j=0}^{[\frac{n-1}2]-1}(-1)^{2j+1}\ran\big(H_{2j+1}(\RR P^n;\ZZ)\big)+(-1)\ran\big(H_{2j+1}(\RR P^n;\ZZ)\big)
=\begin{cases}1 & n\text{ par}\\0 & \text{si no.}\end{cases}
\]
\end{itemize}
\item \textbf{Tres incisos~\ref{thm:tres-incisos}.} Por inducción sobre
$n$. Para $n=1$:
\[
H_q(\CC P^n)=H_q(\{*\})=\begin{cases}R & q=0\\0 & \text{e.o.c.}\end{cases}
\]
Supongamos que se tiene la homología de $\CC P^{n-1}$. Se tiene que
$\CC P^n\cong\CC P^{n-1}\cup_fD^{2n}$. Por el
Teorema~\ref{thm:tres-incisos}:
\begin{enumerate}[label=\arabic*)]
\item $\tilde H_q(\CC P^n)\cong\tilde H_q(\CC P^{n-1})$, $q\ne2n,2n-1$.
\item $\tilde H_{2n-1}(\CC P^n)\cong\tilde H_{2n-1}(\CC P^{n-1})/\im f_*$,
en donde $f_*:\tilde H_{2n-1}(S^{2n-1})\to\tilde H_{2n-1}(\CC P^{n-1})=0$
\item $0\to\tilde H_{2n}(\CC P^{n-1})=0\to\tilde H_{2n}(\CC P^n)\xrightarrow{\cong}\ker f_*\cong\tilde H_{2n-1}(S^{2n-1})\cong R$
\end{enumerate}
\begin{itemize}
\item \textbf{Celular.} Se tiene que $\CC P^n=\bigcup_{k=0}^n\Gamma^{2k}$
en donde $\Gamma^{2k}=\{e^{2k}\}$ para $0\le2k\le2n$. Por el
Teorema~\ref{thm:cw-celular}, $H_q(\CC P^n)=0$ para $q>2n$. También, ya que
las células de $\CC P^n$ no están en dimensiones adyacentes, la
Observación~\ref{obs:cw-observaciones} implica que
\[
H_q(\CC P^n)=\begin{cases}R & q=0,\dots,2n\\0 & \text{e.o.c.}\end{cases}
\]
\item \textbf{Cálculo de $\chi(\CC P^n)$.}
\[
\chi(\CC P^n)=\sum_{q\ge0}(-1)^q\ran\big(H_q(\CC P^n;\ZZ)\big)
=\sum_{j=0}^n(-1)^{2j}\ran\big(H_{2j}(\CC P^n;\ZZ)\big)
=n+1.
\]
\end{itemize}
\item \textbf{Tres incisos.} Tenemos que
$T_g=\bigvee_{2g}S^1\cup_{F|_A}D^2$. Por el Teorema~\ref{thm:tres-incisos}
\begin{enumerate}[label=\arabic*)]
\item $\tilde H_q(T_g)\cong\tilde H_q\left(\bigvee_{2g}S^1\right)\cong0$.
\item $\tilde H_1(T_g)\cong\tilde H_1\left(\bigvee_{2g}S^1\right)/\im f_*$,
donde $f_*:H_1(S^1)\to\tilde H_1\left(\bigvee_{2g}S^1\right)$. Determinemos
$\im f_*$. Sea $[e]\in H_1(S^1)$ un generador. Tomemos la cadena
\[
c=\langle t_0,t_1\rangle+\cdots+\langle t_{4g-1},t_{4g}\rangle=\sigma_1+\cdots+\sigma_{4g}.
\]
Por el ejercicio 14, $[e]=[e_\#(c)]\in H_1(S^1)$. Entonces,
\[
f_*[e]=f_*[e_\#(c)]=f_*[e_\#(\sigma_1+\cdots+\sigma_{4g})]=[fe\sigma_1]+\cdots+[fe\sigma_{4g}].
\]
Cada $[fe\sigma_i]$ es un generador de $H_1\left(\bigvee_{2g}S^1\right)$.
Por el ejercicio 14c, el $1$-simplejo singular recorrido en dirección
contraria es el negativo en homología. Por lo tanto $f_*[e]=0$. Así,
\[
\tilde H_1(T_g)\cong\tilde H_1\left(\bigvee_{2g}S^1\right)\Big/\im f_*\cong R^{2g}/0=R^{2g}.
\]
\item $0\to H_2\left(\bigvee_{2g}S^1\right)=0\to H_2(T_g)\xrightarrow{\cong}\ker f_*=R\to0$
\end{enumerate}
En resumen,
\[
H_q(T_g)=\begin{cases}R & q=0,2\\ R^{2g} & q=1\\0 & \text{e.o.c}\end{cases}
\]
\begin{itemize}
\item \textbf{Celular.} La descomposición celular de $T_g$ consiste de:
$\Gamma^0=\{e^0\}$, $\Gamma^1=\{e^1_{a_i}\cup e^1_{b_i}\}_{i=1}^g$,
$\Gamma^2=\{e^2\}$. Así, $T_g=\bigcup_{k=0}^2\Gamma^k$. Por lo tanto, por
el Teorema~\ref{thm:cw-celular}, se tiene que $H_q(T_g)=0$ para $q>2$. De
esta manera, se tiene la sucesión exacta:
\[
0\to R\xrightarrow{d_2}R^{2g}\xrightarrow{d_1}R\to0
\]
Por la Observación~\ref{obs:cw-observaciones} 3., ya que $\#\Gamma^0=1$, se
tiene que $d_1=0$. Calculemos $d_2$ mediante la fórmula del
Teorema~\ref{thm:grado-pegado}. Se tiene que
\[
d_2(e^2)=\sum_\beta d_{\alpha\beta}e^1_\beta
=e^1_{a_1}+e^1_{b_1}+\cdots+e^1_{a_g}+e^1_{b_g}-e^1_{a_1}-e^1_{b_1}-\cdots-e^1_{a_g}-e^1_{b_g}=0.
\]
Por lo tanto, $d_1=d_2=0$.
\[
H_0(T_g)=\ker d_0/\im d_1=R/0=R; \qquad
H_1(T_g)=\ker d_1/\im d_2=R^{2g}; \qquad
H_2(T_g)=\ker d_2/\im d_3=R/0=R.
\]
Finalmente,
\[
H_q(T_g)=\begin{cases}R & q=0,2\\ R^{2g} & q=1\\0 & \text{e.o.c.}\end{cases}
\]
\item \textbf{Cálculo de $\chi(T_g)$.}
\[
\chi(T_g)=\sum_{q\ge0}(-1)^q\ran\big(H_q(T_g;\ZZ)\big)=1-2g+1=2-2g.
\]
\end{itemize}
\item \textbf{Tres incisos.} Tenemos que $P_g=\bigvee_gS^1\cup_{F|_A}D^2$.
Por el Teorema~\ref{thm:tres-incisos},
\begin{enumerate}[label=\arabic*)]
\item $\tilde H_q(P_g)\cong\tilde H_q\left(\bigvee_gS^1\right)\cong\bigoplus_g\tilde H_q(S^1)=0$ si $q\ne1,2$
\item $\tilde H_1(P_g)\cong\tilde H_1\left(\bigvee_gS^1\right)/\im f_*$, en
donde $f_*:H_1(S^1)\to H_1\left(\bigvee_gS^1\right)$. Sea $[e]\in H_1(S^1)$
un generador. Tomemos,
\[
c=\langle t_0,t_1\rangle+\cdots+\langle t_{2g-1},t_{2g}\rangle=\sigma_1+\cdots+\sigma_{2g}.
\]
Entonces,
\[
f_*[e]=f_*[e_\#(c)]=[fe\sigma_1]+\cdots[fe\sigma_{2g}]=2\big([fe\sigma_1]+\cdots+[fe\sigma_g]\big).
\]
Así, $H_1(P_g)\cong H_1\left(\bigvee_gS^1\right)\Big/\im f_*\cong R^g/2R$.

Sea
\[
\phi:R^g\twoheadrightarrow R^{g-1}\oplus R/2R, \qquad
(r_1,\dots,r_g)\mapsto(r_1-r_g,r_2-r_g\dots,r_{g-1}-r_g,r_g).
\]
Entonces, $\ker\phi=\im f_*$. De esta manera, se sigue el isomorfismo:
\[
R^g/2R\cong R^{g-1}\oplus R/2R.
\]
\item $0\to H_2\left(\bigvee_gS^1\right)=0\to H_2(P_g)\xrightarrow{\cong}\ker f_*=R_2\to0$
\end{enumerate}
En resumen,
\[
H_q(P_g)=\begin{cases}R & q=0\\ R^{g-1}\oplus R/2R & q=1\\ R_2 & q=2\\0 & \text{e.o.c.}\end{cases}
\]
\begin{itemize}
\item \textbf{Celular.} La descomposición celular de $P_g$ consiste de:
$\Gamma^0=\{e^0\}$, $\Gamma^1=\{e^1_{a_i}\}_{i=1}^g$, $\Gamma^2=\{e^2\}$.
Así, $P_g=\bigcup_{k=0}^2\Gamma^k$. Por lo tanto, por el
Teorema~\ref{thm:cw-celular}, $H_q(P_g)=0$ para $q>2$. Se tiene la
sucesión exacta:
\[
0\to R\xrightarrow{d_2}R^g\xrightarrow{d_1}R\to0
\]
Ya que $\#\Gamma^0=1$, la Observación~\ref{obs:cw-observaciones} implica
que $d_1=0$. La fórmula del Teorema~\ref{thm:grado-pegado} da
\[
d_2(e^2)=2\big(e_{a_1}+\cdots+e_{a_g}\big).
\]
Entonces,
\[
H_0(P_g)=\ker d_0/\im d_1=R; \qquad H_1(P_g)=\ker d_1/\im d_2=R^g/2R; \qquad H_2(P_g)=\ker d_2/\im d_3=R_2.
\]
En resumen,
\[
H_q(P_g)=\begin{cases}R & q=0\\ R^{g-1}\oplus R/2R & q=1\\ R_2 & q=2\\0 & \text{e.o.c.}\end{cases}
\]
\item \textbf{Cálculo de $\chi(P_g)$.}
\[
\chi(P_g)=\sum_{q\ge0}(-1)^q\ran\big(H_q(P_g;\ZZ)\big)=1-(g-1)=2-g.
\]
\end{itemize}
\end{enumerate}

\item (\textbf{Moore}) Sean $n\ge2$ y $x$ un punto fijo de $S^{n-1}$. Sea
$f:S^{n-1}\to S^{n-1}$ una aplicación continua. Sea $X=S^{n-1}\cup_fD^n$.
\begin{enumerate}[label=\alph*)]
\item Existe $d>0$ tal que $\deg(f)=d$.
\item Calcular $H_\bullet(X;\ZZ)$.
\item Si $n>2$, entonces $X$ no es del mismo tipo de homotopía de
$\RR P^n$. Si $d=2$ y $n>1$ es impar, entonces $X$ no es del mismo tipo de
homotopía de $\RR P^n$.
\item Mostrar que existe un espacio topológico simplemente conexo $X$ tal
que $H_q(X;\ZZ)\cong\ZZ/d\ZZ$ si $q=n$ y $H_q(X;\ZZ)=0$ si $q\ne0,n$.
\item Mostrar que para todo grupo abeliano de tipo finito $G$, existe un
espacio topológico simplemente conexo $X$ tal que
\[
H_q(X;\ZZ)=\begin{cases}G & q=n\\0 & \text{en otro caso.}\end{cases}
\]
\end{enumerate}

\noindent\textbf{Solución.}
\begin{enumerate}[label=\alph*)]
\item Sea $c_d:S^n\to\bigvee_dS^n$ la aplicación cociente que se obtiene
por colapsar el complemento de $d$ bolas abiertas disjuntas en $S^n$ a un
punto. Se tiene una aplicación canónica $e_d:\bigvee_dS^n\to S^n$ que
identifica todos los sumandos a una sola esfera. Consideremos la
composición $f:e_d\circ c_d:S^n\to S^n$. Tomando la estructura CW de $S^n$
con una sola célula de dimensión $n$ y una sola célula de dimensión $0$,
obtenemos una estructura CW sobre $\bigvee_dS^n$ con $d$ células de
dimensión $n$ y una sola célula de dimensión $0$. Entonces,
\[
e_{k*}:\ZZ^d\to\ZZ, \qquad (x_1,\dots,x_d)\mapsto x_1+\cdots+x_d.
\]
Ahora, sea $\phi_j$ la proyección sobre la $j$-ésima esfera del bouquet
$\bigvee_dS^n$. Entonces, por construcción, $\phi_j\circ c_d$ es homótopa a
la identidad y por lo tanto $(\phi_j\circ c_d)_*$ es isomorfismo.
Deducimos entonces que la composición:
\[
H_n(S^n)\to H_n\left(\bigvee_dS^n\right)\cong\bigoplus_dH_n(S^n)\to H_n(S^n)
\]
es la identidad. Se sigue entonces que
$c_{d*}:H_n(S^n;\ZZ)\to H_n\left(\bigvee_dS^n;\ZZ\right)$ es la aplicación
diagonal $x\mapsto(x,\cdots,x)$. Así, $f_*(x)=x+\cdots+x=dx$.
\item \textbf{Celular.} El espacio $X$ es un complejo CW finito que admite
la descomposición celular: $\Gamma^0=\{e^0\}$, $\Gamma^{n-1}=\{e^{n-1}\}$,
$\Gamma^n=\{e^n\}$. Así, $X=\Gamma^0\cup\Gamma^{n-1}\cup\Gamma^n$. Por lo
tanto, por el Teorema~\ref{thm:cw-celular}, se tiene la sucesión exacta:
\[
0\to\ZZ\xrightarrow{d_n}\ZZ\xrightarrow{d_{n-1}}0\to\cdots\to0\to\ZZ\to0
\]
Se tiene que $d_{n-1}=0$ y $d_n=\times d$. Así las cosas,
\[
H_0(X;\ZZ)=\ZZ; \qquad H_n(X;\ZZ)=\ker d_{n-1}/\im d_n=\ZZ/d\ZZ.
\]
Por lo tanto,
\[
H_q(X;\ZZ)=\begin{cases}\ZZ & q=0\\ \ZZ/d\ZZ & q=n-1\\0 & \text{en otro caso.}\end{cases}
\]
\textbf{Tres incisos.} Tenemos que
\begin{enumerate}[label=\arabic*)]
\item $\tilde H_q(X;\ZZ)\cong\tilde H_q(S^{n-1};\ZZ)\cong0$ para
$q\ne n,n-1$.
\item $0\to\im f_*\cong d\ZZ\to\tilde H_{n-1}(S^{n-1};\ZZ)\cong\ZZ\to\tilde H_{n-1}(X;\ZZ)\to0$.
\item $0\to\tilde H_n(S^{n-1};\ZZ)=0\to\tilde H_n(X;\ZZ)\xrightarrow{\cong}\ker f_*=0$
\end{enumerate}
En resumen,
\[
H_q(X;\ZZ)=\begin{cases}\ZZ & q=0\\ \ZZ/d\ZZ & q=n-1\\0 & \text{en otro caso.}\end{cases}
\]
\item Sabemos que
\[
H_q(\RR P^n;\ZZ)=\begin{cases}\ZZ & q=0\text{ o }q=n\text{ impar}\\ \ZZ/2\ZZ & q\text{ impar },1\le q\le n-1\\0 & \text{en otro caso.}\end{cases}
\]
Entonces,
\[
H_1(\RR P^n;\ZZ)=\begin{cases}\ZZ & n=1\\ \ZZ/2\ZZ & n\ge2\\0 & \text{en otro caso.}\end{cases}
\]
Por otra parte,
\[
H_1(X;\ZZ)=\begin{cases}\ZZ/d\ZZ & n=2\\0 & \text{en otro caso.}\end{cases}
\]
Por lo tanto, $H_1(\RR P^n;\ZZ)=H_1(X;\ZZ)$ solamente cuando $d=2=n$; en
cualquier otro caso $H_1(\RR P^n)\ne H_1(X;\ZZ)$. De esta manera, bajo las
hipótesis del ejercicio, $X$ no puede ser del tipo de homotopía de
$\RR P^n$.
\item Tomemos $X=S^n\cup_fD^{n+1}$. Por lo anterior, $X$ cumple los
requerimientos.
\item Por el teorema de clasificación de los grupos abelianos de tipo
finito, existen enteros $m,r\in\NN$ y $d_1,\dots,d_r\in\NN-\{0,1\}$ tales
que $G\cong\ZZ^m\times\ZZ/d_1\ZZ\times\cdots\times\ZZ/d_r\ZZ$. Tomemos
$X_i=S^n$ para $i=1,\dots,m$ y $X_k=S^n\cup_fD^{n+1}$ para $k=1,\dots,r$ y
$\deg(f)=d_k$. Sea $X=\bigvee_{i=1}^mX_i\vee\bigvee_{k=1}^rX_k$. Entonces,
\[
H_q(X;\ZZ)=\bigoplus_{i=1}^mH_q(X_i;\ZZ)\oplus\bigoplus_{k=1}^rH_q(X_k;\ZZ)
=\bigoplus_{i=1}^m\ZZ\oplus\bigoplus_{k=1}^r\ZZ/d_k\ZZ
=G.
\]
\end{enumerate}

\item (\textbf{Homología de un complementario y curva de Jordan})
\begin{enumerate}[label=\alph*)]
\item Consideremos $\RR^{n+1}$ y $\RR^n=\RR^n\times\{0\}\subset\RR^{n+1}$.
Sea $A\subsetneq\RR^n$ cerrado. Entonces
\[
\tilde H_q(\RR^{n+1}-A)\cong\tilde H_{q-1}(\RR^n-A).
\]
\item Sean $A\subset\RR^p$ y $B\subset\RR^q$ subespacios cerrados y
homeomorfos. Consideremos $A\times\{0\}\subset\RR^p\times\RR^q$ y
$\{0\}\times B\subset\RR^p\times\RR^q$. Entonces,
\[
\RR^p\times\RR^q-(A\times\{0\})\cong\RR^p\times\RR^q-(\{0\}\times B).
\]
\item Sea $A\subsetneq\RR^n$ cerrado. Entonces,
\[
\tilde H_{q+i}(\RR^{n+i}-A)\cong\tilde H_q(\RR^n-A), \qquad i>0.
\]
\item Sean $A,B\subset\RR^n$ cerrados homeomorfos. Entonces,
$H_q(\RR^n-A)\cong H_q(\RR^n-B)$.
\item (\textbf{Jordan--Brouwer}) Sea $B\subset\RR^n$, con $n>1$, un
subespacio homeomorfo a $S^{n-1}$. Entonces, $\RR^n-B$ tiene exactamente
dos componentes.
\item Sea $n\ge2$ y $f:D^n\to\RR^n$ continua e inyectiva. Entonces, $f$ es
un encaje que mapea $\mathring D^n$ en la componente acotada de
$\RR^n-f(S^{n-1})$.
\item (\textbf{Invarianza del dominio}) Sea $n\ge2$. Sea $U\subset\RR^n$ un
abierto y $f:U\to\RR^n$ inyectivo. Entonces, $f$ es un encaje y $f(U)$ es
abierto. En particular, si $U\subset\RR^n$ es abierto y $f$ es un
homeomorfismo, entonces $f(U)\subset\RR^n$ es abierto.
\item Un nudo es un encaje $f:S^1\hookrightarrow\RR^3$. Entonces,
$\pi_1(\RR^3-f(S^1),*)^{\mathrm{ab}}\cong\ZZ$.
\end{enumerate}

\noindent\textbf{Solución.}
\begin{enumerate}[label=\alph*)]
\item Sea $Z=\RR^{n+1}-A$. Sean
\[
Z_+=\{(x,r)\in\RR^n\times\RR\mid r>0\text{ o }x\in\RR^n-A\}, \qquad
Z_-=\{(x,r)\in\RR^n\times\RR\mid r<0\text{ o }x\in\RR^n-A\}.
\]
Se tiene que $Z_+$ y $Z_-$ son abiertos en $Z$ y, además, $Z=Z_+\cup Z_-$ y
$Z_+\cap Z_-=(\RR^n-A)\times\RR$. Sea $\mathcal H_1=\{(x,1)\in\RR^{n+1}\}$.
Defínase,
\[
\rho:Z_+\to\mathcal H_1, \qquad (x,r)\mapsto(x,1).
\]
Entonces, $\rho\circ i=\mathrm{id}$. También, la aplicación continua
$H:\RR^{n+1}\times I\to\RR^{n+1}$ dada por $H(x,r,t)=(x,(1-t)r+t)$ es una
homotopía $H:i\circ\rho\simeq\mathrm{id}_{Z_+}$. Por lo tanto $Z_+$ es
contraíble. De la misma manera, $Z_-$ es contraíble.

Se tiene la sucesión de Mayer--Vietoris:
\[
\cdots\to\tilde H_{q+1}(Z)\to\tilde H_q(Z_+\cap Z_-)\to\tilde H_q(Z_+)\oplus\tilde H_q(Z_-)\to\tilde H_q(Z)\xrightarrow{\cong}\tilde H_{q-1}(Z_+\cap Z_-)\to0
\]
Por lo tanto, se sigue el isomorfismo
\[
\tilde H_q(\RR^{n+1}-A)\cong\tilde H_{q-1}\big((\RR^n-A)\times\RR\big)\cong\tilde H_{q-1}(\RR^n-A).
\]
\item Sea $\varphi:A\to B$ un homeomorfismo. Por el teorema de extensión
de Tietze, existe $\Phi:\RR^p\to\RR^q$ extensión de $\varphi$. Sea
$\psi:B\to A$ inverso de $\varphi$. Por Tietze, existe $\Psi:\RR^q\to\RR^p$
extensión de $\psi$. Defínase $L,R:\RR^p\times\RR^q\to\RR^p\times\RR^q$
como
\[
L(x,y)=(x,y-\Phi(x)) \qquad\text{y}\qquad R(x,y)=(x-\Psi(y),y).
\]
$L$ es homeomorfismo con inverso $\overline L(u,v)=(u,\Phi(u)+v)$ y $R$ es
homeomorfismo con inverso $\overline R(u,v)=(u+\Psi(v),v)$. Considérese el
conjunto
\[
\Gamma=\{(x,y)\in\RR^p\times\RR^q\mid x\in A,\ y=\varphi(x)\}=\{(x,y)\in\RR^p\times\RR^q\mid y\in B,\ x=\Psi(y)\}.
\]
Entonces, $L(\Gamma)=A\times\{0\}$ ya que $L(x,y)=(x,\varphi(x))=(x,0)$
pues $\Phi$ es extensión de $\varphi$. De la misma manera
$R(\Gamma)=\{0\}\times B$.
\item Por inducción sobre $i$. Notemos que $A\ne\RR^n$ ya que si
$A=\RR^n$ y $q=0$,
\[
\tilde H_{-1}(\varnothing)=0\not\cong\tilde H_0(\RR^{n+1}-\RR^n)\cong R.
\]
\item Supongamos que $A\ne\RR^n$ y $B\ne\RR^n$. Se tiene, por los incisos
anteriores,
\[
\tilde H_q(\RR^n-A)=\tilde H_{q+n}(\RR^{n+n}-A)=\tilde H_{q+n}(\RR^{n+n}-B)=\tilde H_q(\RR^n-B).
\]
Por lo tanto,
\[
H_q(\RR^n-A)=\tilde H_q(\RR^n-A)\oplus H_q(\{*\})\cong\tilde H_q(\RR^n-B)\oplus H_q(\{*\})=H_q(\RR^n-B).
\]
Supongamos ahora que $A=\RR^n$. Usando el argumento anterior para el caso
$A,B\subsetneq\RR^{n+1}$, se obtiene
$\tilde H_q(\RR^{n+1}-A)\cong\tilde H_q(\RR^{n+1}-B)$ para toda $q$. En
particular, para $q=0$,
$R\cong\tilde H_0(\RR^{n+1}-\RR^n)\cong\tilde H_0(\RR^{n+1}-B)$.

Por lo tanto, $\RR^{n+1}-B$ tiene dos componentes. Si $B\subsetneq\RR^n$,
entonces $\RR^{n+1}-B$ tendría una sola componente, lo cual es absurdo.
Por lo tanto, $B=\RR^n$. Así, $A=\RR^n$ y $B=\RR^n$ y el resultado es
cierto.
\item Sea $A=\{x\in\RR^n\mid\|x\|=1\}$. Entonces, $\RR^n-A$ tiene
exactamente dos componentes, es decir
$H_0(\RR^n-A)\cong R\oplus R\cong H_0(\RR^n-B)$.
\item Como $D^n$ es compacto, $f$ es cerrada y por lo tanto es un encaje.
Consideremos $f(D^n)\subset\RR^n$. Por lo anterior, el número de
componentes de $\RR^n-f(D^n)$ es el mismo que el número de componentes de
$\RR^n-D^n$ ($n\ge2$) que es una. Como $\RR^n-f(D^n)$ es no acotado, está
contenido en la componente no acotada de $\RR^n-f(S^{n-1})$, llamémosle
$N$. Esto es, $\RR^n-f(D^n)\subset N$, y por lo tanto $\RR^n-N\subset f(D^n)$.
Tenemos que
\[
A\sqcup f(S^{n-1})=\RR^n-N\subset f(D^n)=f(\mathring D^n)\sqcup f(S^{n-1}).
\]
Entonces $A\subset f(\mathring D^n)$. Como $f(\mathring D^n)$ es conexo y
$A$ es componente $A=f(\mathring D^n)$. Además, $\RR^n$ es localmente
conexo y $\RR^n-f(S^n-f(\mathring D^n))$ es abierto. Por lo tanto, sus
componentes $A$ y $N$ son abiertos y por lo tanto $f(\mathring D^n)$ es
abierto.
\item Sea $x\in U$ y $V\subset U$ una vecindad abierta de $x$. Sea
$e^n\subset V$ una $n$-bola cerrada con centro en $x$ tal que
$e^n\subset V$ es una vecindad de $f(x)$. Por el inciso anterior
$f(\mathring e^n)$ es una vecindad de $f(x)$.
\item Sea $A=f(S^1)$. Ya que $f$ es un encaje, entonces $A$ es homeomorfo
a $S^1$. Sabemos que el complemento del círculo unitario $S^1$ en $S^3$
tiene el mismo tipo de homotopía que $S^1$. Entonces, tenemos que
$\tilde H_q(\RR^3-A;\ZZ)=0$ para todo $q\ne1$ y
$\tilde H_1(\RR^3-A;\ZZ)\cong\ZZ$. Usando Hurewicz se tiene lo deseado.
\end{enumerate}

\item (\textbf{Complejos CW, característica de Euler sobre un campo y
cubrientes}) Sea $X$ un complejo CW \emph{finito} y $(X^{(n)})_{n\in\NN}$
su descomposición celular. La característica de Euler de $X$ con valores
en un campo $k$ se define como el entero
\[
\chi(X,k)=\sum_{q=0}^\infty(-1)^q\dim\big(H_q(X;k)\big).
\]
\begin{enumerate}[label=\alph*)]
\item Demostrar que los grupos de homología $H_q(X,k)$ son de dimensión
finita y nulos para $q$ suficientemente grande. Concluir que $\chi(X,k)$
está bien definido.
\item Demostrar que
$\chi(X,k)=\sum_{q\ge0}(-1)^q\dim\big(H_q(\CWq_q(X;k))\big)$.
\item Demostrar que $\chi(X,k)$ es independiente de $k$.
\item Sean $A,B$ dos subcomplejos celulares de $X$ con $X=A\cup B$.
Demostrar que $\chi(X)=\chi(A)+\chi(B)-\chi(A\cap B)$.
\item Sea $p:X\to B$ un cubriente entre dos complejos CW tal que $X$ es
finito y $B$ conexo. Demostrar que $p$ es un cubriente finito con número
de hojas constante y determinar la característica de Euler de $X$ en
función de $B$.
\item ¿Existen cubrientes de un toro de género $3$ por un toro de género
$6$?
\item Calcular $\chi(\RR P^n)$.
\end{enumerate}

\noindent\textbf{Solución.} Denotamos por $\alpha_q(X)=\#\Gamma^q(X)$ el
número de células de $X$. Recordemos que
\[
\CWq_q(X;k)=H_q(X^{(q)},X^{(q-1)};k)\cong H_q\left(\bigvee_{\alpha_q(X)}S^1;k\right)=k^{\alpha_q(X)}.
\]
\begin{enumerate}[label=\alph*)]
\item Por hipótesis, $X$ siendo finito, $\alpha_q(X)<+\infty$ y
$k^{\alpha_q(X)}$ es de dimensión finita. Por consiguiente,
$\dim\big(H_q(\CWq_\bullet(X;k))\big)<+\infty$. Además, como $X$ es finito,
$\alpha_q(X)=0$ para cierto $q$ suficientemente grande.
\item Si denotamos $a_q=\dim\big(\partial(\CWq_q(X;k))\big)$ y
$b_q=\dim\big(\ker(\partial:\CWq_q(X;k)\to\CWq_{q-1}(X;k))\big)$, tenemos
\[
\dim\big(H_q(X;k)\big)=\dim\big(H_q(\CWq_\bullet(X;k))\big)=b_q-a_{q+1}
\]
y además $\alpha_q=\dim\big(\CWq_q(X;k)\big)=a_q+b_q$. Deducimos entonces
que
\[
\sum_{q=0}^\infty(-1)^q(b_q-a_{q+1})=\sum_{q=0}^\infty\big(b_q+(-1)^{q+1}a_{q+1}\big)=\sum_{q=0}^\infty(-1)^q(b_q+a_q).
\]
\item Se tiene que $\dim(\CWq_q(X;k))=\alpha_q(X)$ y se sigue.
\item Denotemos por $(A^{(n)})_{n\ge0}$ y $(B^{(n)})_{n\ge0}$ las
filtraciones celulares de $A$ y $B$ que verifican, por hipótesis,
$A^{(n)}\subset X^{(n)}$ y $B^{(n)}\subset X^{(n)}$ para todo $n\ge0$. Toda
célula $e\subset A^{(n)}\cap B^{(n)}$ se obtiene recolectando un disco
$D^n$ sobre $X^{(n-1)}$ vía $f:S^{n-1}\to X^{(n-1)}$. Como $A$ y $B$ son
subcomplejos, $f$ tiene imagen en $A^{(n-1)}$ y $B^{(n-1)}$. Se sigue que
$A\cap B$ junto con la filtración $(A\cap B)^{(n)}:=A^{(n)}\cap B^{(n)}$ es
un subcomplejo CW de $X$ (y de $A$ y $B$). Por tanto, $\chi(A\cap B)$ está
bien definida. Por consiguiente, como $X=A\cup B$, las células de
dimensión $n$ de $X$ se particionan en aquellas que están en $A\cap B$,
aquellas que están en $A-(A\cap B)$ y aquellas que están en $B-(A\cap B)$.
Se sigue que
\[
\alpha_n(X)=\alpha_n(A)+\alpha_n(B)-\alpha_n(A\cap B).
\]
\item Como $X$ es un complejo CW finito, es compacto. Se sigue que las
fibras de $p$ son compactas y discretas, y por tanto finitas. Por conexidad
de $B$, la cardinalidad es constante. Sea $k=p^{-1}(b)$. Notemos que $B$ es
compacto debido a que $p$ es sobre y $B$ separado. La estructura de
complejo CW de $B$ se levanta (de manera única) en una estructura celular
sobre $X$ que hace de $p$ una aplicación celular. En efecto, denotemos por
$B^{(n)}$ al $n$-esqueleto de $B$. Pongamos $X^{(0)}=p^{-1}(B^{(0)})$ que
es compacto discreto. Es claro que $p(X^{(0)})\subset B^{(0)}$. Sea
$e_1\subset B^{(1)}$ una $1$-célula de $B$. Tal célula se obtiene
recolectando la frontera de $D^1=[0,1]$ sobre $B^{(0)}$. Sea
$q_1:D^1\to B$ la aplicación así obtenida. Como $D^1$ es simplemente
conexo, para toda elección de $x_0\in p^{-1}(q_1(0))$, esta aplicación se
levanta a una única aplicación $p_1:[0,1]\to X$ con $p(0)=x_0$ y
$p\circ p_1=q_1$. Se sigue de aquí que $p^{-1}(1)\in X^{(0)}$. Obtenemos de
esta manera que $X^{(1)}=p^{-1}(B^{(1)})$ es un complejo CW de dimensión
$1$ en donde las $1$-células están dadas por medio de las aplicaciones
$p_1$. Iterando la construcción precedente obtenemos que existen $k$
células de dimensión $q$ en $X$ por encima de una sola célula de dimensión
$q$ de $B$; en efecto, tal única célula está únicamente determinada por la
elección del levantamiento de un punto base de la célula de $B$. Por
consiguiente, $\alpha_n(X)=k\alpha_n(B)$ para cada $n\ge0$ y entonces
\[
\chi(X)=k\chi(B).
\]
\item Supongamos que existe $p:T_6\to T_3$. Entonces, existe $k\in\NN$ tal
que
\[
-10=2-2(6)=\chi(T_6)=k\chi(T_3)=k(2-2(3))=k(-4).
\]
Esto último es imposible.
\item Consideremos al cubriente $p:S^n\to\RR P^n$. Sabemos que
$\#p^{-1}(b)=2$ para cualquier $b\in\RR P^n$. Entonces,
\[
2\chi(\RR P^n)=\chi(S^n)=\begin{cases}2 & n\text{ par}\\0 & \text{e.o.c.}\end{cases}
\]
Entonces,
\[
\chi(\RR P^n)=\begin{cases}1 & n\text{ par}\\0 & \text{e.o.c.}\end{cases}
\]
\end{enumerate}

\end{enumerate}

\appendix

\section{Cohomología singular, celular y de Rham: interludio sobre espacios lente}

Esta sección añade a las notas precedentes el aparato de cohomología (dual
del de la Sección~\ref{sec:homologia-espacios}) y lo aplica a una familia
de ejemplos --- los \emph{espacios lente} --- ya introducida como ejemplo
de cubriente universal en el Ejemplo~\ref{ex:cubrientes-universales} y
resuelta parcialmente en el Ejercicio 5 de la Sección~2. La razón
pedagógica de detenerse en ellos es que su homología entera se calcula con
un complejo celular extremadamente simple, mientras que distinguir sus
tipos de homotopía exige información más fina: la del anillo de
cohomología con coeficientes finitos.
\label{sec:apendice-cohomologia}

\subsection{Cohomología singular}

\begin{definition}[Complejo singular de cocadenas]
\label{def:cocadenas-singulares}
Sea $X$ un espacio topológico y $G$ un grupo abeliano. Definimos
\[
C^q(X;G):=\Hom(S_q(X;\ZZ),G).
\]
La aplicación \emph{coborde} $\delta^q:C^q(X;G)\to C^{q+1}(X;G)$ está dada
por
\[
\delta^q(\varphi):=\varphi\circ\partial_{q+1}.
\]
Como $\partial_q\circ\partial_{q+1}=0$, se tiene $\delta^{q+1}\circ\delta^q=0$;
así, $(C^\bullet(X;G),\delta^\bullet)$ es un complejo de cocadenas.
\end{definition}

\begin{definition}[Cohomología singular]
\label{def:cohomologia-singular}
El $q$-ésimo grupo de cohomología singular de $X$ con coeficientes en $G$
es
\[
H^q(X;G):=\ker(\delta^q)/\im(\delta^{q-1}).
\]
Para una pareja $(X,A)$ se define análogamente
$H^q(X,A;G):=H^q\big(S_\bullet(X;\ZZ)/S_\bullet(A;\ZZ);G\big)$, dualizando
el complejo de cadenas relativo.
\end{definition}

\begin{proposition}[Funtorialidad]
\label{prop:funtorialidad-cohomologia}
Toda aplicación continua $f:X\to Y$ induce, por precomposición,
homomorfismos $f^\#:C^q(Y;G)\to C^q(X;G)$, $f^\#(\varphi)=\varphi\circ f_\#$,
y en consecuencia morfismos en cohomología $f^*:H^q(Y;G)\to H^q(X;G)$.
Además, $(\mathrm{id}_X)^*=\mathrm{id}$ y $(g\circ f)^*=f^*\circ g^*$.
\end{proposition}

\begin{proof}
$f^\#$ es morfismo de cocadenas: para $\varphi\in C^q(Y;G)$,
\[
f^\#(\delta^q\varphi)=(\delta^q\varphi)\circ f_\#=\varphi\circ\partial_{q+1}\circ f_\#=\varphi\circ f_\#\circ\partial_{q+1}=\delta^q(f^\#\varphi),
\]
usando que $f_\#$ conmuta con $\partial$ (Definición~\ref{def:homologia-singular}).
Así $f^\#$ preserva cociclos y cofronteras, y desciende a
$f^*:H^q(Y;G)\to H^q(X;G)$. Que $(\mathrm{id}_X)^\#=\mathrm{id}$ es
inmediato, y $(g\circ f)^\#(\varphi)=\varphi\circ(g\circ f)_\#=\varphi\circ g_\#\circ f_\#=f^\#(g^\#\varphi)=(f^\#\circ g^\#)(\varphi)$
da $(g\circ f)^*=f^*\circ g^*$ al pasar a cohomología.
\end{proof}

\begin{theorem}[Invarianza homotópica]
\label{thm:invarianza-homotopica-cohomologia}
Si $f,g:X\to Y$ son homotópicas, entonces $f^*=g^*:H^q(Y;G)\to H^q(X;G)$
para toda $q$. En particular, la cohomología es un invariante del tipo de
homotopía.
\end{theorem}

\begin{proof}
Sea $P:S_q(X)\to S_{q+1}(Y)$ el operador prisma de la demostración del
Teorema~\ref{thm:invarianza-homotopica}, con
$\partial P+P\partial=g_\#-f_\#$. Dualizando (aplicando $\Hom(-,G)$),
definimos $P^*:C^{q+1}(Y;G)\to C^q(X;G)$, $P^*(\varphi):=\varphi\circ P$.
Para $\varphi\in C^q(Y;G)$ (cociclo o no),
\[
(\delta^{q-1}P^*+P^*\delta^q)(\varphi)=\varphi\circ P\circ\partial+\varphi\circ\partial\circ P=\varphi\circ(P\partial+\partial P)=\varphi\circ(g_\#-f_\#)=g^\#(\varphi)-f^\#(\varphi).
\]
Así $P^*$ es una homotopía de cocadenas entre $f^\#$ y $g^\#$; si $\varphi$
es un cociclo ($\delta^q\varphi=0$), entonces
$g^\#(\varphi)-f^\#(\varphi)=\delta^{q-1}(P^*\varphi)$ es una cofrontera,
así que $[g^\#\varphi]=[f^\#\varphi]$ en $H^q(X;G)$, es decir
$g^*=f^*$.
\end{proof}

\begin{theorem}[Sucesión exacta larga de una pareja]
\label{thm:sucesion-larga-cohomologia}
Para toda pareja $(X,A)$ y todo grupo abeliano $G$ existe una sucesión
exacta larga natural
\[
\cdots\to H^q(X,A;G)\to H^q(X;G)\to H^q(A;G)\xrightarrow{\delta}H^{q+1}(X,A;G)\to H^{q+1}(X;G)\to\cdots
\]
\end{theorem}

\begin{proof}
La sucesión corta exacta de complejos de cadenas
$0\to S_\bullet(A)\to S_\bullet(X)\to S_\bullet(X)/S_\bullet(A)\to0$
(demostración del Teorema~\ref{thm:sucesion-larga}) consiste, en cada
grado, en grupos abelianos libres: $S_q(A)$ y $S_q(X)$ son libres por
definición, y $S_q(X)/S_q(A)$ es libre porque $S_q(A)$ es generado por un
subconjunto de la base de símplices singulares de $S_q(X)$ (un
``sumando coordenado''), de modo que la sucesión se escinde en cada
grado (como grupos abelianos, olvidando la diferencial). Aplicar
$\Hom(-,G)$ a una sucesión corta exacta escindida de grupos abelianos
preserva la exactitud (la escisión se preserva bajo cualquier funtor
aditivo), dando una nueva sucesión corta exacta de complejos de
cocadenas
\[
0\to C^\bullet(X/A;G):=\Hom\big(S_\bullet(X)/S_\bullet(A),G\big)\to C^\bullet(X;G)\to C^\bullet(A;G)\to0,
\]
en donde el primer término calcula $H^q(X,A;G)$ por definición
(Definición~\ref{def:cohomologia-singular}). El lema del zig-zag
(exactamente como en la demostración del Teorema~\ref{thm:sucesion-larga},
ahora con las flechas de conexión invertidas por la naturaleza
contravariante de $\Hom$) da la sucesión exacta larga enunciada.
\end{proof}

\begin{theorem}[Escisión en cohomología]
\label{thm:escision-cohomologia}
Si $U\subset A\subset X$ y $\overline U\subset\mathring A$, entonces la
inclusión $(X-U,A-U)\hookrightarrow(X,A)$ induce isomorfismos
\[
H^q(X,A;G)\xrightarrow{\cong}H^q(X-U,A-U;G)
\]
para toda $q$.
\end{theorem}

\begin{proof}
Por la demostración del Teorema~\ref{thm:escision}, la inclusión de
complejos de cadenas $j_\#:S_\bullet(X-U)/S_\bullet(A-U)\hookrightarrow S_\bullet(X)/S_\bullet(A)$
es una equivalencia de homotopía de cadenas (se construyó explícitamente
un inverso homotópico $\rho$ y homotopías $D$ con
$\partial D+D\partial=\mathrm{id}-\rho j_\#$, restringidas al cociente
relativo). Aplicando $\Hom(-,G)$, los operadores $D^*=\Hom(D,G)$ dan una
homotopía de cocadenas entre $\mathrm{id}$ y $(\rho j_\#)^*=j_\#^*\rho^*$,
mostrando que $j_\#^*:C^\bullet(X,A;G)\to C^\bullet(X-U,A-U;G)$ es
también una equivalencia de homotopía de cocadenas (con inversa
$\rho^*$). Por lo tanto induce isomorfismo en cohomología,
$j^*:H^q(X,A;G)\xrightarrow{\cong}H^q(X-U,A-U;G)$.
\end{proof}

\begin{theorem}[Mayer--Vietoris en cohomología]
\label{thm:mayer-vietoris-cohomologia}
Si $X=U\cup V$ con $U,V$ abiertos, entonces existe una sucesión exacta
larga natural
\[
\cdots\to H^q(X;G)\to H^q(U;G)\oplus H^q(V;G)\to H^q(U\cap V;G)\to H^{q+1}(X;G)\to\cdots
\]
\end{theorem}

\begin{proof}
Igual que en la demostración del Teorema~\ref{thm:mayer-vietoris} (caso
$A_1=B_1=\varnothing$), se tiene la sucesión corta exacta de cadenas
\[
0\to S_\bullet(U\cap V)\xrightarrow{(i_\#,-j_\#)}S_\bullet(U)\oplus S_\bullet(V)\xrightarrow{\psi}S_\bullet(U)+S_\bullet(V)\to0,
\]
que se escinde en cada grado como grupos abelianos libres (mismo
argumento que en la prueba anterior: cada término es libre, generado por
símplices, y la sucesión se escinde por bases). Aplicando $\Hom(-,G)$ se
obtiene una sucesión corta exacta de cocadenas, y el lema del zig-zag
(con flechas invertidas) da una sucesión exacta larga en
$H^\bullet(S_\bullet(U)+S_\bullet(V);G)$, $H^\bullet(U;G)\oplus H^\bullet(V;G)$
y $H^\bullet(U\cap V;G)$. Como $\{U,V\}$ es escisiva (abiertos que cubren
$X$, Nota~\ref{nota:mayer-vietoris-abiertos}), la inclusión
$S_\bullet(U)+S_\bullet(V)\hookrightarrow S_\bullet(X)$ es una
equivalencia de homotopía de cadenas (mismo argumento del Teorema de
escisión, con la cubierta $\{U,V\}$ directamente, sin necesidad de
excindir nada), y dualizando como en la demostración de
escisión en cohomología, se identifica
$H^\bullet(S_\bullet(U)+S_\bullet(V);G)\cong H^\bullet(X;G)$, dando la
sucesión enunciada.
\end{proof}

\subsection{Coeficientes universales y cohomología celular}

\begin{theorem}[Coeficientes universales para cohomología]
\label{thm:coeficientes-universales}
Para todo espacio $X$ y todo grupo abeliano $G$ existe una sucesión exacta
corta natural
\[
0\to\mathrm{Ext}^1_\ZZ\big(H_{q-1}(X;\ZZ),G\big)\to H^q(X;G)\to\Hom\big(H_q(X;\ZZ),G\big)\to0.
\]
Esta sucesión se escinde (pues $\Hom(H_q(X;\ZZ),G)$ es libre siempre que
$H_q(X;\ZZ)$ es finitamente generado), aunque no de manera natural.
\end{theorem}

\begin{proof}
Sea $C_\bullet=S_\bullet(X;\ZZ)$, con $Z_q=\ker\partial_q$ (ciclos) y
$B_{q-1}=\im\partial_q\subset Z_{q-1}$ (fronteras), y sea
$\iota_q:B_q\hookrightarrow Z_q$ la inclusión, de modo que
$H_q(X;\ZZ)=Z_q/B_q$ por definición.

\emph{Paso 1 (Ext vía la sucesión de $B_q,Z_q$).} Como $C_q$ es libre,
$Z_q\subset C_q$ es libre, y $B_{q-1}=\partial_q(C_q)\cong C_q/Z_q$
también (subgrupo, resp.\ cociente por un sumando directo, de un grupo
libre). Aplicando $\Hom(-,G)$ a
$0\to B_q\xrightarrow{\iota_q}Z_q\to H_q(X;\ZZ)\to0$ se obtiene la
sucesión exacta larga de Ext,
\[
0\to\Hom(H_q(X;\ZZ),G)\to\Hom(Z_q,G)\xrightarrow{\iota_q^*}\Hom(B_q,G)\to\mathrm{Ext}^1_\ZZ(H_q(X;\ZZ),G)\to\mathrm{Ext}^1_\ZZ(Z_q,G)=0
\]
(el último término se anula porque $Z_q$, siendo libre, es proyectivo).
De aquí,
\[
\Hom(H_q(X;\ZZ),G)\cong\ker(\iota_q^*), \qquad \mathrm{Ext}^1_\ZZ(H_q(X;\ZZ),G)\cong\mathrm{coker}(\iota_q^*). \tag{$*$}
\]

\emph{Paso 2 (descomposición del complejo de cocadenas).} Como
$B_{q-1}$ es libre, la sucesión $0\to Z_q\to C_q\xrightarrow{\partial_q}B_{q-1}\to0$
se escinde: elíjase, para cada $q$, una sección $s_q:B_{q-1}\to C_q$ de
$\partial_q$ (es decir $\partial_q\circ s_q=\mathrm{id}_{B_{q-1}}$), dando
un isomorfismo (no natural) $C_q\cong Z_q\oplus s_q(B_{q-1})$, y por lo
tanto
\[
C^q(X;G)=\Hom(C_q,G)\cong\Hom(Z_q,G)\oplus\Hom(B_{q-1},G).
\]
Bajo esta descomposición, para $(\varphi,\psi)\in\Hom(Z_q,G)\oplus\Hom(B_{q-1},G)$
(vista como elemento de $C^q(X;G)$ vía la proyección $C_q\to Z_q$ y
$s_q^{-1}$ en el sumando complementario), el coborde
$\delta^q(\varphi,\psi)=(\varphi,\psi)\circ\partial_{q+1}$ se calcula
sobre $C_{q+1}=Z_{q+1}\oplus s_{q+1}(B_q)$: en $Z_{q+1}$, $\partial_{q+1}$
es cero (son ciclos); en $s_{q+1}(B_q)$, $\partial_{q+1}\circ s_{q+1}=\mathrm{id}_{B_q}$
seguido de la inclusión $\iota_q:B_q\hookrightarrow Z_q\subset C_q$. Así,
\[
\delta^q(\varphi,\psi)=(0,\ \varphi\circ\iota_q)\ \in\ \Hom(Z_{q+1},G)\oplus\Hom(B_q,G)=C^{q+1}(X;G)
\]
--- nótese que $\psi$ no interviene: es exactamente la ambigüedad que
Ext detecta.

\emph{Paso 3 (cociclos y cofronteras).} Un elemento $(\varphi,\psi)\in C^q(X;G)$
es cociclo si y solo si $\varphi\circ\iota_q=0$, es decir
$\varphi\in\ker(\iota_q^*)$ (y $\psi$ es arbitrario); por $(*)$,
\[
Z^q(X;G)\cong\Hom(H_q(X;\ZZ),G)\oplus\Hom(B_{q-1},G).
\]
Las cofronteras en grado $q$ son, por el Paso 2 aplicado en grado $q-1$,
exactamente los elementos de la forma $(0,\varphi'\circ\iota_{q-1})$ para
$\varphi'\in\Hom(Z_{q-1},G)$, es decir
\[
B^q(X;G)\cong\{0\}\oplus\im(\iota_{q-1}^*)\subset\Hom(H_q(X;\ZZ),G)\oplus\Hom(B_{q-1},G).
\]
Por lo tanto,
\[
H^q(X;G)=Z^q(X;G)/B^q(X;G)\cong\Hom(H_q(X;\ZZ),G)\ \oplus\ \Hom(B_{q-1},G)/\im(\iota_{q-1}^*),
\]
y por $(*)$ (aplicado en grado $q-1$), $\Hom(B_{q-1},G)/\im(\iota_{q-1}^*)\cong\mathrm{Ext}^1_\ZZ(H_{q-1}(X;\ZZ),G)$.
Esto da
\[
H^q(X;G)\cong\Hom(H_q(X;\ZZ),G)\oplus\mathrm{Ext}^1_\ZZ(H_{q-1}(X;\ZZ),G),
\]
en particular la sucesión corta exacta del enunciado, escindida (la
escisión depende de la elección no canónica de las secciones $s_q$, por
lo que no es natural en $X$).
\end{proof}

\begin{corollary}
\label{cor:cohomologia-libre}
Si $H_{q-1}(X;\ZZ)$ es libre, entonces $H^q(X;G)\cong\Hom(H_q(X;\ZZ),G)$.
En particular, si toda la homología entera de $X$ es libre, la cohomología
entera es el dual de la homología.
\end{corollary}

\begin{proof}
Si $H_{q-1}(X;\ZZ)$ es libre, entonces
$\mathrm{Ext}^1_\ZZ(H_{q-1}(X;\ZZ),G)=0$ (todo grupo abeliano libre es
proyectivo), así que el Teorema~\ref{thm:coeficientes-universales} da
directamente $H^q(X;G)\cong\Hom(H_q(X;\ZZ),G)$.
\end{proof}

\begin{definition}[Complejo celular de cocadenas]
\label{def:cocadenas-celulares}
Sea $X$ un complejo CW. Definimos
\[
C^q_{\CWq}(X;G):=\Hom\big(C_q^{\CWq}(X;\ZZ),G\big).
\]
Si $d_{q+1}:C_{q+1}^{\CWq}(X;\ZZ)\to C_q^{\CWq}(X;\ZZ)$ es el borde
celular, definimos el coborde celular por
$\delta^q:=\Hom(d_{q+1},G)$.
\end{definition}

\begin{theorem}[Cohomología celular]
\label{thm:cohomologia-celular}
Sea $X$ un complejo CW. Entonces $H^q_{\CWq}(X;G)\cong H^q(X;G)$ para toda
$q\ge0$.
\end{theorem}

\begin{proof}
Los grupos celulares $C_q^{\CWq}(X;\ZZ)=H_q(X^q,X^{q-1};\ZZ)$ son libres
(Teorema~\ref{thm:cw-celular}.1). El mismo argumento que en la
demostración del Teorema~\ref{thm:coeficientes-universales}, aplicado al
complejo celular en lugar del singular (posible porque cada
$C_q^{\CWq}(X;\ZZ)$ es libre, la hipótesis que se usó), muestra que
$H^q_{\CWq}(X;G)$ encaja en una sucesión corta exacta natural con
$\mathrm{Ext}^1_\ZZ(H_{q-1}^{\CWq}(X;\ZZ),G)$ y $\Hom(H_q^{\CWq}(X;\ZZ),G)$;
por el Teorema~\ref{thm:cw-iso-singular}, $H_q^{\CWq}(X;\ZZ)\cong H_q(X;\ZZ)$
para toda $q$, así que esta sucesión coincide (naturalmente) con la del
Teorema~\ref{thm:coeficientes-universales} para $H^q(X;G)$. Por unicidad
del término medio de una sucesión corta exacta con extremos dados (hasta
isomorfismo compatible), $H^q_{\CWq}(X;G)\cong H^q(X;G)$.
\end{proof}

\begin{corollary}[Esferas]
\label{cor:cohomologia-esferas}
Para todo grupo abeliano $G$ y $n\ge1$,
\[
H^q(S^n;G)=\begin{cases}G & q=0,n\\0 & \text{en otro caso.}\end{cases}
\]
\end{corollary}

\begin{proof}
$S^n$ admite una estructura CW con una $0$-célula y una $n$-célula
(Teorema~\ref{thm:cw-lente} con $m=1$, o directamente: $S^n\cong\{*\}\cup_fD^n$
con $f$ constante). El complejo celular de cadenas es
$0\to\ZZ\to0\to\cdots\to0\to\ZZ\to0$ (con $\ZZ$ en grados $0$ y $n$, y
diferenciales cero, pues no hay dos células en dimensiones adyacentes),
dando $H_q(S^n;\ZZ)=\ZZ$ para $q=0,n$ y $0$ en otro caso, con
$H_{q-1}(S^n;\ZZ)=0$ libre para todo $q$. Por el
Corolario~\ref{cor:cohomologia-libre},
$H^q(S^n;G)\cong\Hom(H_q(S^n;\ZZ),G)$, que es $\Hom(\ZZ,G)\cong G$ para
$q=0,n$, y $0$ en otro caso.
\end{proof}

\begin{corollary}[Bouquets]
\label{cor:cohomologia-bouquets}
Si $\{X_\alpha\}_{\alpha\in I}$ es una familia de complejos CW conexos,
entonces
\[
\tilde H^q\left(\bigvee_{\alpha\in I}X_\alpha;G\right)\cong\bigoplus_{\alpha\in I}\tilde H^q(X_\alpha;G)
\]
para toda $q\ge1$.
\end{corollary}

\begin{proof}
Por el Corolario~\ref{cor:homologia-bouquet},
$\tilde H_q\big(\bigvee_\alpha X_\alpha;\ZZ\big)\cong\bigoplus_\alpha\tilde H_q(X_\alpha;\ZZ)$
para toda $q$; en particular, si cada $\tilde H_{q-1}(X_\alpha;\ZZ)$ es
libre para toda $\alpha$, la suma directa también lo es (suma directa de
grupos libres es libre), y el Corolario~\ref{cor:cohomologia-libre} da
\[
\tilde H^q\Big(\bigvee_\alpha X_\alpha;G\Big)\cong\Hom\Big(\bigoplus_\alpha\tilde H_q(X_\alpha;\ZZ),G\Big)\cong\prod_\alpha\Hom(\tilde H_q(X_\alpha;\ZZ),G).
\]
Para $I$ finito (el caso relevante para complejos CW con un número finito
de células en cada dimensión, que es la situación de interés en estas
notas --- por ejemplo los bouquets de esferas usados en la homología
celular), $\prod_\alpha=\bigoplus_\alpha$, y cada factor es
$\Hom(\tilde H_q(X_\alpha;\ZZ),G)\cong\tilde H^q(X_\alpha;G)$ de nuevo por
el Corolario~\ref{cor:cohomologia-libre} (aplicado a cada $X_\alpha$
individualmente, bajo la misma hipótesis de libertad en grado $q-1$, que
se cumple automáticamente si cada $\tilde H_{q-1}(X_\alpha;\ZZ)$ es libre,
en particular si es $0$). En el caso general (homología no
necesariamente libre en grado $q-1$), el resultado se sigue en su lugar
directamente de la Nota~\ref{nota:cociente-pareja-buena} y su análogo en
cohomología: $\bigvee_\alpha X_\alpha$ es buena pareja respecto al punto
base, y $H^\bullet(X,x_0;G)\cong\tilde H^\bullet(X;G)$ (dual de la
identificación usada para homología), reduciendo el enunciado al
análogo, ya establecido, del Ejercicio 1f/g en cohomología.
\end{proof}

\begin{remark}
\label{obs:homologia-vs-cohomologia}
La homología describe la estructura aditiva del espacio; la cohomología,
además de esa estructura aditiva, retiene productos entre clases (el
producto de taza, \S\ref{sec:anillo-cohomologia-lente}). Por eso dos
espacios pueden compartir la misma homología y no tener el mismo anillo de
cohomología --- fenómeno que se ilustra más abajo con
$\CC P^2\not\simeq S^2\vee S^4$ (Ejercicio~\ref{ex:anillo-proyectivos}) y,
entre los espacios lente, con $L(5,1)\not\simeq L(5,2)$
(Ejemplo~\ref{ex:L51-L52}).
\end{remark}

\subsection{Espacios lente: definición y complejo celular}
\label{sec:espacios-lente-apendice}

Los espacios lente forman una familia fundamental de variedades cerradas
orientables de dimensión impar. Desde el punto de vista pedagógico son
decisivos porque su homología integral se calcula por un complejo celular
extremadamente simple, mientras que su clasificación homotópica exige
información más fina proveniente de la cohomología con coeficientes
finitos y de la estructura multiplicativa del producto de taza.

\begin{definition}[Espacio lente]
\label{def:espacio-lente-apendice}
Sean $m\ge2$ y $\ell_1,\dots,\ell_n\in\ZZ$ con $\gcd(\ell_j,m)=1$ para todo
$j$. Sea $\xi=e^{2\pi i/m}$. El grupo cíclico $\ZZ/m\ZZ=\langle\xi\rangle$
actúa libremente en $S^{2n-1}\subset\CC^n$ por
\[
\xi\cdot(z_1,\dots,z_n)=(\xi^{\ell_1}z_1,\dots,\xi^{\ell_n}z_n).
\]
El cociente
\[
L_m(\ell_1,\dots,\ell_n):=S^{2n-1}/(\ZZ/m\ZZ)
\]
se llama \emph{espacio lente} de dimensión $2n-1$. Cuando $n=2$ se escribe
$L(m,\ell):=L_m(1,\ell)$, el \emph{espacio lente clásico} tridimensional.
\end{definition}

\begin{center}
\begin{tikzpicture}[
  scale=1,
  pole/.style={circle,fill=black,inner sep=1.1pt},
  meridian/.style={thick},
  meridianB/.style={thick,dashed},
  >=Stealth
]

  \coordinate (N) at (0,2.1);
  \coordinate (S) at (0,-2.1);


  \draw[thick] (N) to[out=-25,in=25] (S);
  \draw[thick] (N) to[out=205,in=155] (S);


  \draw[thick] (0,0) ellipse [x radius=1.9, y radius=0.55];


  \foreach \k in {1,2,3} {
    \pgfmathsetmacro{\ang}{90-\k*60}
    \pgfmathsetmacro{\eqx}{1.9*cos(\ang)}
    \pgfmathsetmacro{\eqy}{0.55*sin(\ang)}
    \pgfmathsetmacro{\outa}{-90+\ang*0.28}
    \pgfmathsetmacro{\ina}{90-\ang*0.28}
    \draw[meridian] (N) to[out=\outa,in=90] (\eqx,\eqy);
    \draw[meridian] (\eqx,\eqy) to[out=-90,in=\ina] (S);
  }

  \pgfmathsetmacro{\bqx}{1.9*cos(210)}
  \pgfmathsetmacro{\bqy}{0.55*sin(210)}
  \draw[meridianB] (N) to[out=-108,in=90] (\bqx,\bqy);
  \draw[meridianB] (\bqx,\bqy) to[out=-90,in=108] (S);


  \node[pole,label={above:$N$}] at (N) {};
  \node[pole,label={below:$S$}] at (S) {};


  \draw[->,thick] (2.3,0.15) to[out=20,in=-70] node[midway,right] {\footnotesize $\xi\cdot$} (2.55,1.05);
  \node[right] at (2.65,1.15) {\footnotesize rotación $2\pi\ell/m$};


  \node at (0,-2.75) {\footnotesize $D^3=\{$casquete superior$\}\cup\{$casquete inferior$\}$, identificados vía $\xi$};
\end{tikzpicture}
\end{center}
\begin{center}
\footnotesize\emph{El espacio lente clásico $L(m,\ell)$: la bola $D^3$
dibujada como una lente biconvexa (de ahí el nombre), con su casquete
superior pegado al inferior mediante la rotación $\xi\cdot$ de ángulo
$2\pi\ell/m$ alrededor del eje polar --- restricción real de la acción de
$\ZZ/m\ZZ=\langle\xi\rangle$ en $S^3\subset\CC^2$ de la
Definición~\ref{def:espacio-lente-apendice}.}
\end{center}

\begin{remark}
\label{obs:accion-libre-lente}
La hipótesis $\gcd(\ell_j,m)=1$ garantiza que la acción es libre: si
$\xi^k\cdot z=z$ y alguna coordenada $z_j$ es no nula, entonces
$\xi^{k\ell_j}=1$, de donde $m\mid k\ell_j$. Como $\ell_j$ es invertible
módulo $m$, se concluye $m\mid k$, luego $\xi^k=1$. Esto es exactamente el
argumento ya usado en el Ejercicio 5c de la Sección~2 para calcular
$\pi_1(L_m(\ell_1,\dots,\ell_n))\cong\ZZ/m\ZZ$.
\end{remark}

\begin{theorem}[Complejo celular de un espacio lente]
\label{thm:cw-lente}
Todo espacio lente $L=L_m(\ell_1,\dots,\ell_n)$ admite una estructura CW
con exactamente una celda en cada dimensión $0,1,2,\dots,2n-1$. Con
respecto a bases adecuadas, el complejo celular de cadenas tiene la forma
\[
0\to C_{2n-1}\cong\ZZ\xrightarrow{0}C_{2n-2}\cong\ZZ\xrightarrow{m}C_{2n-3}\cong\ZZ\xrightarrow{0}\cdots\xrightarrow{m}C_1\cong\ZZ\xrightarrow{0}C_0\cong\ZZ\to0.
\]
Equivalentemente, $d_{2j+1}=0$ y $d_{2j}=m$ para $1\le j\le n-1$.
\end{theorem}

\begin{proof}
La esfera $S^{2n-1}$ admite una descomposición CW equivariante bajo la
acción de $\ZZ/m\ZZ$, y al pasar al cociente aparece exactamente una celda
en cada dimensión. El hecho geométrico clave es que la cara de una celda
impar se pega de manera que las contribuciones se cancelan, mientras que en
dimensión par el borde recubre $m$ veces la celda anterior. Traducido al
complejo celular, eso produce precisamente la alternancia $0,m,0,m,\dots,0$.
El resto del cálculo es puramente algebraico.
\end{proof}

\begin{corollary}[Homología integral de los espacios lente]
\label{cor:homologia-lente}
Para $L=L_m(\ell_1,\dots,\ell_n)$ se tiene
\[
H_q(L;\ZZ)\cong\begin{cases}
\ZZ & q=0\text{ o }q=2n-1,\\
\ZZ/m\ZZ & q\text{ impar},\ 1\le q\le2n-3,\\
0 & q\text{ par},\ 0<q<2n-1.
\end{cases}
\]
\end{corollary}

\begin{proof}
Del Teorema~\ref{thm:cw-lente} se obtiene, para $1\le j\le n-1$,
\[
H_{2j}(L;\ZZ)=\frac{\ker(d_{2j})}{\im(d_{2j+1})}=\frac{0}{0}=0,
\]
porque $d_{2j}=m:\ZZ\to\ZZ$ es inyectiva y $d_{2j+1}=0$, mientras que
\[
H_{2j-1}(L;\ZZ)=\frac{\ker(d_{2j-1})}{\im(d_{2j})}=\frac{\ZZ}{m\ZZ}\cong\ZZ/m\ZZ.
\]
Además, $H_{2n-1}(L;\ZZ)=\ker(d_{2n-1})=\ZZ$ y $H_0(L;\ZZ)=\ZZ$, ya que $L$
es conexo.
\end{proof}

\begin{corollary}[Caso tridimensional]
\label{cor:homologia-lente-3d}
Para todo espacio lente clásico $L(m,\ell)$,
\[
H_q(L(m,\ell);\ZZ)\cong\begin{cases}\ZZ & q=0,3\\ \ZZ/m\ZZ & q=1\\0 & q=2.\end{cases}
\]
\end{corollary}

\begin{proof}
Se trata del caso $n=2$. El complejo celular queda
$0\to\ZZ\xrightarrow{0}\ZZ\xrightarrow{m}\ZZ\xrightarrow{0}\ZZ\to0$, y el
cálculo es inmediato del Corolario~\ref{cor:homologia-lente}.
\end{proof}

\subsection{Cohomología celular de espacios lente, coeficientes universales y anillo}
\label{sec:anillo-cohomologia-lente}

La cohomología de los espacios lente se obtiene dualizando el complejo
celular anterior. Éste es el punto en el que la teoría de cohomología se
vuelve tan efectiva como lo fue la homología en la subsección precedente.

\begin{theorem}[Cohomología integral de los espacios lente]
\label{thm:cohomologia-lente}
Para $L=L_m(\ell_1,\dots,\ell_n)$ se tiene
\[
H^q(L;\ZZ)\cong\begin{cases}
\ZZ & q=0\text{ o }q=2n-1,\\
\ZZ/m\ZZ & q\text{ par},\ 2\le q\le2n-2,\\
0 & q\text{ impar},\ 1\le q\le2n-3.
\end{cases}
\]
\end{theorem}

\begin{proof}
Dualizando el complejo celular de cadenas, obtenemos el complejo de
cocadenas
\[
0\to C^0\cong\ZZ\xrightarrow{0}C^1\cong\ZZ\xrightarrow{m}C^2\cong\ZZ\xrightarrow{0}C^3\cong\ZZ\xrightarrow{m}\cdots\xrightarrow{0}C^{2n-1}\cong\ZZ\to0,
\]
es decir, $\delta^{2j-1}=m$ y $\delta^{2j}=0$ para $1\le j\le n-1$.
Entonces
\[
H^{2j-1}(L;\ZZ)=\frac{\ker(\delta^{2j-1})}{\im(\delta^{2j-2})}=0
\]
porque $\delta^{2j-1}=m:\ZZ\to\ZZ$ es inyectiva, y
\[
H^{2j}(L;\ZZ)=\frac{\ker(\delta^{2j})}{\im(\delta^{2j-1})}=\frac{\ZZ}{m\ZZ}\cong\ZZ/m\ZZ.
\]
Finalmente, $H^0(L;\ZZ)=\ZZ$ y $H^{2n-1}(L;\ZZ)=\ZZ$.
\end{proof}

Obsérvese que la torsión $\ZZ/m\ZZ$ aparece en homología en grados
\emph{impares} ($1\le q\le2n-3$) y en cohomología en grados \emph{pares}
($2\le q\le2n-2$): un grado más arriba. Este desplazamiento es exactamente
el que produce el término $\mathrm{Ext}^1$ del Teorema de Coeficientes
Universales, y es el mismo fenómeno, a menor escala, que distingue
$H_2(\RR P^2;\ZZ)=0$ de $H^2(\RR P^2;\ZZ)\cong\ZZ/2\ZZ$.

\begin{corollary}[Caso tridimensional]
\label{cor:cohomologia-lente-3d}
Para todo espacio lente clásico $L(m,\ell)$,
\[
H^q(L(m,\ell);\ZZ)\cong\begin{cases}\ZZ & q=0,3\\ \ZZ/m\ZZ & q=2\\0 & q=1.\end{cases}
\]
\end{corollary}

\begin{proof}
Es el caso $n=2$ del Teorema~\ref{thm:cohomologia-lente}.
\end{proof}

\begin{theorem}[Coeficientes universales en el caso lente]
\label{thm:ucl-lente}
Para todo espacio lente $L=L_m(\ell_1,\dots,\ell_n)$ y todo grupo abeliano
$G$,
\[
0\to\mathrm{Ext}^1_\ZZ\big(H_{q-1}(L;\ZZ),G\big)\to H^q(L;G)\to\Hom\big(H_q(L;\ZZ),G\big)\to0.
\]
En particular:
\begin{enumerate}[label=(\roman*)]
\item si $G=\QQ$, entonces
\[
H^q(L;\QQ)\cong\begin{cases}\QQ & q=0,2n-1\\0 & \text{en otro caso;}\end{cases}
\]
\item si $G=\ZZ/m\ZZ$, entonces $H^q(L;\ZZ/m\ZZ)\cong\ZZ/m\ZZ$ para todo
$0\le q\le2n-1$.
\end{enumerate}
\end{theorem}

\begin{proof}
La afirmación general es el Teorema de Coeficientes
Universales~\ref{thm:coeficientes-universales} aplicado a la homología del
Corolario~\ref{cor:homologia-lente}. Si $G=\QQ$, la torsión desaparece
porque $\Hom(\ZZ/m\ZZ,\QQ)=0$ y $\mathrm{Ext}^1_\ZZ(\ZZ/m\ZZ,\QQ)=0$. Si
$G=\ZZ/m\ZZ$, se tiene $\Hom(\ZZ,\ZZ/m\ZZ)\cong\ZZ/m\ZZ$,
$\Hom(\ZZ/m\ZZ,\ZZ/m\ZZ)\cong\ZZ/m\ZZ$ y
$\mathrm{Ext}^1_\ZZ(\ZZ/m\ZZ,\ZZ/m\ZZ)\cong\ZZ/m\ZZ$, de modo que en cada
grado aparece un grupo isomorfo a $\ZZ/m\ZZ$ (con extensión trivial cuando
el término de $\Hom$ correspondiente se anula).
\end{proof}

\begin{theorem}[Anillo de cohomología módulo $m$]
\label{thm:anillo-cohomologia-lente}
Sea $L=L_m(\ell_1,\dots,\ell_n)$ de dimensión $2n-1$.
\begin{enumerate}[label=(\roman*)]
\item Si $m$ es impar, entonces
\[
H^*(L;\ZZ/m\ZZ)\cong\Lambda(\alpha)\otimes\frac{(\ZZ/m\ZZ)[\beta]}{(\beta^n)},
\qquad |\alpha|=1,\ |\beta|=2,
\]
con $\alpha\beta^{n-1}$ generando $H^{2n-1}(L;\ZZ/m\ZZ)$.
\item Si $m=2$, entonces
\[
H^*(L;\ZZ/2\ZZ)\cong\frac{(\ZZ/2\ZZ)[\alpha]}{(\alpha^{2n})}, \qquad |\alpha|=1.
\]
\end{enumerate}
\end{theorem}

\begin{proof}[Idea de la demostración --- no reproducida en detalle]
Los grupos aditivos subyacentes ya están determinados por el
Teorema~\ref{thm:cohomologia-lente} (o el Teorema~\ref{thm:ucl-lente}(ii)
para $G=\ZZ/m\ZZ$); lo que falta es el \emph{producto de taza}. Éste no
puede calcularse solo con la sucesión exacta larga o el complejo
celular de cadenas --- ambos son invariantes puramente aditivos, como
señala la Observación~\ref{obs:homologia-vs-cohomologia} --- sino que
requiere modelar $S^{2n-1}$ como el borde de una resolución libre
explícita de $\ZZ$ sobre el anillo de grupo $\ZZ[\ZZ/m\ZZ]$: el
\emph{complejo estándar periódico}
\[
\cdots\xrightarrow{N}\ZZ[\ZZ/m\ZZ]\xrightarrow{T-1}\ZZ[\ZZ/m\ZZ]\xrightarrow{N}\ZZ[\ZZ/m\ZZ]\xrightarrow{T-1}\ZZ[\ZZ/m\ZZ]\to\ZZ\to0,
\]
donde $T$ es un generador de $\ZZ/m\ZZ$ actuando por multiplicación y
$N=1+T+\cdots+T^{m-1}$, y comparar (vía una equivalencia de cadenas
$\ZZ/m\ZZ$-equivariante) el complejo celular de $S^{2n-1}$ con este
complejo algebraico para obtener representantes explícitos de cociclos
cuyo producto de taza se computa directamente. Esta construcción es
estándar pero extensa; no se reproduce aquí y se cita como resultado
clásico (Reidemeister; véase Hatcher, \emph{Algebraic Topology}, Ejemplo
3.41 y Teorema 3.44, para el argumento completo). Los
Teoremas~\ref{thm:cw-lente}, \ref{thm:cohomologia-lente} y
\ref{thm:ucl-lente} de esta nota, en cambio, sí están demostrados de
manera autocontenida.
\end{proof}

\begin{corollary}[Caso tridimensional]
\label{cor:anillo-lente-3d}
Si $L=L(m,\ell)$ tiene dimensión $3$, entonces:
\begin{enumerate}[label=(\roman*)]
\item si $m$ es impar,
\[
H^*(L;\ZZ/m\ZZ)\cong\Lambda(\alpha)\otimes\frac{(\ZZ/m\ZZ)[\beta]}{(\beta^2)}, \qquad |\alpha|=1,\ |\beta|=2;
\]
en particular, $\alpha^2=0$, $\beta^2=0$, $\alpha\beta\ne0$;
\item si $m=2$,
\[
H^*(L(2,1);\ZZ/2\ZZ)\cong\frac{(\ZZ/2\ZZ)[\alpha]}{(\alpha^4)}.
\]
\end{enumerate}
\end{corollary}

\begin{proof}
Es el caso $n=2$ del Teorema~\ref{thm:anillo-cohomologia-lente}: para
$m$ impar, $(\ZZ/m\ZZ)[\beta]/(\beta^n)$ con $n=2$ es
$(\ZZ/m\ZZ)[\beta]/(\beta^2)$, y el generador $\alpha\beta^{n-1}=\alpha\beta$
de $H^{2n-1}=H^3$ da $\alpha\beta\ne0$; como $\Lambda(\alpha)$ impone
$\alpha^2=0$, y $\beta^2\in H^4=0$ (la dimensión de $L$ es $3$) obliga
$\beta^2=0$ en el anillo truncado. Para $m=2$, $(\ZZ/2\ZZ)[\alpha]/(\alpha^{2n})$
con $n=2$ es $(\ZZ/2\ZZ)[\alpha]/(\alpha^4)$.
\end{proof}

\begin{theorem}[Clasificación por tipo de homotopía]
\label{thm:clasificacion-lente}
Sean $L(m,\ell)$ y $L(m,\ell')$ espacios lente de dimensión $3$.
\begin{enumerate}[label=(\roman*)]
\item Son homeomorfos si y solo si
\[
\ell'\equiv\pm\ell^{\pm1}\pmod m.
\]
\item Son del mismo tipo de homotopía si y solo si
\[
\ell\ell'\equiv\pm n^2\pmod m
\]
para algún entero $n$.
\end{enumerate}
\end{theorem}

\begin{proof}[Idea de la demostración --- no reproducida en detalle]
La necesidad del criterio de (ii) es el contenido del
Ejercicio~\ref{ex:distincion-lente}(a), demostrado arriba a partir del
Teorema~\ref{thm:anillo-cohomologia-lente} (el anillo de cohomología
módulo $m$ es un invariante de tipo de homotopía, y su clase en grado $2$
codifica $\ell$ salvo cuadrados y signo). La suficiencia de (ii) --- que
la congruencia \emph{basta} para construir una equivalencia de homotopía
explícita --- y ambas direcciones de (i) requieren la teoría clásica de
la \emph{torsión de Reidemeister--Franz--de~Rham}: a cada posible
equivalencia simple entre los complejos de cadenas $\ZZ[\ZZ/m\ZZ]$-libres
asociados a $L(m,\ell)$ y $L(m,\ell')$ se le asocia un invariante
algebraico (la torsión de Reidemeister) valuado en
$(\ZZ/m\ZZ)^*/\{\pm\ell'^{\ell'}\}$, y se muestra que dos espacios lente
son homeomorfos (resp.\ simple-homotópicamente equivalentes) si y solo si
sus torsiones coinciden hasta la acción del grupo de unidades adecuado;
para dimensión $3$ esto se reduce a la congruencia (i) (Reidemeister
1935, de~Rham 1936, clasificación de homeomorfismo completa por
J.~H.~C.~Whitehead 1941 y Brody 1960 para el caso topológico). Esta
maquinaria --- construir la torsión, verificar su invariancia, y calcular
el conjunto de valores realizables --- excede el alcance de estas notas;
se cita sin reproducir. Véase Cohen, \emph{A Course in Simple-Homotopy
Theory}, Cap.~II, \S 22--23, y Milnor, ``Whitehead torsion'', Bull.\
Amer.\ Math.\ Soc.\ 72 (1966), \S 12.
\end{proof}

\begin{remark}[Sobre el signo en (ii)]
\label{obs:signo-clasificacion}
El signo $\pm$ en (ii) es necesario, no opcional: la conjugación compleja
$(z_1,z_2)\mapsto(\bar z_1,\bar z_2)$ en $S^3\subset\CC^2$ intercambia la
acción de parámetro $\ell$ con la de parámetro $-\ell$, de modo que
$L(m,\ell)\cong L(m,-\ell)$ \emph{siempre}, vía un homeomorfismo que
invierte la orientación --- caso particular de (i) con exponente $-1$
compuesto con signo. Si el criterio de (ii) se escribiera sin el signo
($\ell\ell'\equiv n^2\pmod m$ solamente), dejaría de ser invariante bajo
$\ell'\mapsto-\ell'$: aplicado a $\ell'$ y a $-\ell'$ (que dan espacios
homeomorfos, luego ciertamente del mismo tipo de homotopía que
$L(m,\ell')$) daría en general dos respuestas distintas, una
contradicción. El argumento del Ejercicio~\ref{ex:distincion-lente} usando
el anillo de cohomología produce, correctamente, la condición con el
$\pm$: si $f^*(\alpha')=n\alpha$ para un isomorfismo de anillos $f^*$,
comparando los generadores de $H^2$ se obtiene
$\ell\equiv\pm\ell'n^2\pmod m$, equivalente a $\ell\ell'\equiv\pm(\ell'n)^2\pmod m$.
\end{remark}

\begin{corollary}
\label{cor:misma-homologia-lente}
Dos espacios lente de la misma dimensión y con el mismo valor de $m$ tienen
la misma homología integral y la misma cohomología integral como grupos
abelianos.
\end{corollary}

\begin{proof}
Inmediato: las fórmulas del Corolario~\ref{cor:homologia-lente} y del
Teorema~\ref{thm:cohomologia-lente} para $H_q(L_m(\ell_1,\dots,\ell_n);\ZZ)$
y $H^q(L_m(\ell_1,\dots,\ell_n);\ZZ)$ dependen únicamente de $m$ y de $n$
(equivalentemente, de la dimensión $2n-1$), sin mencionar en absoluto los
parámetros $\ell_1,\dots,\ell_n$; dos espacios lente con el mismo $m$ y la
misma dimensión producen así, término a término, los mismos grupos.
\end{proof}

\begin{corollary}
\label{cor:homologia-no-clasifica-lente}
La homología integral y la cohomología integral no bastan para clasificar
espacios lente hasta tipo de homotopía.
\end{corollary}

\begin{proof}
Los grupos aditivos calculados en el Corolario~\ref{cor:homologia-lente} y
el Teorema~\ref{thm:cohomologia-lente} solo dependen de $m$ y de la
dimensión, no de los parámetros $\ell_1,\dots,\ell_n$. Por ello, para
distinguir tipos de homotopía se requiere información adicional
proveniente del anillo de cohomología con coeficientes finitos y, en
dimensión $3$, de la teoría clásica de la torsión de Reidemeister--Franz.
\end{proof}

\begin{example}[Espacios lente no homotópicamente equivalentes con la misma homología]
\label{ex:L51-L52}
Los espacios $L(5,1)$ y $L(5,2)$ tienen la misma homología integral, la
misma cohomología integral y el mismo anillo abstracto de cohomología
módulo $5$, pero $L(5,1)\not\simeq L(5,2)$, pues la condición de
equivalencia homotópica exigiría $1\cdot2\equiv\pm n^2\pmod5$ para algún
$n$, y los cuadrados módulo $5$ son $\{0,1,4\}$, mientras que
$\{2,-2\}=\{2,3\}\pmod5$: ni $2$ ni $3$ son cuadrados módulo $5$.
\end{example}

\subsection{Cohomología de de Rham}

Para variedades diferenciables, la cohomología de de Rham ofrece una
presentación mediante formas diferenciales.

\begin{definition}[Complejo de de Rham]
\label{def:complejo-derham}
Sea $M$ una variedad suave. El \emph{complejo de de Rham} es
\[
0\to\Omega^0(M)\xrightarrow{d}\Omega^1(M)\xrightarrow{d}\Omega^2(M)\xrightarrow{d}\cdots,
\]
donde $\Omega^k(M)$ denota el espacio de $k$-formas diferenciales suaves y
$d$ es la diferencial exterior. Los \emph{grupos de cohomología de de
Rham} son
\[
H^k_{\mathrm{dR}}(M):=\ker\big(d:\Omega^k\to\Omega^{k+1}\big)\Big/\im\big(d:\Omega^{k-1}\to\Omega^k\big).
\]
\end{definition}

\begin{theorem}[de Rham, 1931]
\label{thm:derham}
Para toda variedad suave $M$ existe un isomorfismo natural (la
integración)
\[
H^k_{\mathrm{dR}}(M)\xrightarrow{\cong}H^k(M;\RR),
\]
donde el lado derecho es la cohomología singular con coeficientes reales.
\end{theorem}

\begin{proof}[Idea de la demostración]
El morfismo se define integrando una $k$-forma cerrada $\omega$ sobre
cadenas singulares suaves: $\omega\mapsto\big(\sigma\mapsto\int_\sigma\omega\big)$,
que por el Teorema de Stokes es efectivamente un cociclo del complejo
singular (con coeficientes reales) cuando $d\omega=0$, y una cofrontera
si $\omega$ es exacta; esto define
$I:H^k_{\mathrm{dR}}(M)\to H^k(M;\RR)$.

Que $I$ sea isomorfismo se demuestra, para $M$ que admite un \emph{buen
cubrimiento finito} $\{U_i\}_{i=1}^N$ (cada $U_i$ y cada intersección
finita no vacía $U_{i_1}\cap\cdots\cap U_{i_r}$ contráctil; todo
compacto suave lo admite), por inducción sobre $N$ usando el
Teorema~\ref{thm:mayer-vietoris-cohomologia} (Mayer--Vietoris en
cohomología singular) y su análogo en de Rham (Ejercicio 6 de esta
sección), comparados por el \emph{lema de los cinco} vía las
transformaciones naturales $I$ en cada término:
\[
\begin{tikzcd}
H^{k-1}_{\mathrm{dR}}(U\cap V) \arrow[r] \arrow[d,"I","\cong"'] & H^k_{\mathrm{dR}}(U\cup V) \arrow[r] \arrow[d,"I"] & H^k_{\mathrm{dR}}(U)\oplus H^k_{\mathrm{dR}}(V) \arrow[r] \arrow[d,"I","\cong"'] & H^k_{\mathrm{dR}}(U\cap V) \arrow[d,"I","\cong"'] \\
H^{k-1}(U\cap V;\RR) \arrow[r] & H^k(U\cup V;\RR) \arrow[r] & H^k(U;\RR)\oplus H^k(V;\RR) \arrow[r] & H^k(U\cap V;\RR)
\end{tikzcd}
\]
\emph{Caso base} ($N=1$, $M=U_1$ contráctil): $I$ es isomorfismo porque
ambos lados son $\RR$ en grado $0$ y $0$ en grados positivos (Lema de
Poincaré para de Rham, Ejercicio 4e; invarianza homotópica para
$H^k(\cdot;\RR)$, Teorema~\ref{thm:invarianza-homotopica-cohomologia}).
\emph{Paso inductivo}: escribiendo $M=\big(\bigcup_{i<N}U_i\big)\cup U_N$,
ambos abiertos ($\bigcup_{i<N}U_i$ con un buen cubrimiento de $N-1$
elementos, y $U_N$ con uno de un elemento) y su intersección
$\big(\bigcup_{i<N}U_i\big)\cap U_N=\bigcup_{i<N}(U_i\cap U_N)$ (que
también admite un buen cubrimiento de menos de $N$ elementos, a saber
$\{U_i\cap U_N\}_{i<N}$), la hipótesis de inducción da que $I$ es
isomorfismo en $\bigcup_{i<N}U_i$, en $U_N$ y en la intersección; el
lema de los cinco aplicado al diagrama de arriba da entonces que $I$
también es isomorfismo en $M$.

Para $M$ compacta esto concluye la demostración (todo compacto suave
admite un buen cubrimiento finito, por un argumento estándar con la
métrica riemanniana y convexidad geodésica). Para $M$ no compacta se
necesita además un argumento de paso al límite directo sobre una
exhaución por compactos, o el uso de cubrimientos buenos numerables
localmente finitos junto con particiones de la unidad subordinadas; esto
no se detalla aquí y se cita como el argumento clásico (de~Rham, 1931;
véase Bott--Tu, \emph{Differential Forms in Algebraic Topology}, \S I.8,
Teorema 8.9, para la demostración completa incluyendo el caso no
compacto).
\end{proof}

\begin{example}[de Rham en espacios lente]
\label{ex:derham-lente}
Como $L(m,\ell)$ es una variedad suave compacta de dimensión $3$ y
orientable,
\[
H^k_{\mathrm{dR}}(L(m,\ell))\cong H^k(L(m,\ell);\RR)\cong\begin{cases}\RR & k=0,3\\0 & k=1,2.\end{cases}
\]
El grupo de torsión $\ZZ/m\ZZ$ desaparece al tensorizar con $\RR$, de modo
que la cohomología de de Rham no puede distinguir espacios lente de
distintos $m$: para este propósito se necesitan coeficientes enteros o
módulo $m$.
\end{example}

\begin{remark}[Dualidad de Poincaré para espacios lente]
\label{obs:dualidad-poincare-lente}
Para una variedad cerrada orientable $M$ de dimensión $n$, la dualidad de
Poincaré establece $H^k(M;\ZZ)\cong H_{n-k}(M;\ZZ)$ a menos de torsión (más
precisamente, la dualidad opera sobre la cohomología libre y sobre la
parte de torsión de modo entrelazado). Para $L(m,\ell)$ de dimensión $3$:
\[
H^0\cong H_3\cong\ZZ, \qquad H^3\cong H_0\cong\ZZ, \qquad H^2\cong H_1\cong\ZZ/m\ZZ
\]
(la torsión se desplaza: $H^2\cong H_1$, no $H^2\cong H_2$, coherente con
el corrimiento de grado ya señalado tras el
Teorema~\ref{thm:cohomologia-lente}).
\end{remark}

\subsection*{Ejercicios}

\begin{enumerate}[label=\arabic*.,leftmargin=*]

\item (\textbf{Coeficientes universales y dualidad})
\begin{enumerate}[label=\alph*)]
\item Calcular $H^q(L(m,\ell);\ZZ)$ para toda $q$ usando el Teorema de
Coeficientes Universales.
\item Calcular $H^q(L(m,\ell);\ZZ/m\ZZ)$ para toda $q$.
\item Calcular $H^q(\RR P^n;\ZZ)$ y $H^q(\RR P^n;\ZZ/2\ZZ)$.
\item Calcular $H^q(\CC P^n;\ZZ)$.
\item Para un espacio lente de dimensión $2n-1$ cualquiera, establecer la
dualidad de Poincaré entre $H^q$ y $H_{2n-1-q}$.
\end{enumerate}

\noindent\textbf{Solución.}
\begin{enumerate}[label=\alph*)]
\item Aplicamos el T.C.U. a los grupos de homología de $L(m,\ell)$:
$H_0=\ZZ$, $H_1=\ZZ/m\ZZ$, $H_2=0$, $H_3=\ZZ$.
\[
H^0\cong\Hom(\ZZ,\ZZ)=\ZZ, \qquad
H^1\cong\Hom(\ZZ/m\ZZ,\ZZ)\oplus\mathrm{Ext}^1(\ZZ,\ZZ)=0\oplus0=0,
\]
\[
H^2\cong\Hom(0,\ZZ)\oplus\mathrm{Ext}^1(\ZZ/m\ZZ,\ZZ)=0\oplus\ZZ/m\ZZ=\ZZ/m\ZZ, \qquad
H^3\cong\Hom(\ZZ,\ZZ)\oplus\mathrm{Ext}^1(0,\ZZ)=\ZZ\oplus0=\ZZ.
\]
Esto coincide con el Corolario~\ref{cor:cohomologia-lente-3d}.
\item Con $R=\ZZ/m\ZZ$: $\mathrm{Ext}^1(\ZZ/m\ZZ,\ZZ/m\ZZ)\cong\ZZ/m\ZZ$ y
$\Hom(\ZZ/m\ZZ,\ZZ/m\ZZ)\cong\ZZ/m\ZZ$, luego
$H^q(L(m,\ell);\ZZ/m\ZZ)\cong\ZZ/m\ZZ$ para todo $0\le q\le3$, coincidiendo
con el Teorema~\ref{thm:ucl-lente}(ii).
\item Para $\RR P^n$ con la estructura CW estándar ($d_{2k}=0$,
$d_{2k+1}=2$):
\[
H^q(\RR P^n;\ZZ)=\begin{cases}
\ZZ & q=0,\\
\ZZ/2\ZZ & q\text{ par},\ 2\le q\le n\ (\text{si }n\text{ par})\text{ o }\le n-1\ (\text{si }n\text{ impar}),\\
\ZZ & q=n\ (\text{si }n\text{ impar}),\\
0 & \text{e.o.c.}
\end{cases}
\]
Con coeficientes $\ZZ/2\ZZ$: $H^q(\RR P^n;\ZZ/2\ZZ)\cong\ZZ/2\ZZ$ para
$0\le q\le n$, y $0$ en otro caso.
\item Como $H_q(\CC P^n;\ZZ)=\ZZ$ para $q$ par, $0\le q\le2n$, y $0$ para
$q$ impar (sin torsión), el T.C.U. da $H^q(\CC P^n;\ZZ)\cong\ZZ$ para $q$
par, $0\le q\le2n$, y $0$ en otro caso.
\item Sea $M=L_m(\ell_1,\dots,\ell_n)$. Con $R=\ZZ$, la dualidad de
Poincaré afirma que el morfismo $\frown[M]:H^k(M;\ZZ)\to H_{2n-1-k}(M;\ZZ)$
es isomorfismo en la parte libre y entrelaza los factores de torsión (el
grupo $\ZZ/m\ZZ$ en $H^{2j}$ corresponde al mismo grupo en
$H_{2n-1-2j}=H_{2(n-j)-1}$, que es impar y está en el rango donde el
Corolario~\ref{cor:homologia-lente} da torsión $\ZZ/m\ZZ$). Se verifica
inspeccionando los grupos calculados en los Teoremas~\ref{thm:cohomologia-lente}
y \ref{cor:homologia-lente}.
\end{enumerate}

\item (\textbf{Anillo de cohomología de espacios proyectivos})
\label{ex:anillo-proyectivos}
\begin{enumerate}[label=\alph*)]
\item Demostrar que $H^*(\CC P^n;\ZZ)\cong\ZZ[\alpha]/(\alpha^{n+1})$ con
$|\alpha|=2$.
\item Demostrar que $H^*(\RR P^n;\ZZ/2\ZZ)\cong(\ZZ/2\ZZ)[\alpha]/(\alpha^{n+1})$
con $|\alpha|=1$.
\item Concluir que $\CC P^m\not\simeq\CC P^n$ para $m\ne n$ y
$\RR P^m\not\simeq\RR P^n$ para $m\ne n$.
\item Mostrar que $\CC P^2\not\simeq S^2\vee S^4$ aunque tengan la misma
homología.
\end{enumerate}

\noindent\textbf{Solución.}
\begin{enumerate}[label=\alph*)]
\item Sea $\alpha\in H^2(\CC P^n;\ZZ)$ el generador. La fórmula de Künneth
y el axioma del producto de taza en la construcción celular de $\CC P^n$
muestran que la potencia $\alpha^k\in H^{2k}(\CC P^n;\ZZ)$ es generadora
para $0\le k\le n$ y que $\alpha^{n+1}=0$. La estructura de álgebra
graduada conmutativa es la del anillo de polinomios truncado.
\item Análogo, reemplazando $\ZZ$ por $\ZZ/2\ZZ$ y usando el generador
$\alpha\in H^1(\RR P^n;\ZZ/2\ZZ)$. La clave es que la clase $\alpha$ no es
de $2$-torsión (en el sentido de anillo: $\alpha\ne0$), y $\alpha^k$ genera
$H^k(\RR P^n;\ZZ/2\ZZ)$.
\item Los anillos $\ZZ[\alpha]/(\alpha^{m+1})$ y $\ZZ[\alpha]/(\alpha^{n+1})$
no son isomorfos para $m\ne n$ (difieren en su nilpotencia). Luego
$\CC P^m\not\simeq\CC P^n$; análogamente para $\RR P^m$ y $\RR P^n$.
\item $H^*(S^2\vee S^4;\ZZ)$ tiene producto de taza nulo en grados
positivos (la reducida de un bouquet satisface
$\tilde H^*(X\vee Y)\cong\tilde H^*(X)\oplus\tilde H^*(Y)$ como
$\ZZ$-módulos, y los productos cruzados entre factores son cero vía la
fórmula de Künneth). En cambio, en $\CC P^2$ el generador $\alpha\in H^2$
satisface $\alpha^2\ne0\in H^4$. Luego $\CC P^2\not\simeq S^2\vee S^4$.
\end{enumerate}

\item (\textbf{Cohomología celular: cálculos explícitos}) Calcular la
cohomología celular de los siguientes espacios (incluyendo el complejo de
cocadenas y los morfismos coborde):
\begin{enumerate}[label=\alph*)]
\item $L(4,1)$ y $L(4,3)$ con coeficientes en $\ZZ$.
\item $\mathbb T^2$ (toro), $K$ (botella de Klein) y $\RR P^2$.
\item La superficie orientable $T_g$ de género $g$ con coeficientes en
$\ZZ$ y en $\ZZ/2\ZZ$.
\item El espacio $X=S^1\times S^2$.
\end{enumerate}

\noindent\textbf{Solución.}
\begin{enumerate}[label=\alph*)]
\item Para $L(4,\ell)$ el complejo celular es $C_0=C_1=C_2=C_3=\ZZ$ con
bordes $d_1=0$, $d_2=4$, $d_3=0$. El complejo dual (cocadenas) es
idéntico en rango, con coborde $\delta^0=0$, $\delta^1=4\cdot$,
$\delta^2=0$. Luego $H^0=\ZZ$, $H^1=0$, $H^2=\ZZ/4\ZZ$, $H^3=\ZZ$ para
ambos $L(4,1)$ y $L(4,3)$ (la homología/cohomología sobre $\ZZ$ no los
distingue --- para eso se necesita el anillo de cohomología, y aquí no
distingue tampoco: $L(4,1)\cong L(4,3)$ como espacios, pues
$3\equiv-1\equiv1^{-1}\pmod4$, luego son homeomorfos por el
Teorema~\ref{thm:clasificacion-lente}(i)).
\item $\mathbb T^2$: complejo de cocadenas con $\delta^0=0$, $\delta^1=0$,
$\delta^2=0$ (todos los bordes son cero por orientabilidad). Luego
$H^q(\mathbb T^2;\ZZ)=\ZZ,\ZZ^2,\ZZ,0,\dots$ para $q=0,1,2,\dots$.

$K$ (botella de Klein): con un vértice, dos aristas $a,b$ y una cara pegada
por la palabra $abab^{-1}$, el complejo celular de cadenas es
$C_0=\ZZ,C_1=\ZZ^2,C_2=\ZZ$ con $d_1=0$ y $d_2(e^2)=2a$ (en la base
$(a,b)$ de $C_1$), de modo que $H_0(K;\ZZ)=\ZZ$, $H_1(K;\ZZ)=\ZZ\oplus\ZZ/2\ZZ$
y $H_2(K;\ZZ)=0$. Dualizando, $\delta^0=0$ y $\delta^1:\ZZ^2\to\ZZ$ es la
funcional $(x,y)\mapsto2x$. Por el T.C.U. (o directamente del complejo
dual),
\[
H^0(K;\ZZ)=\ZZ, \qquad H^1(K;\ZZ)=\ker(\delta^1)=\ZZ, \qquad
H^2(K;\ZZ)=\mathrm{coker}(\delta^1)\oplus\Hom(H_2,\ZZ)\cong\ZZ/2\ZZ.
\]
Es decir, $H^2(K;\ZZ)\cong\ZZ/2\ZZ$ --- \textbf{no} $0$: la torsión de
$H_1(K;\ZZ)=\ZZ\oplus\ZZ/2\ZZ$ reaparece, desplazada un grado, en
$H^2(K;\ZZ)$, exactamente el mismo fenómeno de
$H_2(\RR P^2;\ZZ)=0$ frente a $H^2(\RR P^2;\ZZ)\cong\ZZ/2\ZZ$ (calculado a
continuación) y de la torsión de los espacios lente
(\S\ref{sec:anillo-cohomologia-lente}). Sería un error --- fácil de
cometer si solo se mira el rango del complejo dual sin aplicar
correctamente el T.C.U.\ --- concluir $H^2(K;\ZZ)=0$.

$\RR P^2$: $\delta^1=2\cdot\mathrm{id}:\ZZ\to\ZZ$, luego $H^0=\ZZ$,
$H^1=0$, $H^2=\ZZ/2\ZZ$.
\item $H^q(T_g;\ZZ)\cong\ZZ,\ZZ^{2g},\ZZ$ para $q=0,1,2$, y cero en otro
caso (todos los bordes celulares son cero por orientabilidad, igual que
para $\mathbb T^2$). Con $\ZZ/2\ZZ$: ídem, reemplazando $\ZZ$ por
$\ZZ/2\ZZ$.
\item $S^1\times S^2$ tiene una estructura CW producto. Por Künneth:
$H^q(S^1\times S^2;\ZZ)=\ZZ,\ZZ,\ZZ,\ZZ$ para $q=0,1,2,3$, y cero en otro
caso. El producto de taza está determinado por el generador
$\alpha\in H^1(S^1)$ y $\beta\in H^2(S^2)$: $\alpha\cup\beta$ genera
$H^3$, y $\alpha^2=\beta^2=0$.
\end{enumerate}

\item (\textbf{Cohomología de de Rham: cálculos})
\begin{enumerate}[label=\alph*)]
\item Calcular $H^k_{\mathrm{dR}}(\RR^n)$ para toda $k$ usando el Lema de
Poincaré.
\item Calcular $H^k_{\mathrm{dR}}(S^n)$ para toda $k$.
\item Calcular $H^k_{\mathrm{dR}}(\mathbb T^n=(S^1)^n)$ para toda $k$, e
identificar clases generadoras explícitas.
\item Sea $\omega=\dfrac{x\,dy-y\,dx}{x^2+y^2}$ en $\RR^2\setminus\{0\}$.
\begin{enumerate}[label=\arabic*)]
\item Mostrar que $\omega$ es cerrada.
\item Mostrar que $\omega$ no es exacta en $\RR^2\setminus\{0\}$.
\item Mostrar que $[\omega]$ genera $H^1_{\mathrm{dR}}(S^1)\cong\RR$.
\end{enumerate}
\item Enunciar y demostrar el Lema de Poincaré: toda forma cerrada es
localmente exacta.
\end{enumerate}

\noindent\textbf{Solución.}
\begin{enumerate}[label=\alph*)]
\item Por el Lema de Poincaré: $H^0_{\mathrm{dR}}(\RR^n)\cong\RR$
(constantes), y $H^k_{\mathrm{dR}}(\RR^n)=0$ para $k\ge1$ (toda forma
cerrada de grado positivo es exacta en un espacio contráctil).
\item $H^0_{\mathrm{dR}}(S^n)\cong\RR$ y $H^n_{\mathrm{dR}}(S^n)\cong\RR$
(generada por la forma de volumen
$\omega_n=\sum_i(-1)^ix_i\,dx_0\wedge\cdots\widehat{dx_i}\cdots\wedge dx_n$
restringida a $S^n$). Los demás grupos son cero.
\item Para $\mathbb T^n=(S^1)^n$ con coordenadas angulares
$\theta_1,\dots,\theta_n$, las formas $d\theta_{i_1}\wedge\cdots\wedge d\theta_{i_k}$
para $i_1<\cdots<i_k$ son cerradas y sus clases generan
$H^k_{\mathrm{dR}}(\mathbb T^n)\cong\RR^{\binom nk}$, consistente con el
isomorfismo de de Rham: $H^k_{\mathrm{dR}}(\mathbb T^n)\cong H^k(\mathbb T^n;\RR)\cong\RR^{\binom nk}$.
\item
\begin{enumerate}[label=\arabic*)]
\item En coordenadas polares $r,\theta$, $\omega=d\theta$, y $d(d\theta)=0$.
En coordenadas cartesianas:
$d\omega=\dfrac{(x^2+y^2)(dy\wedge dy-dx\wedge dx)+\cdots}{(x^2+y^2)^2}=0$
por antisimetría del producto exterior.
\item Si $\omega=df$ para alguna $f\in C^\infty(\RR^2\setminus\{0\})$,
entonces $\oint_{S^1}\omega=\oint_{S^1}df=f(2\pi)-f(0)=0$, pero
$\oint_{S^1}\omega=\int_0^{2\pi}d\theta=2\pi\ne0$. Contradicción.
\item $[\omega]\ne0$ por (2), y $H^1_{\mathrm{dR}}(S^1)\cong H^1(S^1;\RR)\cong\RR$
tiene dimensión $1$, luego $[\omega]$ lo genera.
\end{enumerate}
\item \textbf{Lema de Poincaré.} Sea $U\subset\RR^n$ un abierto
estrellado con centro $0$, y $\omega\in\Omega^k(U)$ con $d\omega=0$,
$k\ge1$. Se define el operador de homotopía
\[
(K\omega)(x)=\int_0^1t^{k-1}\iota_x\omega(tx)\,dt,
\]
donde $\iota_x$ denota la contracción con el vector posición $x$ como
campo vectorial radial. Una verificación directa da la fórmula
$d(K\omega)+K(d\omega)=\omega$. Si $d\omega=0$, entonces $d(K\omega)=\omega$,
es decir, $\omega$ es exacta en $U$.
\end{enumerate}

\item (\textbf{Distinción de espacios lente por cohomología})
\label{ex:distincion-lente}
\begin{enumerate}[label=\alph*)]
\item Usando el anillo de cohomología con coeficientes en $\ZZ/m\ZZ$,
demostrar que $L(m,\ell)\not\simeq L(m,\ell')$ si no existe $n\in\ZZ$ con
$\ell\ell'\equiv\pm n^2\pmod m$.
\item Concluir que $L(5,1)\not\simeq L(5,2)$.
\item Demostrar que $L(5,1)$ y $L(5,2)$ \emph{son} del mismo tipo de
homotopía racional (es decir, $L(5,1)\otimes_\ZZ\QQ\simeq L(5,2)\otimes_\ZZ\QQ$
en el sentido de que sus cohomologías de de Rham coinciden).
\item Demostrar que $H^*_{\mathrm{dR}}(L(m,\ell))\cong H^*_{\mathrm{dR}}(L(m,\ell'))$
para todo $m,\ell,\ell'$, y concluir que la cohomología de de Rham no
puede distinguir espacios lente.
\end{enumerate}

\noindent\textbf{Solución.}
\begin{enumerate}[label=\alph*)]
\item Por el Teorema~\ref{thm:anillo-cohomologia-lente}, el anillo
$H^*(L(m,\ell);\ZZ/m\ZZ)$ está determinado en grado $2$ por el cuadrado del
generador $\alpha$ de $H^1$, el cual está dado por la forma $\ell\alpha^2$.
Una equivalencia de homotopía $f:L(m,\ell)\to L(m,\ell')$ debe inducir un
isomorfismo de anillos en cohomología. Si $f^*(\alpha')=n\alpha$, entonces
$f^*(\ell'\alpha'^2)=\ell'n^2\alpha^2$, mientras que en $L(m,\ell)$ el
generador de $H^2$ es $\ell\alpha^2$. Para compatibilidad,
$\ell\equiv\ell'n^2\pmod m$; equivalentemente,
$\ell\ell'\equiv(\ell'n)^2\pmod m$. El signo $\pm$ de la
Observación~\ref{obs:signo-clasificacion} entra al notar que $\ell'$ mismo
está definido solo salvo signo ($L(m,\ell')\cong L(m,-\ell')$), de modo que
la condición necesaria correcta, invariante bajo esa ambigüedad, es
$\ell\ell'\equiv\pm k^2\pmod m$ para algún entero $k$. Si no existe tal
$k$, $f^*$ no puede ser isomorfismo de anillos, contradiciendo la
existencia de la equivalencia de homotopía.
\item Con $\ell=1$, $\ell'=2$, $m=5$: necesitamos $n^2\equiv\pm2\pmod5$
para algún $n$. Los cuadrados módulo $5$ son $\{0,1,4\}$, y
$\{2,-2\}\equiv\{2,3\}\pmod5$: ninguno de los dos está en $\{0,1,4\}$.
Luego $L(5,1)\not\simeq L(5,2)$.
\item $H^k(L(5,1);\RR)\cong H^k_{\mathrm{dR}}(L(5,1))\cong H^k_{\mathrm{dR}}(L(5,2))$
para toda $k$, pues la torsión $\ZZ/5\ZZ$ desaparece al tensorizar con
$\RR$. La equivalencia racional entre ambos espacios se sigue de que sus
cohomologías reales (y, de hecho, sus tipos de homotopía racional en el
sentido de la teoría de Sullivan) coinciden.
\item Para cualquier $m,\ell,\ell'$, la cohomología
$H^k_{\mathrm{dR}}(L_m(\ell))\cong\RR$ solo para $k=0,2n-1$ y $0$ en otro
caso (independientemente de $\ell$). La forma de torsión $\ZZ/m\ZZ$ no
contribuye a la cohomología real. Luego la cohomología de de Rham es la
misma para todos los espacios lente de dimensión $2n-1$ y valor de $m$
fijo, y no puede distinguirlos --- consistente con el
Ejemplo~\ref{ex:derham-lente}.
\end{enumerate}

\item (\textbf{Sucesión de Mayer--Vietoris en cohomología})
\begin{enumerate}[label=\alph*)]
\item Enunciar la sucesión exacta larga de Mayer--Vietoris en cohomología
para $X=U\cup V$.
\item Calcular $H^k_{\mathrm{dR}}(S^n)$ usando Mayer--Vietoris de de Rham.
\item Calcular $H^k(S^1\times S^n;\ZZ)$ usando Mayer--Vietoris en
cohomología singular.
\item Sea $f:S^n\to S^n$. Demostrar que $f^*:H^n(S^n;\ZZ)\to H^n(S^n;\ZZ)$
es la multiplicación por $\deg(f)$.
\end{enumerate}

\noindent\textbf{Solución.}
\begin{enumerate}[label=\alph*)]
\item Es el Teorema~\ref{thm:mayer-vietoris-cohomologia}.
\item Tomamos $U=S^n\setminus\{N\}$, $V=S^n\setminus\{S\}$, ambos
contráctiles, y $U\cap V\simeq S^{n-1}$. Para $k<n-1$, la sucesión da
$H^k_{\mathrm{dR}}(S^n)\cong H^{k-1}_{\mathrm{dR}}(S^{n-1})$ (inducción).
Para $k=n$, el segmento
\[
0\to H^{n-1}_{\mathrm{dR}}(U\cap V)\xrightarrow{\delta^{-1}}H^n_{\mathrm{dR}}(S^n)\to H^n_{\mathrm{dR}}(U)\oplus H^n_{\mathrm{dR}}(V)=0
\]
da $H^n_{\mathrm{dR}}(S^n)\cong H^{n-1}_{\mathrm{dR}}(S^{n-1})\cong\cdots\cong H^1_{\mathrm{dR}}(S^1)\cong\RR$.
\item Por Künneth: $H^k(S^1\times S^n;\ZZ)\cong H^k(S^1;\ZZ)\oplus H^{k-1}(S^1;\ZZ)\otimes H^1(S^n;\ZZ)\oplus\cdots$.
Con el resultado de Künneth para cohomología, se obtiene $\ZZ$ en
$k=0,1,n,n+1$ y $0$ en otro caso.
\item Como $H^n(S^n;\ZZ)\cong\ZZ$, el homomorfismo $f^*$ es la
multiplicación por algún entero $d$. Por dualidad con la homología (vía el
coeficiente universal), $d$ coincide con $f_*$ en $H_n(S^n;\ZZ)$, que es
por definición $\deg(f)$.
\end{enumerate}

\item (\textbf{Fórmula de Künneth y cohomología del producto})
\begin{enumerate}[label=\alph*)]
\item Enunciar el Teorema de Künneth para cohomología singular con
coeficientes en un campo $k$.
\item Calcular $H^*(S^p\times S^q;k)$ como anillo para $p\ne q$.
\item Calcular $H^*(S^n\times S^n;k)$ y describir su estructura de álgebra
(¿es isomorfo al anillo de cohomología de $S^{2n}\vee S^n\vee S^n$?).
\item Sea $\Delta:S^n\to S^n\times S^n$ la diagonal $\Delta(x)=(x,x)$.
Calcular $\Delta^*:H^*(S^n\times S^n;k)\to H^*(S^n;k)$ y usar esto para
mostrar que $\Delta$ no es homotópica a la inclusión en ningún factor.
\end{enumerate}

\noindent\textbf{Solución.}
\begin{enumerate}[label=\alph*)]
\item Si $k$ es un campo, $H^*(X\times Y;k)\cong H^*(X;k)\otimes_kH^*(Y;k)$
como álgebras graduadas.
\item Para $p\ne q$: $H^*(S^p\times S^q;k)\cong k[\alpha,\beta]/(\alpha^2,\beta^2)$
con $|\alpha|=p$, $|\beta|=q$. Como $p\ne q$, $\alpha\beta=(-1)^{pq}\beta\alpha$
y el álgebra es un álgebra exterior en dos generadores de distintos grados.
\item Para $p=q=n$: $H^*(S^n\times S^n;k)\cong k\oplus k^2\oplus k$ en
grados $0,n,2n$. El generador de $H^{2n}$ es $\alpha\otimes\beta$
(producto de los generadores de cada factor). En cambio,
$H^*(S^{2n}\vee S^n\vee S^n;k)$ tiene producto nulo en grados positivos,
así los anillos no son isomorfos.
\item $\Delta^*(\alpha\otimes1)=\Delta^*(1\otimes\alpha)=\alpha$
(generador de $H^n(S^n;k)$). En particular
$\Delta^*(\alpha\otimes\beta)=\alpha^2$ ($=0$ si $n$ impar, o
$=\alpha^2$ si $n$ par). La inclusión en el primer factor
$i_1:S^n\to S^n\times S^n$ satisface $i_1^*(1\otimes\beta)=0$, mientras
que $\Delta^*(1\otimes\beta)=\alpha\ne0$. Luego $\Delta^*\ne i_1^*$ como
homomorfismos de anillos, de modo que $\Delta\not\simeq i_1$ (y
análogamente para $i_2$).
\end{enumerate}

\item (\textbf{Cohomología de de Rham y formas de volumen})
\begin{enumerate}[label=\alph*)]
\item Demostrar que una variedad orientable compacta sin frontera de
dimensión $n$ satisface $H^n_{\mathrm{dR}}(M)\ne0$, usando la forma de
volumen.
\item Deducir que no existe una retracción $r:D^{n+1}\to S^n$. (Nueva
demostración del Teorema de Punto Fijo de Brouwer.)
\item Demostrar que una variedad compacta sin frontera no puede embeberse
en $\RR^n$ como subvariedad de dimensión completa.
\item Sea $M$ una variedad compacta orientable sin frontera. Mostrar que
cualquier aplicación $f:M\to S^n$ con $\deg(f)\ne0$ es sobreyectiva, usando
formas de volumen.
\end{enumerate}

\noindent\textbf{Solución.}
\begin{enumerate}[label=\alph*)]
\item Sea $\Omega\in\Omega^n(M)$ una forma de volumen (cerrada, $d\Omega=0$,
pues $\Omega^{n+1}=0$). Si $\Omega=d\eta$, el Teorema de Stokes daría
$\int_M\Omega=\int_Md\eta=\int_{\partial M}\eta=0$ (pues $\partial M=\varnothing$).
Pero $\int_M\Omega>0$ por definición de forma de volumen. Contradicción.
Luego $[\Omega]\ne0$ en $H^n_{\mathrm{dR}}(M)$.
\item Si existiera una retracción $r:D^{n+1}\to S^n$, entonces
$H^n_{\mathrm{dR}}(S^n)\hookrightarrow H^n_{\mathrm{dR}}(D^{n+1})=0$,
contradiciendo $H^n_{\mathrm{dR}}(S^n)\cong\RR\ne0$.
\item Si $M^n\subset\RR^n$ (dimensión completa), entonces $M$ sería abierta
en $\RR^n$ (invarianza del dominio), luego no compacta o con frontera.
Contradicción.
\item Sea $\omega\in\Omega^n(S^n)$ la forma de volumen. Entonces
$\int_Mf^*\omega=\deg(f)\int_{S^n}\omega\ne0$. Si $f$ no fuera
sobreyectiva, existiría $y_0\notin f(M)$, y $f$ factorizaría a través de
$S^n\setminus\{y_0\}\simeq\RR^n$ (contráctil), dando $f^*\omega=d(\cdots)$
y $\int_Mf^*\omega=0$. Contradicción.
\end{enumerate}

\end{enumerate}

\end{document}

